% =============================================================================
%  main.tex
%  Theoretical Foundations of tpeanuts
%  Fast Simulation of Neutrino Oscillations in Matter (MSc Thesis)
% =============================================================================
\documentclass[11pt,a4paper,openany]{report}

\usepackage[a4paper,margin=2.6cm,marginparwidth=2.1cm,marginparsep=4mm]{geometry}
\usepackage{tpeanuts-theory}

\addbibresource{references.bib}

\title{\Huge Theoretical Foundations of \texttt{tpeanuts}\\[4pt]
  \Large Fast Simulation of Neutrino Oscillations in Matter}
\author{Juan Ramon Diaz Santos}
\date{July 2026 --- living document, generated from the source code}

\begin{document}

\begin{titlepage}
  \centering
  \vspace*{2cm}
  {\Huge\bfseries\color{coreColor} Theoretical Foundations\\[4pt] of \texttt{tpeanuts}\par}
  \vspace{0.8cm}
  {\Large Fast Simulation of Neutrino Oscillations in Matter\par}
  \vspace{1.5cm}
  {\large MSc Thesis (Trabajo Fin de M\'aster)\par}
  \vspace{0.3cm}
  {\large Juan Ramon Diaz Santos \\ \texttt{diazjuan@alumni.uv.es}\par}
  \vspace{0.6cm}
  {\normalsize Supervisors: Roberto Ruiz de Austri Bazan \quad\&\quad Michele Lucente\par}
  \vfill
  {\normalsize July 2026\par}
  \vspace{0.5cm}
  \begin{tcolorbox}[colback=coreColor!5,colframe=coreColor,width=0.85\textwidth,sharp corners=all,rounded corners=all,arc=1mm]
    \small
    This document accompanies the \texttt{tpeanuts} code repository and aims
    to connect, package by package and module by module, the physics of
    neutrino oscillations with its concrete implementation. Blue boxes
    (\textbf{Theoretical Foundation}) contain the derivations; green boxes
    (\textbf{Implementation}) link that theory to actual functions and
    classes in the code; orange boxes collect approximations and ranges of
    validity; text highlighted in \keyresult{yellow} marks key results.
    Sections flagged as \textcolor{pendingColor}{\textbf{pending}} are
    scaffolding ready to be completed incrementally.
  \end{tcolorbox}
\end{titlepage}

\pagenumbering{roman}
\tableofcontents
\listoftodos[Pending sections and notes]
\clearpage
\pagenumbering{arabic}

% =============================================================================
\part{General Framework}
\chapter{General Theoretical Framework and Conventions}
\label{ch:intro}

% This chapter's Implementation-in-tpeanuts blocks use full descriptive
% paragraphs rather than the one-line module purpose used from Chapter 2
% onwards, so the shared \modulehead's blanket italics (meant for a short
% phrase) reads poorly here. Locally drop it to plain upright text -- italics
% stay available for deliberately highlighting a single term -- and restore
% the shared definition at the end of the chapter so later chapters are
% unaffected.
\let\savedmoduleheadchOne\modulehead
\renewcommand{\modulehead}[2]{%
  \noindent\pkgpath{#1}\par
  \noindent#2\par\medskip
}

\section{Purpose and how to read this document}

\texttt{tpeanuts} is a \texttt{PyTorch} rewrite (GPU-first, with batching
support) of the original \texttt{peanuts} code \parencite{LucentePeanuts},
extended with BSM scenarios (NSI and a 3+1 sterile neutrino), a proprietary
perturbative evolution scheme, and interface layers to external reference
engines (\texttt{nuSQuIDS}, \texttt{MCEq}, Honda/HKKM tables,
\texttt{pymsis}).

Every section in the chapters that follow corresponds to \emph{one source
file} in the repository and always follows the same structure:

\begin{itemize}
  \item \modulehead{package/module.py}{one line describing its purpose}
    identifies the exact file.
  \item The blue box \textbf{Theoretical Foundation} contains the physics and
    mathematical derivations that justify what the module computes.
  \item The green box \textbf{Implementation in \texttt{tpeanuts}} translates
    that theory into the actual functions/classes, with their essential
    signature and the sign, unit, or tensor-shape conventions that must be
    respected to use them correctly.
  \item The orange box \textbf{Approximations and Limitations} (when
    applicable) documents the range of validity, limiting cases, and
    numerical simplifications.
  \item The grey box \textbf{Key References} links the specific bibliography
    for that piece of physics.
  \item Central results are \keyresult{highlighted}; sections that do not yet
    have a complete derivation are explicitly flagged as pending and are
    listed in the pending-tasks index right after the table of contents.
\end{itemize}

\section{The physics problem: neutrino oscillation in matter}

A neutrino is produced and detected in a flavour eigenstate
$\ket{\nu_\alpha}$, $\alpha \in \{e,\mu,\tau\}$ (plus $\ket{\nu_s}$ in the
sterile sector, \cref{ch:core-bsm}), but propagates as a superposition of
mass eigenstates $\ket{\nu_i}$. The PMNS mixing matrix
\parencite{Pontecorvo1957,MakiNakagawaSakata1962} connects the two bases,
\begin{equation}
  \ket{\nu_\alpha} = \sum_i U_{\alpha i}\,\ket{\nu_i}.
\end{equation}
In vacuum each mass eigenstate evolves with a kinetic phase $e^{-iE_i L}$;
when crossing matter an additional effective potential from coherent forward
interaction appears (the MSW effect,
\parencite{Wolfenstein1978,MikheyevSmirnov1985}). In both cases the evolution
is summarised by an effective Hamiltonian $H(x)$ acting on the vector of
flavour amplitudes,
\begin{equation}
  i \frac{d}{dx}\, \psi(x) = H(x)\, \psi(x),
  \qquad
  \psi = \begin{pmatrix}\psi_e\\ \psi_\mu\\ \psi_\tau\end{pmatrix},
  \label{eq:schrodinger-flavour}
\end{equation}
where $x$ is the dimensionless propagation coordinate defined in
\cref{sec:units}. This document organises \emph{all} propagation code around
three objects that appear systematically in every module:

\begin{enumerate}
  \item the \textbf{Hamiltonian} $H = H_{\rm kin} + H_{\rm mat}$
    (\cref{ch:core-common,ch:core-bsm}),
  \item the \textbf{evolution operator} (evolutor) $S$ that formally solves
    \labelcref{eq:schrodinger-flavour},
    $\psi(x_2) = S(x_2, x_1)\,\psi(x_1)$, obtained either exactly (constant
    density, \cref{ch:core-perturbative}) or numerically via segmented
    exponentiation (\cref{ch:core-numerical}),
  \item the \textbf{transition probability}
    $P(\nu_\alpha \to \nu_\beta) = |S_{\beta\alpha}|^2$
    (\cref{ch:core-common}), the final observable quantity.
\end{enumerate}

The \emph{Physical Media} chapters (\cref{ch:vacuum,ch:solar,ch:earth,%
ch:atmosphere}) instantiate these three objects for each concrete
propagation environment (vacuum, solar interior, Earth, atmosphere), and the
\emph{External Engines} chapters document how those results are
cross-checked against independent reference codes.

\section{Notation and global conventions}
\label{sec:units}

\begin{theorybox}
Physics is worked out in natural units $\hbar = c = 1$ unless explicitly
converted to laboratory units. Public input quantities always use:
\begin{itemize}
  \item neutrino energy $E$ in $\mathrm{MeV}$,
  \item squared mass differences $\Delta m^2$ in $\mathrm{eV}^2$,
  \item path lengths $L$ in metres,
  \item electron density $n_e$ in $\mathrm{mol\,cm^{-3}}$ (molar electron
    density, not mass density).
\end{itemize}
The Hamiltonian is made dimensionless by introducing an \emph{evolution
scale} $L_{\rm scale}$ (by default the Earth radius $R_\oplus$) and the
propagation coordinate
\begin{equation}
  x \equiv \frac{L}{L_{\rm scale}},
\end{equation}
so that $H(x)$ in \labelcref{eq:schrodinger-flavour} is unitless. The two
Hamiltonian terms are made dimensionless as
\begin{align}
  k_i &= L_{\rm scale}\,\frac{\Delta m^2_{i1}}{2E}\,\frac{1}{\hbar c}
       &&\text{(kinetic term, see \cref{ch:core-common}),}\\
  V &= \pm\sqrt{2}\,G_F\, n_e\, L_{\rm scale}
       &&\text{(MSW matter potential, see \cref{ch:core-common}),}
\end{align}
with sign $+$ for neutrinos and $-$ for antineutrinos (the standard sign
convention for the charged-current term).
\end{theorybox}

\begin{implbox}
\begin{itemize}
  \item \modulehead{util/constant.py}{single source of truth for physical
    constants: $G_F$ in $\mathrm{MeV^{-2}}$, $\hbar$, $c$, $\hbar c$, $N_A$,
    Earth and Sun radii, the astronomical unit, the proton mass, and the
    mass-to-nucleon-density conversion factor.}
  \item \modulehead{util/context.py}{\pyclass{RuntimeContext} bundles
    \code{(device, dtype)} and is passed explicitly throughout the stack
    instead of duplicating those two arguments in every function.}
  \item \modulehead{util/type.py}{tensor conversion and validation: resolving
    the complex dtype paired with a real dtype
    (\pyfunc{cdtype\_from\_real}), building/validating 3- or 4-flavour state
    vectors (\pyfunc{state\_tensor}), and broadcasting flavour vectors to
    batched shapes (\pyfunc{broadcast\_last3}).}
  \item \modulehead{util/torch\_util.py}{default device/dtype resolution,
    generic tensor broadcasting, batch flattening, and the convention
    \pyfunc{enforce\_identity\_for\_zero\_length} (a zero-length evolutor
    must collapse exactly to the identity, which numerical exponentials do
    not automatically guarantee in the $x\to 0$ limit).}
  \item \modulehead{core/common/neutrino.py}{canonical flavour-index
    convention: $e\to0$, $\mu\to1$, $\tau\to2$ (plus $s\to3$ in the sterile
    sector), used throughout the code to index tensors.}
\end{itemize}
\end{implbox}

\subsection*{Probability matrix convention}

\begin{theorybox}
Given the flavour evolutor $S$, the final amplitude is
$\psi_\beta(x_2) = \sum_\alpha S_{\beta\alpha}\,\psi_\alpha(x_1)$, so
\begin{equation}
  \keyresult{P(\nu_\alpha \to \nu_\beta) = \left|S_{\beta\alpha}\right|^2},
\end{equation}
with the \emph{initial flavour} index in the column and the \emph{final
flavour} index in the row. This is the convention fixed, once and for the
whole project, by \texttt{core.common.probability}.
\end{theorybox}

\begin{implbox}
\modulehead{core/common/probability.py}{full derivation in
\cref{ch:core-common}. Every \code{*\_probability\_state(...)} function in
the medium modules (\code{vacuum\_probability\_state},
\code{solar\_probability\_state}, \code{earth\_probability\_state},
\code{atmosphere\_probability\_state}) delegates to the utilities in this
module to keep the index convention identical across the whole project.}
\end{implbox}

\section{Coherent states versus incoherent mixtures}
\label{sec:coherent-vs-incoherent}

\begin{theorybox}
The code systematically distinguishes two representations of the state being
propagated, selected with the \code{massbasis} parameter:
\begin{description}
  \item[Coherent propagation] (\code{massbasis=False}). The state is a
    flavour-amplitude vector $\psi_\alpha$ carrying physical relative
    phases. The evolutor is applied directly,
    $\psi_{\rm final} = S\,\psi_{\rm inicial}$, and
    $P_\alpha = |\psi_{{\rm final},\alpha}|^2$. This is the correct regime
    when trajectories are short compared to the coherence length (vacuum,
    atmosphere, Earth in a single leg).
  \item[Incoherent mass mixture] (\code{massbasis=True}). The state is a
    vector of real weights $w_i \geq 0$, $\sum_i w_i = 1$, over mass
    eigenstates \emph{without} relative-phase information. The final flavour
    probability is
    \begin{equation}
      P_\alpha = \sum_i \left| \bigl(S\, U_{\rm PMNS}\bigr)_{\alpha i}\right|^2 w_i.
    \end{equation}
    This is the physically correct regime for solar neutrinos: the spread at
    the production point and the enormous Sun--Earth distance completely
    decohere the mass superposition before it reaches Earth
    (\cref{ch:solar}), so only the populations $w_i$ survive, not the
    phases.
\end{description}
\end{theorybox}

\begin{notebox}
The choice \code{massbasis=True/False} is not a numerical detail: it is a
physical statement about whether the neutrino wave-packet's coherence length
is larger or smaller than the scale of the problem. Using the incoherent
mode on a short leg (e.g.\ only inside the atmosphere) would discard real
interference; using the coherent mode for Sun--Earth propagation would
overstate coherence and produce spurious interference patterns that are not
observed.
\end{notebox}

\section{Named presets and parameter traceability}

\begin{implbox}
\modulehead{core/common/presets.py}{single registry
(\pyfunc{register\_preset}/\pyfunc{get\_preset}/\pyfunc{list\_presets}) for
every named parameter set in the project: \pyclass{OSCILLATION\_PRESETS}
(mixing angles and $\Delta m^2$, for both the standard 3-flavour sector and
the 3+1 sterile extension) and \pyclass{NSI\_PRESETS} (non-standard
couplings). Each preset carries its physical justification and bibliographic
citation (e.g.\ \cite{Esteban2020NuFit} for standard best-fit oscillation
parameters, \cite{Esteban2018LMADark} for the \emph{LMA-Dark} preset, and
\cite{IceCube2022NSI} for the IceCube DeepCore NSI bounds), so any numerical
result can be traced back to its origin in the literature.}
\end{implbox}

\begin{notebox}
Counting CP phases for the sterile extension requires care: an $N$-generation
Dirac mixing matrix has $(N-1)(N-2)/2$ independent physical phases. For
$N=4$ (the 3+1 case) this gives exactly \keyresult{three} Dirac phases
($\delta_{13}$, $\delta_{14}$, $\delta_{24}$); the rotation $R_{34}$ is
always real because the fourth phase can be absorbed into a redefinition of
the charged-lepton fields, and hence sterile presets do not include
\code{delta34\_deg} (see \cref{ch:core-bsm}).
\end{notebox}

% Restore the shared \modulehead (one-line italic description) for every
% chapter from here on.
\let\modulehead\savedmoduleheadchOne


% =============================================================================
\part{Propagation Core --- \texttt{core/}}
\chapter{\texttt{core.common} --- shared Hamiltonian, mixing, and probabilities}
\label{ch:core-common}

This chapter covers the physics infrastructure shared by \emph{every} flavour
scenario (the 3-flavour SM and its BSM extensions) and by \emph{every}
propagation medium. It is the conceptual core of the whole project: this is
where the factorisation trick that makes matter propagation cheap to compute
is defined.

\section{\texttt{potential.py} --- kinetic and matter potentials}

\modulehead{core/common/potential.py}{dimensionless (matter + kinetic)
potentials, GPU-first.}

\begin{theorybox}
The propagation Hamiltonian in natural units splits into a kinetic part
(vacuum oscillation) and a coherent matter-interaction part (the MSW
effect),
\begin{equation}
  H = H_{\rm kin} + H_{\rm mat}.
\end{equation}
\textbf{Kinetic term.} For a mass eigenstate $i$ with energy $E \gg m_i$, the
phase accumulated relative to the lightest mode after travelling a distance
$L$ is $\Delta m^2_{i1} L / (2E)$ (in natural units, then divided by
$\hbar c$ to convert to laboratory units). Introducing the evolution scale
$L_{\rm scale}$ (by default the Earth radius $R_\oplus$) and the
dimensionless coordinate $x=L/L_{\rm scale}$, the dimensionless kinetic
eigenvalue entering the Hamiltonian is
\begin{equation}
  \keyresult{k_i = L_{\rm scale}\,\frac{\Delta m^2_i}{2E}\,\frac{1}{\hbar c}}.
\end{equation}
\textbf{Charged- versus neutral-current matter potentials.} Coherent
forward scattering off the background medium proceeds through two distinct
weak-interaction channels, each contributing an independent term to
$H_{\rm mat}$ and computed by its own dedicated function:
\begin{itemize}
  \item \textbf{Charged current (CC).} Only $\nu_e$ ($\bar\nu_e$) can
    exchange a $W^\pm$ with the ambient electrons (there are no positrons,
    muons, or taus to scatter off in ordinary matter), so this channel
    singles out the electron flavour alone. It adds an effective
    Wolfenstein-type potential \parencite{Wolfenstein1978},
    \begin{equation}
      V_{CC} = \sqrt{2}\, G_F\, n_e,
    \end{equation}
    which only affects the $ee$ entry of the flavour-basis Hamiltonian.
    Made dimensionless with the scale $L_{\rm scale}$ and computed by
    \pyfunc{matter\_potential\_cc},
    \begin{equation}
      \keyresult{V_{CC} = \pm\sqrt{2}\,G_F\,n_e\,L_{\rm scale}}, \qquad (+\text{ for } \nu,\ -\text{ for } \bar\nu).
    \end{equation}
  \item \textbf{Neutral current (NC).} A $Z^0$ is exchanged with electrons,
    protons, \emph{and} neutrons alike, and identically across all three
    \emph{active} flavours (lepton universality). In electrically neutral
    matter the proton and electron NC contributions cancel exactly, leaving
    \begin{equation}
      V_{NC} = -\frac{1}{\sqrt2}\,G_F\,n_n,
    \end{equation}
    where $n_n$ is the neutron density. Made dimensionless the same way and
    computed by \pyfunc{matter\_potential\_nc},
    \begin{equation}
      \keyresult{V_{NC} = \mp\frac{1}{\sqrt2}\,G_F\,n_n\,L_{\rm scale}}, \qquad (-\text{ for } \nu,\ +\text{ for } \bar\nu).
    \end{equation}
\end{itemize}
For antineutrinos both potentials flip sign, $V \to -V$ (the weak
interaction has opposite sign for particle/antiparticle), which is why both
functions take the same \code{antinu} flag. In the pure 3-flavour
Standard Model, $V_{NC}\,\mathbb{1}_3$ is a common phase across all three
active flavours and therefore has \emph{no effect on any oscillation
probability}: this is why $V_{NC}$ never appears at all in the SM matter
Hamiltonian, and why $n_n$ is nowhere plumbed through the SM code path. The
3+1 sterile extension changes this qualitatively, since the sterile flavour
is a gauge singlet and \emph{does not} receive $V_{NC}$: subtracting the
same global phase $V_{NC}\,\mathbb{1}_4$ from
$\diag(V_{CC}+V_{NC},\,V_{NC},\,V_{NC},\,0)$ leaves a genuine relative term
on the sterile diagonal entry,
\begin{equation}
  \diag(V_{CC}+V_{NC},\,V_{NC},\,V_{NC},\,0) - V_{NC}\,\mathbb{1}_4
  = \diag(V_{CC},\,0,\,0,\,-V_{NC}),
\end{equation}
which \emph{does} affect sterile-active oscillation probabilities, with
magnitude comparable to $V_{CC}$ itself: for isoscalar matter (equal proton
and neutron density, e.g.\ Earth's mantle) $V_{NC}/V_{CC}=n_n/(2n_e)
\approx1/2$, independent of the absolute density. See
\texttt{hamiltonian.py} below for how this term enters the reduced
Hamiltonian, and \cref{ch:earth,ch:atmosphere} for how $n_n$ is exposed from
the Earth and atmosphere density profiles.
\end{theorybox}

\begin{implbox}
\begin{itemize}
  \item \pyfunc{kinetic\_potential(DeltamSq, E, evolution\_scale\_m, ...)}
    implements $k = L_{\rm scale}\,\Delta m^2/(2E\hbar c)$, accepting scalar
    or tensor inputs with automatic \code{device}/\code{dtype} inference.
  \item \pyfunc{matter\_potential\_cc(n\_e, antinu, evolution\_scale\_m, ...)}
    implements $V_{CC}=\pm\sqrt2 G_F n_e L_{\rm scale}$; the sign is set by
    the boolean flag \code{antinu}.
  \item \pyfunc{matter\_potential\_nc(n\_n, antinu, evolution\_scale\_m, ...)}
    implements $V_{NC}=\mp\tfrac{1}{\sqrt2} G_F n_n L_{\rm scale}$ (same
    \code{antinu} sign convention as the CC potential). It has no
    \code{legacy\_precision} counterpart (see below), since the legacy
    \texttt{peanuts} reference never modelled the sterile extension.
  \item Public API input units: $n_e$, $n_n$ in $\mathrm{mol\,cm^{-3}}$,
    $\Delta m^2$ in $\mathrm{eV^2}$, $E$ in $\mathrm{MeV}$,
    \code{evolution\_scale\_m} in metres (default $R_\oplus$, see
    \code{util.constant.R\_E}).
\end{itemize}
\end{implbox}

\begin{notebox}
\textbf{The \code{legacy\_precision} flag.} Both \pyfunc{matter\_potential\_cc}
and \pyfunc{kinetic\_potential} accept a boolean \code{legacy\_precision}
keyword, but only \pyfunc{matter\_potential\_cc} does anything with it: it
selects which numerical value of the matter-potential prefactor
$\sqrt2\,G_F\,n_e$ is used.
\begin{itemize}
  \item \code{legacy\_precision=False} (the default): the prefactor
    $\sqrt2\,G_F\,N_A\times10^{6}\,(\hbar c)^2$ is computed once, at import
    time, from the full double-precision constants in
    \code{util/constant.py} (see the table below).
  \item \code{legacy\_precision=True}: the prefactor is replaced by the
    single hardcoded value $3.868\times10^{-7}\ \mathrm{MeV\,m^2}$,
    reproduced from Eq.~4.17 of the legacy \texttt{peanuts} reference
    implementation \parencite{LucentePeanuts}, itself only quoted to 4
    significant digits. This mode exists exclusively to reproduce
    \texttt{peanuts} output bit-for-bit in the legacy validation tests
    (\code{medium/*/test/*legacy*}); it is \textbf{not} intended for
    physics production use, since it discards precision that is otherwise
    available at no cost.
  \item \pyfunc{kinetic\_potential} accepts the same keyword purely for
    call-site symmetry with \pyfunc{matter\_potential} (so both can be
    driven by one shared \code{legacy\_precision} switch); no
    legacy-rounded kinetic prefactor exists, so the flag is a no-op there.
\end{itemize}

\begin{center}
\small
\begin{tabular}{@{}llll@{}}
\toprule
Symbol & Description & Value & Units \\
\midrule
$G_F$      & Fermi coupling constant     & $1.1663787\times10^{-11}$ & $\mathrm{MeV^{-2}}$ \\
$N_A$      & Avogadro constant           & $6.02214076\times10^{23}$ & $\mathrm{mol^{-1}}$ \\
$\hbar$    & Reduced Planck constant     & $6.582119569\times10^{-22}$ & $\mathrm{MeV\,s}$ \\
$c$        & Speed of light (exact)      & $2.99792458\times10^{8}$  & $\mathrm{m\,s^{-1}}$ \\
$\hbar c$  & derived from $\hbar$, $c$   & $1.9732698\times10^{-13}$ & $\mathrm{MeV\,m}$ \\
\midrule
\multicolumn{4}{@{}l}{Matter-potential prefactor $\sqrt2\,G_F\,N_A\times10^6(\hbar c)^2$:} \\
\midrule
full precision & derived from the constants above & $3.86793\times10^{-7}$ & $\mathrm{MeV\,m^2}$ \\
legacy (Eq.~4.17,~\parencite{LucentePeanuts}) & hardcoded, 4 sig.\ digits & $3.868\times10^{-7}$ & $\mathrm{MeV\,m^2}$ \\
\bottomrule
\end{tabular}
\end{center}

\textbf{Precision difference.} The two prefactors differ by a relative
$1.85\times10^{-5}$ ($\sim 15$ parts per million) --- five orders of
magnitude larger than double-precision machine epsilon
($\approx2.2\times10^{-16}$). This confirms the discrepancy is not a
floating-point rounding artefact but a deliberate consequence of the legacy
reference value carrying only 4 significant digits; \code{legacy\_precision}
therefore trades a small, well-quantified accuracy loss for exact
reproducibility of the legacy \texttt{peanuts} numbers.
\end{notebox}

\begin{notebox}
\begin{itemize}
  \item \textbf{Ultrarelativistic kinetic term.} \emph{Formulation:} the
    mass-eigenstate energy is linearised as $E_i\approx E+\Delta m_i^2/(2E)$
    to obtain $k_i=L_{\rm scale}\,\Delta m_i^2/(2E\hbar c)$
    (\pyfunc{kinetic\_potential}). \emph{Justification and validity:}
    requires $E\gg m_i$ for every propagating mass eigenstate; for the
    solar/reactor/atmospheric energies simulated here
    ($E\gtrsim\mathrm{MeV}$) and sub-eV neutrino masses this holds to better
    than one part in $10^{12}$, so the linearisation is exact for every
    scenario in this project.
  \item \textbf{Coherent forward elastic scattering.} \emph{Formulation:}
    only the forward ($\theta=0$) elastic-scattering amplitude contributes
    to $V_{CC}$/$V_{NC}$ --- the standard Wolfenstein index-of-refraction
    argument \parencite{Wolfenstein1978}; incoherent (large-angle or
    inelastic) scattering is neglected entirely.
    \emph{Justification and validity:} holds whenever the neutrino's mean
    free path for incoherent scattering is far longer than the propagation
    baseline, true throughout the Sun, Earth, and atmosphere at the
    MeV--GeV energies simulated here.
  \item \textbf{Electrically neutral, unpolarised matter for $V_{NC}$.}
    \emph{Formulation:} $n_p=n_e$ is assumed when deriving
    $V_{NC}=-\tfrac1{\sqrt2}G_F n_n$, so the proton and electron NC
    contributions cancel exactly and only the neutron density survives; the
    medium is further assumed unpolarised and at rest, so no macroscopic
    spin- or bulk-velocity-dependent term is added to $H_{\rm mat}$.
    \emph{Justification and validity:} exact for ordinary,
    non-magnetised astrophysical matter --- the Sun, Earth, and atmosphere
    all satisfy local charge neutrality to extremely high precision, and
    carry no coherent spin polarisation on the scales relevant here.
  \item \textbf{Local, instantaneous density.} \emph{Formulation:} $V_{CC}$
    and $V_{NC}$ are evaluated from the density $n_e(x)$, $n_n(x)$ at the
    neutrino's instantaneous position, treating the medium as locally
    static and homogeneous on the scale of the local oscillation length.
    \emph{Justification and validity:} requires the density profile to vary
    slowly compared to that oscillation wavelength; the segmented-evolutor
    machinery (\texttt{evolutor.py} below, and the numerical and
    perturbative chapters, \cref{ch:core-numerical,ch:core-perturbative})
    exists precisely to control the residual error from this approximation
    by refining the segmentation where the density changes fastest.
\end{itemize}
\end{notebox}

\section{\texttt{pmns.py} --- shared mixing infrastructure}

\modulehead{core/common/pmns.py}{abstract class \pyclass{PMNS} and container
\pyclass{PMNSParams}, common to every flavour scenario.}

\begin{theorybox}
\subsection*{Mixing angles and the CP-violating phase}
The active-flavour sector is parametrised by three mixing angles and one CP
phase, all real (angles in radians) \parencite{ParticleDataGroup}:
\begin{itemize}
  \item $\theta_{12}$ (\textbf{solar} angle) sets the $(e,\mu)$-plane
    rotation; the mixing angle measured by solar and by KamLAND reactor
    data, and the largest of the three.
  \item $\theta_{13}$ (\textbf{reactor} angle) sets the $(e,\tau)$-plane
    rotation; the smallest angle, measured by short-baseline reactor
    experiments (Daya Bay, RENO, Double Chooz). Its non-zero value is what
    makes CP violation observable at all in the lepton sector.
  \item $\theta_{23}$ (\textbf{atmospheric} angle) sets the $(\mu,\tau)$-plane
    rotation; measured by atmospheric and long-baseline accelerator
    experiments, close to (possibly exactly) maximal, $\theta_{23}\approx\pi/4$.
  \item $\delta$ (\textbf{Dirac CP phase}) is the only source of CP
    violation carried by this parametrisation (Majorana phases, relevant
    only if neutrinos are Majorana particles, do not affect oscillation
    probabilities and are absent from it). $\delta$ has no physical,
    rephasing-invariant effect unless all three angles are non-zero; it
    enters oscillation probabilities through the Jarlskog invariant
    (\pyfunc{PMNS.jarlskog\_invariant}, see below).
\end{itemize}
These four scalars, together with a \pyclass{RuntimeContext} fixing the
target \code{device}/\code{dtype}, are exactly the contents of
\pyclass{PMNSParams} --- nothing else: the mass-squared splittings do not
enter the mixing matrix, so they are kept out of this container and live in
\pyclass{MassSpectrum} instead (next section).

\subsection*{Elementary rotation and phase generators}
Every mixing matrix in this project, SM or BSM, is assembled from two
elementary building blocks that \pyclass{PMNS} implements once, sized
generically by \code{self.n\_flavours} so the same formulas serve the
$3\times3$ active-only case and any larger active+sterile case (BSM
chapter). Restricted to the base $3\times3$ active sector, with the fixed
flavour order $(e,\mu,\tau)=(0,1,2)$, they read explicitly:
\begin{align}
  R_{12}(\theta_{12}) &=
    \begin{pmatrix}
      \cos\theta_{12} & \sin\theta_{12} & 0 \\
      -\sin\theta_{12} & \cos\theta_{12} & 0 \\
      0 & 0 & 1
    \end{pmatrix}
  & \text{rotation in the $(e,\mu)$ plane;}
  \\
  R_{13}(\theta_{13}) &=
    \begin{pmatrix}
      \cos\theta_{13} & 0 & \sin\theta_{13} \\
      0 & 1 & 0 \\
      -\sin\theta_{13} & 0 & \cos\theta_{13}
    \end{pmatrix}
  & \text{rotation in the $(e,\tau)$ plane;}
  \\
  R_{23}(\theta_{23}) &=
    \begin{pmatrix}
      1 & 0 & 0 \\
      0 & \cos\theta_{23} & \sin\theta_{23} \\
      0 & -\sin\theta_{23} & \cos\theta_{23}
    \end{pmatrix}
  & \text{rotation in the $(\mu,\tau)$ plane;}
  \\
  \Delta(\delta) &= \diag(1,\,1,\,e^{i\delta})
  & \text{CP phase, carried at the $\tau$ index.}
\end{align}
$R_{12}$, $R_{13}$, $R_{23}$ are always \textbf{real}: \pyclass{PMNS} never
attaches a phase to them (they call the shared generator with
\code{phase=None}); only $\Delta$ carries $\delta$, and only at position
$(\tau,\tau)$. This is a modelling choice of the parametrisation, not a
computational shortcut: the same underlying generator (\pyfunc{PMNS.\_rot})
accepts an arbitrary rotation plane $(i,j)$ and an arbitrary phase, which is
exactly what lets the 3+1 sterile extension reuse it unchanged for the
active-sterile rotations $R_{14}$, $R_{24}$, $R_{34}$ (each carrying its own
CP phase) without modifying this base class at all (BSM chapter). How
$R_{12}$, $R_{13}$, $R_{23}$, $\Delta$ combine into a full mixing matrix ---
the product order, and which rotations sit inside versus outside the block
that commutes with the matter potential --- is scenario-specific and is
given in the SM and BSM chapters.

\subsection*{Flavour, mass, and reduced bases}
Three bases appear throughout this project:
\begin{itemize}
  \item the \textbf{flavour basis} $\alpha=e,\mu,\tau(,s)$: the physical
    basis in which neutrinos are produced and detected, and the basis in
    which the matter potential is always diagonal. In the plain 3-flavour
    Standard Model, or the 3+1 sterile extension without its optional
    neutral-current term, only the electron entry is populated,
    $H_{\rm mat}=\diag(V_{CC},0,\dots,0)$; once that neutral-current term is
    active, $H_{\rm mat}=\diag(V_{CC},0,0,-V_{NC})$ carries a second,
    non-zero entry on the sterile diagonal (\texttt{potential.py} section
    above, \texttt{hamiltonian.py} section below) --- diagonal in every
    case, but not necessarily supported on the electron entry alone;
  \item the \textbf{mass basis} $i=1,\dots,n$: the eigenbasis of the
    free-propagation (kinetic-only) part of the Hamiltonian, in which the
    propagation phases $e^{-ik_ix}$ are defined;
  \item an intermediate \textbf{reduced basis}, introduced purely for
    computational convenience (\texttt{hamiltonian.py} section below).
\end{itemize}
Every concrete \pyclass{PMNS} scenario provides a unitary $U$
(\pyfunc{pmns\_matrix}) relating the flavour and mass bases: an operator
$A_{\rm flavour}$ and its mass-basis counterpart $A_{\rm mass}$ are related
by $A_{\rm flavour} = U\,A_{\rm mass}\,U^\dagger$. In addition, it
factorises this $U$ as
\begin{equation}
  U = O\,U_{\rm red},
\end{equation}
for a second unitary $U_{\rm red}$ (\pyfunc{reduced}, the \emph{reduced
mixing matrix}) and a third, scenario-specific unitary $O$
(\pyfunc{outer\_block}, the \emph{outer block}). This factorisation defines
the reduced basis: it is the basis reached from the flavour basis by
applying $O$ \emph{alone},
\begin{equation}
  A_{\rm flavour} = O\,A_{\rm reduced}\,O^\dagger,
  \qquad
  A_{\rm reduced} = O^\dagger\,A_{\rm flavour}\,O,
\end{equation}
implemented generically, for any scenario's $O$, by \pyfunc{flavour\_basis}
and \pyfunc{reduced\_basis} (see the implementation box below). $U_{\rm
red}$ is precisely the remaining rotation that carries a reduced-basis
operator the rest of the way to the true mass basis: combining
$A_{\rm flavour}=U\,A_{\rm mass}\,U^\dagger=O\,A_{\rm reduced}\,O^\dagger$
with $U=O\,U_{\rm red}$ gives the reduced-to-mass-basis transformation and
its inverse,
\begin{equation}
  A_{\rm reduced} = U_{\rm red}\,A_{\rm mass}\,U_{\rm red}^\dagger,
  \qquad
  A_{\rm mass} = U_{\rm red}^\dagger\,A_{\rm reduced}\,U_{\rm red},
\end{equation}
so the reduced basis coincides with the mass basis exactly when $O$ is the
identity (equivalently $U_{\rm red}=U$).
That $U$ factorises this way, and that \pyfunc{flavour\_basis}/%
\pyfunc{reduced\_basis} implement exactly the $O$-dressing above, is the
core.common contract every concrete \pyclass{PMNS} subclass must satisfy;
\emph{which} concrete $U$, $O$, $U_{\rm red}$ a given scenario provides is
scenario-specific and is given in the SM and BSM chapters.
\end{theorybox}

\begin{implbox}
\begin{itemize}
  \item \pyclass{PMNSParams}: immutable container for
    $\theta_{12},\theta_{13},\theta_{23},\delta$ and the
    \pyclass{RuntimeContext}.
  \item \pyclass{PMNS} (abstract \pyclass{torch.nn.Module}): builds the
    shared generators \pyfunc{R12}, \pyfunc{R13}, \pyfunc{R23},
    \pyfunc{Delta} above, sized generically by \code{self.n\_flavours} (3
    for the SM, 4 for the 3+1 sterile case), and implements the part of the
    public interface that is \emph{independent of the number of flavours
    and of the product structure}: \pyfunc{refresh}, \pyfunc{flavour\_basis},
    \pyfunc{reduced\_basis}, \pyfunc{select\_antinu} (see box below),
    \pyfunc{conjugate}, \pyfunc{transpose}, \pyfunc{dagger}, the
    reduced-matrix variants
    \pyfunc{reduced\_conjugate}/\pyfunc{reduced\_transpose}/%
    \pyfunc{reduced\_dagger}, and the rephasing-invariant
    \pyfunc{jarlskog\_invariant}.
  \item Concrete subclasses only need to supply \pyfunc{pmns\_matrix},
    \pyfunc{reduced}, and \pyfunc{outer\_block}: how the four generators
    above combine into the full and reduced mixing matrices is genuinely
    different per scenario, so those three methods remain abstract here.
  \item \pyfunc{flavour\_basis} / \pyfunc{reduced\_basis} apply the generic
    dressing $O(\cdot)O^\dagger$ / its inverse $O^\dagger(\cdot)O$, for
    \emph{any} scenario-specific outer block $O$ (\pyfunc{outer\_block}) and
    any reduced-basis operator (not only the Hamiltonian); why this
    reproduces the full-flavour result exactly is shown in the
    \texttt{hamiltonian.py} section below.
\end{itemize}
\end{implbox}

\begin{implbox}
\textbf{\pyfunc{select\_antinu(U, antinu)}.} Applies the standard
antineutrino prescription to a full or reduced mixing matrix: returns $U$
unchanged for neutrinos and its complex conjugate $U^*$ for antineutrinos
(the sign flip of the matter potential, $V\to-V$, is handled separately by
\pyfunc{potential.matter\_potential}, not here). \code{antinu} may be a
plain Python \code{bool} (applies globally) or a boolean tensor broadcast
against the batch dimensions of $U$ via \code{torch.where}, so a single
batched call can mix neutrinos and antineutrinos within the same tensor ---
e.g.\ when propagating an atmospheric flux that contains both. Every mixing
matrix exposed by \pyclass{PMNS} (\pyfunc{pmns\_matrix}, \pyfunc{reduced},
\pyfunc{outer\_block}) routes through this single function, so the
neutrino/antineutrino convention is enforced in exactly one place.
\end{implbox}

\section{\texttt{mass\_spectrum.py} --- mass-squared splittings}

\modulehead{core/common/mass\_spectrum.py}{abstract class \pyclass{MassSpectrum}:
the two measurable mass-squared splittings shared by every oscillation
model.}

\begin{theorybox}
Oscillation probabilities depend on mass-squared \emph{differences}, never
on absolute masses, so the mass sector is parametrised by two independent,
directly measurable splittings \parencite{ParticleDataGroup}:
\begin{itemize}
  \item $\Delta m^2_{21}\equiv m_2^2-m_1^2$ (\textbf{solar} splitting): by
    convention always positive ($m_2>m_1$); measured by solar and KamLAND
    reactor data, $\Delta m^2_{21}\approx7.4\times10^{-5}\ \mathrm{eV^2}$.
    Sets the ``slow'' oscillation frequency.
  \item $\Delta m^2_{3\ell}$ (\textbf{atmospheric} splitting):
    $\Delta m^2_{31}\equiv m_3^2-m_1^2$ if $m_3$ is the heaviest state
    (\textbf{normal ordering}, $\ell=1$), or
    $\Delta m^2_{32}\equiv m_3^2-m_2^2$ if $m_3$ is the lightest
    (\textbf{inverted ordering}, $\ell=2$); measured by atmospheric and
    long-baseline accelerator data,
    $|\Delta m^2_{3\ell}|\approx2.5\times10^{-3}\ \mathrm{eV^2}$. Sets the
    ``fast'' oscillation frequency.
\end{itemize}
These are exactly the two fields stored by \pyclass{MassSpectrum}: no
absolute mass scale, and no third independent splitting (since
$\Delta m^2_{32}=\Delta m^2_{31}-\Delta m^2_{21}$ by construction, only two
of the three pairwise splittings are independent).
\end{theorybox}

\begin{notebox}[title={Normal and Inverted Ordering}]
Whether $m_3$ is the heaviest or the lightest mass eigenstate --- the
\textbf{neutrino mass ordering} --- is not yet known experimentally, and is
not fixed by any mixing angle: it is encoded entirely in the \emph{sign} of
$\Delta m^2_{3\ell}$, which \pyclass{MassSpectrum} takes as a single signed
input rather than as a separate ordering flag. Concretely,
\pyfunc{difference\_vector\_base()} builds the reduced (global-phase-removed)
three-component vector
\begin{equation}
  \bigl(\Delta m^2_1,\ \Delta m^2_2,\ \Delta m^2_3\bigr)
  = \begin{cases}
      \bigl(0,\ \Delta m^2_{21},\ \Delta m^2_{3\ell}\bigr) & \Delta m^2_{3\ell} > 0\ \text{(normal ordering, NO)},\\[2pt]
      \bigl(-\Delta m^2_{21},\ 0,\ \Delta m^2_{3\ell}\bigr) & \Delta m^2_{3\ell} < 0\ \text{(inverted ordering, IO)}.
    \end{cases}
\end{equation}
Both branches use the same stored $\Delta m^2_{21}>0$ and the same
$|\Delta m^2_{3\ell}|$: only the sign of the single scalar
$\Delta m^2_{3\ell}$ passed in switches the ordering, there is no separate
boolean or string argument anywhere in the API. This is a common source of
confusion when porting values from sources that specify ordering
explicitly: the sign must be set on $\Delta m^2_{3\ell}$ itself before it
reaches \pyclass{MassSpectrum}. The batched form evaluates both branches and
selects per-element with \code{torch.where}, so a single tensor of
$\Delta m^2_{3\ell}$ values with mixed signs is handled correctly in one
call.
\end{notebox}

\begin{implbox}
\begin{itemize}
  \item \pyclass{MassSpectrum} (abstract \pyclass{dataclass}): stores
    $\Delta m^2_{21}$, $\Delta m^2_{3\ell}$ as its only fields.
  \item \pyfunc{difference\_vector\_base()}: concrete, shared by every
    scenario; builds the 3-component vector above.
  \item \pyfunc{difference\_vector()}: abstract; sizes the base vector to
    the concrete scenario's flavour count --- returned unchanged for the
    plain 3-flavour case (\pyclass{MassSpectrum\_SM}, SM chapter), with a
    fourth entry $\Delta m^2_{41}$ appended for the 3+1 sterile case
    (\pyclass{MassSpectrum\_BSM}, BSM chapter). This split exists because
    sizing is scenario-specific while the active-sector vector itself is
    not.
  \item Consumed by \pyfunc{kinetic\_eigenvalue\_vector} (\texttt{potential.py})
    to build the dimensionless kinetic phases $k_i$, and bundled into
    \pyclass{OscillationParameters} as the \code{mass\_spectrum} field
    (next section).
\end{itemize}
\end{implbox}

\section{\texttt{oscillation.py} --- bundled oscillation state}

\modulehead{core/common/oscillation.py}{\pyclass{OscillationParameters}: an
immutable container grouping \code{pmns}, \code{mass\_spectrum},
\code{antinu}, and the optional \code{nsi} coupling.}

\begin{implbox}
This introduces no new physics: it replaces the previous pattern of passing
\code{(pmns, DeltamSq21, DeltamSq3l, antinu)} as independent arguments
through the entire propagation stack. \pyclass{OscillationParameters} is a
passive domain container: it does not build itself, and it never imports the
concrete SM/BSM/preset implementations it bundles. Building a fully
assembled \pyclass{OscillationParameters} from a named preset --- choosing
the concrete PMNS/mass-spectrum implementation and reading the parameter
values --- is done by
\code{PropagationConfig.oscillation\_parameters\_from\_preset} in
\code{tpeanuts/config/propagation.py}, from the preset registry in
\code{tpeanuts/config/presets.py}: that composition layer is intentionally
kept outside this module, which stays a plain container.
\end{implbox}

\begin{center}
\small
\begin{tabular}{@{}llll@{}}
\toprule
Attribute & Type & Default & Role \\
\midrule
\code{pmns} & \pyclass{PMNS} & required & \parbox[t]{4.4cm}{Built mixing object (\pyclass{PMNS\_SM} or \pyclass{PMNS\_sterile}).} \\[2pt]
\code{mass\_spectrum} & \pyclass{MassSpectrum} & required & \parbox[t]{4.4cm}{Built splitting object (\pyclass{MassSpectrum\_SM} or \pyclass{MassSpectrum\_BSM}).} \\[2pt]
\code{antinu} & \code{bool}/\code{Tensor} & \code{False} & \parbox[t]{4.4cm}{Antineutrino selector ($U\to U^*$, $V\to-V$).} \\[2pt]
\code{nsi} & \pyclass{NSIConfig} or \code{None} & \code{None} & \parbox[t]{4.4cm}{Optional NSI coupling; \code{None} means the plain SM matter potential.} \\[2pt]
\code{preset\_name} & \code{str} & \code{"custom"} & \parbox[t]{4.4cm}{Name of the preset used to build this object.} \\[2pt]
\code{ordering} & \code{str} & \code{""} & \parbox[t]{4.4cm}{Mass-ordering label carried by the preset (\code{"NO"}/\code{"IO"}).} \\[2pt]
\code{label} & \code{str} & \code{""} & \parbox[t]{4.4cm}{Short identifier string carried by the preset.} \\[2pt]
\code{description} & \code{str} & \code{""} & \parbox[t]{4.4cm}{Human-readable description/reference carried by the preset.} \\[2pt]
\midrule
\multicolumn{4}{@{}l}{Derived BSM-extension flags (read-only \code{@property}, not stored fields):} \\
\midrule
\code{BSM\_extension\_sterile} & \code{bool} & from \code{pmns} & \parbox[t]{4.4cm}{True iff \code{pmns.n\_flavours == 4} (a 3+1 sterile object).} \\[2pt]
\code{BSM\_extension\_NSI} & \code{bool} & from \code{nsi} & \parbox[t]{4.4cm}{True iff \code{nsi is not None}.} \\[2pt]
\code{BSM\_extension} & \code{bool} & \parbox[t]{1.6cm}{OR of the two above} & \parbox[t]{4.4cm}{True iff \emph{any} BSM extension is active; the routing flag checked by \code{core.BSM}, \code{core.numerical}, and \code{core.perturbative} to pick the SM-only fast path or the general BSM path.} \\[2pt]
\bottomrule
\end{tabular}
\end{center}

\begin{notebox}
The three \code{BSM\_extension*} flags are \emph{derived}, never stored:
they are recomputed on every access from \code{pmns}/\code{nsi} rather than
set at construction time, so they can never drift out of sync with the
objects they describe (there is no way to end up with
\code{BSM\_extension\_sterile=True} while \code{pmns} is actually a
3-flavour \pyclass{PMNS\_SM}). The design rationale for keeping
\code{mass\_spectrum} separate from \code{pmns} is that the PMNS mixing
matrix does not depend on $\Delta m^2_{21}/\Delta m^2_{3\ell}$ (those
parameters only enter through the kinetic term, previous sections), so
storing them on \pyclass{PMNSParams} as well would duplicate the same
values.
\end{notebox}

\begin{notebox}[title={Default Configuration}]
Unless a different preset is requested, the reference configuration used
throughout this project's examples and tests is the NuFIT 5.2 (2022)
Standard Model best fit with normal ordering, registered as
\code{"\_SM\_NUFIT52\_NO"} in
\code{tpeanuts.config.presets.OSCILLATION\_PRESETS} \parencite{Esteban2020NuFit,NuFit52}.
Built via
\code{PropagationConfig.oscillation\_parameters\_from\_preset("\_SM\_NUFIT52\_NO")},
it populates \pyclass{OscillationParameters} as follows:

\begin{center}
\small
\begin{tabular}{@{}lll@{}}
\toprule
Attribute & Default value & Note \\
\midrule
$\theta_{12}$ & $33.41^\circ$ & solar angle \\
$\theta_{13}$ & $8.58^\circ$ & reactor angle \\
$\theta_{23}$ & $49.0^\circ$ & atmospheric angle \\
$\delta$ & $197.0^\circ$ & Dirac CP phase \\
$\Delta m^2_{21}$ & $7.41\times10^{-5}\ \mathrm{eV^2}$ & solar splitting \\
$\Delta m^2_{3\ell}$ & $+2.511\times10^{-3}\ \mathrm{eV^2}$ & atmospheric splitting (positive) \\
\code{antinu} & \code{False} & neutrino, not antineutrino \\
\code{ordering} & \code{"NO"} & normal ordering, set by the positive sign above \\
\code{BSM\_extension\_sterile} & \code{False} & \code{pmns} is 3-flavour, not sterile \\
\code{BSM\_extension\_NSI} & \code{False} & \code{nsi = None} \\
\code{BSM\_extension} & \code{False} & no BSM extension active \\
\bottomrule
\end{tabular}
\end{center}
\end{notebox}

\section{\texttt{hamiltonian.py} --- the propagation Hamiltonian}

\modulehead{core/common/hamiltonian.py}{the single Hamiltonian-assembly
module in the project: correct and complete for the 3-flavour SM, the 3+1
sterile extension, Non-Standard Interactions, or any combination, with no
separate SM-only implementation.}

\begin{theorybox}
\subsection*{General construction of the flavour Hamiltonian}
In the flavour basis (\texttt{pmns.py} section above), the propagation
Hamiltonian is, in full generality, the mass-basis kinetic term rotated
into the flavour basis by the full mixing matrix $U$, plus whatever matter
term $H_{\rm mat}$ the active scenario contributes:
\begin{equation}
  H_{\rm flavour} = U\,\diag(k_1,\dots,k_n)\,U^\dagger + H_{\rm mat}.
\end{equation}
This decomposition is always correct, for \emph{any} unitary $U$ and
\emph{any} Hermitian $H_{\rm mat}$, with no assumption on how $U$
factorises or on what $H_{\rm mat}$ concretely is: it is simply the
Schr\"odinger equation written in the flavour basis. What \emph{does}
depend on the scenario is the concrete content of $H_{\rm mat}$
(\texttt{potential.py} section above): plain 3-flavour Standard Model, or
3+1 sterile \emph{without} its optional neutral-current term,
$H_{\rm mat}=\diag(V_{CC},0,\dots,0)$, populated on the electron entry
alone; with Non-Standard Interactions and/or the sterile neutral-current
term active, $H_{\rm mat}$ acquires extra diagonal and/or off-diagonal
entries beyond the electron one (full expression below, in ``The 3+1
sterile neutral-current term'' subsection). Taken at face value, evolving
$H_{\rm flavour}$ is expensive in every case: since $U$ and $H_{\rm mat}$ do
not commute in general, $H_{\rm flavour}$ must be re-diagonalised as a
genuinely complex $n\times n$ matrix at every point along the path where
the electron (and, when the NC term is active, neutron) density changes.

\subsection*{The reduction trick: factorising out the outer block}
Every concrete \pyclass{PMNS} scenario additionally factorises its mixing
matrix as $U = O\,U_{\rm red}$ (\texttt{pmns.py} section above), where the
scenario-specific \textbf{outer block} $O$ (\pyfunc{outer\_block}) is
constructed, in every scenario implemented so far, to leave the electron
flavour entry fixed, and therefore to commute exactly with the
\emph{plain charged-current} matter term:
\begin{equation}
  O\,\diag(V_{CC},0,\dots,0)\,O^\dagger = \diag(V_{CC},0,\dots,0).
\end{equation}
This commutation is a structural property of $O$ alone, and holds for
\emph{any} scenario's $O$ and $U_{\rm red}$ --- but, as stressed below, only
for this plain CC-only piece of $H_{\rm mat}$, not for $H_{\rm mat}$ in
general. Whenever $H_{\rm mat}$ reduces to this plain CC-only form (i.e.\
neither NSI nor the sterile neutral-current term is active), the general
expression above simplifies to
\begin{align}
  H_{\rm flavour} &= U \diag(k_1,\dots,k_n)\, U^\dagger + \diag(V_{CC},0,\dots,0) \nonumber\\
  &= O\,U_{\rm red}\diag(k)\,U_{\rm red}^\dagger\,O^\dagger + O\,\diag(V_{CC},0,\dots,0)\,O^\dagger \nonumber\\
  &= O \underbrace{\Bigl[U_{\rm red}\diag(k)\,U_{\rm red}^\dagger + \diag(V_{CC},0,\dots,0)\Bigr]}_{\displaystyle H_{\rm red}}\, O^\dagger
  \;=\; O\, H_{\rm red}\, O^\dagger.
  \label{eq:H-factorization}
\end{align}
\textbf{Computational consequence (the peanuts ``reduction trick'').} The
$n$-flavour evolution problem reduces to: (1) solving the evolution of the
smaller, cheaper reduced Hamiltonian $H_{\rm red}$, obtaining
$S_{\rm red}(x)$; (2) \emph{dressing} the result once with the constant
matrix $O$ (which depends on neither density nor $x$),
$S_{\rm flavour}(x) = O\, S_{\rm red}(x)\, O^\dagger$. This is exactly what
\pyfunc{pmns.flavour\_basis} does: the spectral calculations that
diagonalise the propagation Hamiltonian only ever need to act on
$H_{\rm red}$, never on $H_{\rm flavour}$. This shortcut is \emph{not}
available once $H_{\rm mat}$ carries an NSI or sterile neutral-current
contribution, since neither commutes with $O$ in general (see below, and
the dedicated implementation box on the explicit BSM $H_{\rm mat}$
construction).

\subsection*{Assembling the reduced Hamiltonian}
Given the reduced mass-squared vector $(\Delta m^2_1,\dots,\Delta m^2_n)$
from \pyclass{MassSpectrum} (\texttt{mass\_spectrum.py} section above) and
the reduced mixing matrix $U_{\rm red}$ of the active \pyclass{PMNS}
scenario, the reduced Hamiltonian is
\begin{equation}
  \keyresult{H_{\rm red} = U_{\rm red}\,\diag(k_1,\dots,k_n)\,U_{\rm red}^{\dagger} + \diag(V_{CC},0,\dots,0)}.
\end{equation}
The matter term populates only the electron-flavour entry, by construction
of $H_{\rm mat}$ (\texttt{potential.py} section above): this holds for any
flavour count $n$, so no special handling is needed as $n$ grows --- every
additional flavour beyond the first simply carries a zero matter coupling.
When Non-Standard Interactions are active, the matter term
generalises to $V_{CC}\,(\diag(1,0,\dots,0)+\varepsilon)$ for a Hermitian
coupling matrix $\varepsilon$, built in the flavour basis and rotated back
into the reduced basis with $O^\dagger(\cdot)O$.

\subsection*{The 3+1 sterile neutral-current term, and why it breaks the
reduction trick}
Optionally supplying a neutron density $n_n$ for a 4-flavour \pyclass{PMNS}
scenario (\texttt{potential.py} section above) adds the relative
neutral-current term derived there to the flavour-basis matter Hamiltonian,
\begin{equation}
  H_{\rm mat} = \diag(V_{CC},0,0,0) + V_{CC}\,\varepsilon - \diag(0,0,0,V_{NC}),
\end{equation}
(the NSI term present only when $\varepsilon$ is set). Crucially, the outer
block $O$ of the 3+1 sterile scenario mixes the $\{\mu,\tau,s\}$ block
(\cref{ch:core-bsm}) and therefore does \emph{not} commute with a term that
singles out the sterile diagonal entry:
\begin{equation}
  O\,\diag(0,0,0,V_{NC})\,O^\dagger \neq \diag(0,0,0,V_{NC}).
\end{equation}
This means the reduction-trick shortcut of the previous subsection ---
skipping the $O^\dagger(\cdot)O$ rotation entirely when $H_{\rm mat}$ is
already diagonal in the reduced basis --- is only valid when \emph{neither}
NSI \emph{nor} the neutral-current term is active. Whenever either is
present, \pyfunc{hamiltonian\_matter\_reduced} always takes the general
path: build $H_{\rm mat}$ in the flavour basis (as above) and rotate it,
$H_{\rm mat}^{\rm red} = O^\dagger\,H_{\rm mat}\,O$. Omitting $n_n$ (the
default) recovers the CC-only sterile matter term used throughout the rest
of this thesis, unchanged; the term has no effect at all for a 3-flavour
\pyclass{PMNS} (there $V_{NC}$ is the unobservable global phase discussed in
\texttt{potential.py}), so passing $n_n$ to a 3-flavour scenario is silently
ignored rather than raising.
\end{theorybox}

\begin{implbox}
\begin{itemize}
  \item \pyfunc{kinetic\_eigenvalue\_vector(oscillation, E, ...)}
    (\code{core/common/potential.py}) converts the mass spectrum's
    \pyfunc{difference\_vector()} (previous section) into the dimensionless
    kinetic phases $k_i$ via \pyfunc{kinetic\_potential}.
  \item \pyfunc{hamiltonian\_kinetic\_reduced(...)} implements
    $U_{\rm red}\diag(k)\,U_{\rm red}^\dagger$ for any flavour count, always
    using the Hermitian conjugate.
  \item \pyfunc{hamiltonian\_matter\_reduced(oscillation, n\_e\_mol\_cm3, *,
    n\_n\_mol\_cm3=None, ...)} implements $\diag(V_{CC},0,\dots,0)$ (no
    rotation) when neither NSI nor a genuine sterile NC term is requested,
    or the general rotated construction above otherwise. The
    \code{n\_n\_mol\_cm3} keyword is optional and only meaningful for a
    4-flavour \code{oscillation.pmns}.
  \item \pyfunc{\_matter\_direction\_reduced(oscillation, *, antinu,
    include\_nc, context)} (module-private): factors the reduced matter
    Hamiltonian into $H_{\rm mat}^{\rm red}=V_{CC}\,P_{CC} + V_{NC}\,P_{NC}$,
    returning the two density-independent ``direction'' matrices
    $P_{CC},P_{NC}$ on their own. Shared by
    \pyfunc{hamiltonian\_matter\_reduced} and by the perturbative
    first-order correction (\cref{ch:core-perturbative}), which needs the
    direction matrices in isolation to sandwich the spectral projectors.
  \item \pyfunc{hamiltonian\_reduced(...)} / \pyfunc{hamiltonian\_flavour(...)}
    sum both terms directly from $n_e$ (and optional $n_n$) and the
    splittings; \pyfunc{hamiltonian\_flavour} additionally applies the
    $O(\cdot)O^\dagger$ dressing via \pyfunc{pmns.flavour\_basis}.
\end{itemize}
\end{implbox}

\begin{implbox}[title={Explicit construction of \texorpdfstring{$H_{\rm mat}$}{H\_mat} under BSM extensions}]
\pyfunc{hamiltonian\_matter\_reduced} never rebuilds $H_{\rm mat}$ term by
term at every call; it factorises it into position-\emph{independent}
``direction'' matrices and position-\emph{dependent} scalar potentials, via
\pyfunc{\_matter\_direction\_reduced}:
\begin{equation}
  H_{\rm mat}^{\rm red} = V_{CC}\,P_{CC} \; [+\; V_{NC}\,P_{NC}].
\end{equation}
Concretely, in code order:
\begin{itemize}
  \item \textbf{Neither NSI nor the sterile NC term requested} (the common
    case). Fast path: $P_{CC}=\diag(1,0,\dots,0)$ directly in the reduced
    basis, with \emph{no} rotation, since this direction is exactly
    $O$-invariant (\labelcref{eq:H-factorization} above); $P_{NC}$ is not
    built at all.
  \item \textbf{NSI and/or the sterile NC term requested.} General path:
    build the flavour-basis direction matrices
    $D_{\rm flavour}=\diag(1,0,\dots,0)+\varepsilon$ (with $\varepsilon=0$
    if NSI is off) and, only if the NC term is requested,
    $D_{n,\rm flavour}=-\diag(0,\dots,0,1)$ (the sterile entry); rotate both
    into the reduced basis with the scenario's outer block $O$,
    \begin{equation}
      P_{CC} = O^\dagger\,D_{\rm flavour}\,O,
      \qquad
      P_{NC} = O^\dagger\,D_{n,\rm flavour}\,O.
    \end{equation}
  \item \pyfunc{hamiltonian\_matter\_reduced} then evaluates the
    position-dependent scalar potentials --- $V_{CC}$ always, via
    \pyfunc{matter\_potential\_cc}, and, only when \code{n\_n\_mol\_cm3} is
    supplied to a 4-flavour scenario, $V_{NC}$ via
    \pyfunc{matter\_potential\_nc} --- and assembles
    $H_{\rm mat}^{\rm red}=V_{CC}\,P_{CC}\,[+\,V_{NC}\,P_{NC}]$ above.
\end{itemize}
This split is exactly what lets \pyfunc{evolutor\_first\_order}
(\cref{ch:core-perturbative}) reuse $P_{CC}$, $P_{NC}$ directly, sandwiching
them against the spectral projectors without recomputing the rotation at
every energy or density sample.
\end{implbox}

\begin{notebox}
The reduced kinetic term is always built as $U_{\rm red}\diag(k)\,U_{\rm
red}^\dagger$ (the Hermitian conjugate, never the plain transpose): this is
the only form for which \pyfunc{pmns.flavour\_basis} applied to it
reproduces $U_{\rm full}\diag(k)\,U_{\rm full}^\dagger$ exactly, since
$U_{\rm full}=O\,U_{\rm red}$ per \labelcref{eq:H-factorization}. Whenever
$U_{\rm red}$ happens to be real for a given scenario, the Hermitian
conjugate and the transpose coincide there numerically; the distinction
only becomes observable once $U_{\rm red}$ is genuinely complex. This is,
in fact, exactly what happened in the original \texttt{peanuts}
implementation \parencite{LucentePeanuts}: since it only ever considered
the plain 3-flavour Standard Model, its reduced mixing matrix was always
real, so using the plain transpose there was correct, not merely
convenient -- the two conventions were indistinguishable in that restricted
setting. It is only once \texttt{tpeanuts} generalises to a genuinely
complex $U_{\rm red}$ (the 3+1 sterile extension with $\delta_{14}\neq0$,
\cref{ch:core-bsm}) that the plain transpose silently becomes wrong, which
is why the code always uses the Hermitian conjugate unconditionally rather
than branching on whether $U_{\rm red}$ happens to be real.
\end{notebox}

\section{\texttt{evolutor.py} (generic) --- operator composition}

\modulehead{core/common/evolutor.py}{applying an evolutor to a state and
composing ordered segment evolutors.}

The concrete construction of a single segment's evolutor $S(x_{k+1},x_k)$
--- closed-form, perturbative, or fully numerical --- is model-dependent
and is given in the corresponding SM and BSM chapters
(\cref{ch:core-sm,ch:core-bsm,ch:core-numerical,ch:core-perturbative}); this
module only composes and applies whatever evolutors those chapters produce.

\begin{theorybox}
\subsection*{Composing segment evolutors}
If the propagation path is split into $n$ consecutive legs
$x_0 \to x_1 \to \dots \to x_n$, the composition property of the
position-dependent Schr\"odinger equation gives
\begin{equation}
  \keyresult{S(x_n, x_0) = S(x_n,x_{n-1})\cdots S(x_2,x_1)\,S(x_1,x_0)},
\end{equation}
i.e.\ the total evolution operator is the \emph{ordered} product (most
recent leg on the left) of the per-leg operators. This identity is what
allows any continuous density profile to be treated as a sequence of
locally solvable segments, each with its own closed-form or numerically
integrated evolutor.

\subsection*{Applying the evolutor to a state}
Once the total evolutor $S(x_n,x_0)$ is known, it is applied to an initial
state as a plain matrix--vector product,
\begin{equation}
  \keyresult{\psi(x_n) = S(x_n,x_0)\,\psi(x_0)},
\end{equation}
implemented directly by \pyfunc{apply\_evolutor\_to\_state}. What ``the
state'' $\psi$ represents, and how a transition probability is read off
from the result, splits into two cases:
\begin{itemize}
  \item \textbf{Coherent propagation.} $\psi(x_0)$ is a genuine flavour
    amplitude vector (a quantum superposition carrying relative phase,
    e.g.\ a single produced flavour, $\psi(x_0)=e_\alpha$). The transition
    amplitude is obtained by applying the formula above directly, and the
    transition probability to flavour $\beta$ follows immediately as
    \begin{equation}
      P_\beta = |\psi_\beta(x_n)|^2 = \bigl|(S\,\psi(x_0))_\beta\bigr|^2.
    \end{equation}
  \item \textbf{Incoherent propagation.} The initial ensemble is a
    classical mixture of \emph{mass} eigenstates with weights
    $w_i\geq0$, $\sum_i w_i=1$, carrying no relative phase between them
    (e.g.\ neutrinos that lost coherence over a very long baseline). There
    is no single amplitude vector $\psi(x_0)$ to apply $S$ to; instead each
    mass eigenstate is propagated and projected onto the final flavour
    separately, and the resulting \emph{probabilities} --- not amplitudes
    --- are summed with the classical weights,
    \begin{equation}
      P_\beta = \sum_i \bigl|(S\,U_{\rm PMNS})_{\beta i}\bigr|^2\, w_i.
    \end{equation}
\end{itemize}
Both formulas, together with the unitarity check that verifies them, are
derived and implemented in full in the \texttt{probability.py} section
below.
\end{theorybox}

\begin{implbox}
\begin{itemize}
  \item \pyfunc{apply\_evolutor\_to\_state(S, psi)}: applies a scalar or
    batched evolutor to a state vector, with automatic broadcasting.
  \item \pyfunc{compose\_segment\_evolutors(segments)}: performs the ordered
    product above; internally uses a binary-tree reduction
    (\code{util.math.tree\_reduce\_matmul}) to minimise the depth of the
    matrix-multiplication chain on GPU.
\end{itemize}
\end{implbox}

\section{\texttt{probability.py} --- from evolutors to observable probabilities}

\modulehead{core/common/probability.py}{the single transition-probability
convention for the whole project.}

\begin{theorybox}
Given the flavour evolutor $S$ (row = final flavour, column = initial
flavour),
\begin{equation}
  P(\nu_\alpha\to\nu_\beta) = \bigl|S_{\beta\alpha}\bigr|^2.
\end{equation}
For a coherent initial state $\psi_\alpha$ (flavour superposition with
physical phase), the final flavour probability is
\begin{equation}
  P_\beta = \Bigl|\sum_\alpha S_{\beta\alpha}\psi_\alpha\Bigr|^2
          = \bigl|(S\psi)_\beta\bigr|^2.
\end{equation}
For an incoherent mass mixture with weights $w_i\geq0$, $\sum_i w_i=1$
(no relative-phase information between mass eigenstates, the regime selected
by \code{massbasis=True} throughout the code),
\begin{equation}
  P_\beta = \sum_i \bigl|(S\,U_{\rm PMNS})_{\beta i}\bigr|^2\, w_i .
\end{equation}
Probability conservation: since $S$ is unitary, $\sum_\beta P_{\beta\alpha}=1$
for every column $\alpha$; this is used as a numerical consistency check,
not as an additional physics ingredient.
\end{theorybox}

\begin{implbox}
\begin{itemize}
  \item \pyfunc{probability\_transition(S, ...)}: the full probability
    matrix, or a selected channel (transition or survival).
  \item \pyfunc{probability\_coherent\_state(psi)} / \pyfunc{probability\_coherent(S, psi)}:
    implement the coherent formula above.
  \item \pyfunc{probability\_incoherent(...)}: implements the incoherent
    formula, either from a precomputed probability matrix or by projecting
    mass weights through $S$ and $U_{\rm PMNS}$.
  \item \pyfunc{probability\_state(...)}: dispatches between both
    modes depending on the type of initial state received. Every medium's
    \code{[medium]\_probability\_state} (\cref{ch:vacuum,ch:solar,ch:earth,ch:atmosphere})
    delegates to this function.
  \item \pyfunc{normalize\_probability\_columns(P, eps=1e-15)}: rescales
    each column of a probability matrix to sum exactly to one, for callers
    that need a strictly conserved matrix after a numerically imperfect
    construction (e.g.\ a truncated spectral sum); \code{eps} guards the
    division against a vanishing column sum.
  \item \pyfunc{check\_probability\_conservation} / \pyfunc{check\_probability\_matrix}:
    numerical unitarity checks, used extensively in the tests
    (\code{core/BSM/test}, \code{core/common/test}).
\end{itemize}
\end{implbox}

\begin{theorybox}
\subsection*{Averaging a probability over a coordinate}
A probability is dimensionless and normalised, so averaging it over any
coordinate (energy, an arbitrary production variable, or the full zenith
range) must \emph{preserve} that normalisation: the natural reduction is a
\textbf{weighted average}, not a plain integral. Given a weight $w$ over a
grid $x$ (an explicit production spectrum for energy, an arbitrary
caller-supplied weight for any other coordinate, or the geometric
solid-angle element for angle),
\begin{equation}
  \keyresult{\langle P\rangle = \frac{\int P(x)\,w(x)\,\mathrm{d}x}{\int w(x)\,\mathrm{d}x}},
\end{equation}
evaluated by the trapezoidal rule, so $\langle P\rangle$ stays in $[0,1]$
regardless of how $w$ is normalised. This is the defining difference from
\texttt{flux.py} below: a flux is not itself normalised, so its energy
reduction is a plain (unnormalised) integral, not a weighted average. Like
every other function in this module, all three reductions below accept a
final flavour dimension of 3 or 4, the latter for the 3+1 sterile
extension.
\end{theorybox}

\begin{implbox}
Three functions implement the weighted-average pattern above, one generic
and two coordinate-specific:
\begin{center}
\small
\begin{tabular}{@{}lll@{}}
\toprule
Function & Weight $w$ & Formula implemented \\
\midrule
\parbox[t]{5.0cm}{\pyfunc{probability\_\allowbreak weighted\_\allowbreak average(\allowbreak P, coordinate, weight, dim=-2)}} &
  \parbox[t]{2.6cm}{arbitrary, caller-supplied} &
  \parbox[t]{5.6cm}{$\displaystyle \langle P\rangle=\frac{\int P(x)\,w(x)\,dx}{\int w(x)\,dx}$
  --- the generic primitive, used e.g.\ for production-height averages
  (\cref{ch:atmosphere})} \\[10pt]
\parbox[t]{5.0cm}{\pyfunc{probability\_integrated(\allowbreak P, E\_grid\_MeV, spectrum, energy\_dim=-2)}} &
  \parbox[t]{2.6cm}{production spectrum $w(E)$, \textbf{required} (no default)} &
  \parbox[t]{5.6cm}{$\displaystyle \langle P\rangle=\frac{\int P(E)\,w(E)\,dE}{\int w(E)\,dE}$} \\[10pt]
\parbox[t]{5.0cm}{\pyfunc{probability\_\allowbreak integrated\_\allowbreak angular(\allowbreak P, theta\_deg, angular\_dim=-2)}} &
  \parbox[t]{2.6cm}{$w(\theta)=\sin\theta$, fixed and geometric} &
  \parbox[t]{5.6cm}{$\displaystyle \langle P\rangle=\frac{\int P(\theta)\sin\theta\,d\theta}{\int \sin\theta\,d\theta}$} \\
\bottomrule
\end{tabular}
\end{center}
\pyfunc{probability\_integrated} has no default \code{spectrum}: a flat
weight over an arbitrary energy grid is a modelling choice, not a safe
default, so every caller must pass one explicitly (an explicit flat tensor
if that is genuinely intended). \pyfunc{probability\_integrated\_angular}
takes nothing in its place, since its weight is purely geometric. Both
energy- and angle-specific functions are self-contained implementations
(their own trapezoidal-rule code), not thin wrappers around
\pyfunc{probability\_weighted\_average}; the latter exists as the shared
pattern for coordinates --- such as production height --- that do not
warrant a dedicated function of their own.
\end{implbox}

\section{\texttt{flux.py} --- from probabilities to fluxes}

\modulehead{core/common/flux.py}{multiplying final flavour probabilities by
flux normalisations and optional spectra.}

\begin{theorybox}
\subsection*{From probabilities to a differential flux}
If $\Phi_\alpha^{\rm prod}$ is the flux (or its normalised spectral shape)
of the produced flavour $\alpha$, the detected flux in flavour $\beta$ is
\begin{equation}
  \Phi_\beta^{\rm det}(E) = \Phi_\alpha^{\rm norm}\cdot f(E)\cdot P(\nu_\alpha\to\nu_\beta; E),
\end{equation}
where $f(E)$ is an optional normalised spectrum. This module knows nothing
about the physical origin of $\Phi^{\rm prod}$ or the medium traversed: it
is purely the final algebraic step, shared by \texttt{medium.vacuum.flux},
\texttt{medium.solar.flux}, \texttt{medium.earth.flux}, and
\texttt{medium.atmosphere.flux}. When the full transition matrix
$P_{\beta\alpha}$ is available rather than an already-selected final-flavour
probability, the equivalent matrix form
$\Phi_\beta^{\rm det}=\sum_\alpha P_{\beta\alpha}\,\Phi_\alpha^{\rm prod}$
is provided directly by \pyfunc{flux\_transition} (see below), mirroring
\pyfunc{probability\_incoherent}'s plain-matrix branch
(\texttt{probability.py} above).

\subsection*{From differential flux to event rate}
$\Phi_\beta^{\rm det}(E)$ above is a \emph{differential} flux: a density in
energy (and, when angularly resolved, in zenith/nadir angle), not yet a
rate. Obtaining a neutrino \textbf{rate} --- the quantity that, multiplied
by an effective area and an exposure time, gives an expected event count
--- requires integrating the energy dependence out:
\begin{equation}
  \keyresult{\Phi_\beta^{\rm rate} = \int \Phi_\beta^{\rm det}(E)\;dE},
\end{equation}
evaluated numerically by the trapezoidal rule over a finite energy grid.
When $\Phi_\beta^{\rm det}$ additionally carries an angular dependence
(e.g.\ the atmospheric or Earth-crossing zenith angle), integrating only
over energy keeps that angular resolution, giving a rate resolved by both
angle and flavour. A further integration over the full zenith range is
provided by this module too (\pyfunc{flux\_integrated\_angular} below), as
the geometric counterpart to energy integration; unlike
\texttt{probability.py} above, both reductions here are \textbf{unnormalised}
totals -- a flux is a physical rate density, not a probability, so nothing
should be divided out. Every function in this module, like every function in
\texttt{probability.py}, accepts a final flavour dimension of 3 or 4 (the
3+1 sterile extension).
\end{theorybox}

\begin{implbox}
\begin{itemize}
  \item \pyfunc{flux\_state(probability, flux, spectrum=None)}:
    multiplies and preserves the final-flavour dimension; if
    \code{spectrum} is given, it is applied as an element-wise
    multiplicative factor before reduction. Every medium's
    \code{[medium]\_flux\_state} (\cref{ch:vacuum,ch:solar,ch:earth,ch:atmosphere})
    delegates to this function.
  \item \pyfunc{flux\_transition(probability\_matrix, initial\_flux)}:
    the matrix form
    $\Phi_\beta=\sum_\alpha P_{\beta\alpha}\Phi_\alpha$ described above, for
    callers that hold the full transition matrix rather than a
    final-flavour probability vector.
\end{itemize}
Three functions implement the unnormalised-integral pattern above, one
generic and two coordinate-specific:
\begin{center}
\small
\begin{tabular}{@{}lll@{}}
\toprule
Function & Coordinate & Formula implemented \\
\midrule
\parbox[t]{5.2cm}{\pyfunc{flux\_integrated\_coordinate(\allowbreak flux, coordinate, dim=-2)}} &
  \parbox[t]{2.2cm}{arbitrary} &
  \parbox[t]{5.8cm}{$\displaystyle \Phi_{\rm out}=\int \Phi(x)\,dx$ --- the generic,
  unnormalised primitive} \\[10pt]
\parbox[t]{5.2cm}{\pyfunc{flux\_integrated(\allowbreak flux, E\_grid\_MeV, energy\_dim=-2)}} &
  \parbox[t]{2.2cm}{energy $E$} &
  \parbox[t]{5.8cm}{$\displaystyle \Phi^{\rm rate}=\int \Phi(E)\,dE$ --- delegates directly to
  \pyfunc{flux\_integrated\_coordinate}} \\[10pt]
\parbox[t]{5.2cm}{\pyfunc{flux\_integrated\_angular(\allowbreak flux, theta\_deg, angular\_dim=-2)}} &
  \parbox[t]{2.2cm}{zenith angle $\theta$} &
  \parbox[t]{5.8cm}{$\displaystyle \keyresult{\Phi_{\rm total} = 2\pi\int \frac{\mathrm{d}\Phi}{\mathrm{d}\Omega}\sin\theta\,\mathrm{d}\theta}$
  --- its own implementation (extra $2\pi$ factor), not a wrapper} \\
\bottomrule
\end{tabular}
\end{center}
\pyfunc{flux\_integrated}'s default \code{energy\_dim=-2} matches a plain
\code{(..., E, 3)} flux; an angle-resolved \code{(..., E, angle, 3)} flux is
integrated with \code{energy\_dim=-3}, preserving the angle axis in the
result. \pyfunc{flux\_integrated\_angular} assumes azimuthal symmetry and
is evaluated by the trapezoidal rule with an explicit $2\pi$ azimuthal
factor -- unlike \pyfunc{probability\_integrated\_angular}
(\texttt{probability.py} above), none of the three reductions here
normalise: a flux is a physical rate density, not a probability, so
nothing should be divided out.
\end{implbox}

\section{\texttt{neutrino.py} --- flavour index convention}

\modulehead{core/common/neutrino.py}{\pyclass{FLAVOUR\_TO\_INDEX} dictionary,
\pyfunc{flavour\_index} validator, and \pyfunc{flavour\_state} builder.}

\begin{implbox}
Fixes, once and for all, the flavour index convention used in every tensor
across the project: $e\to0,\ \mu\to1,\ \tau\to2,\ s\to3$. The sterile index
has no charged-lepton counterpart, but the 3+1 extension
(\pyclass{PMNS\_sterile}, \cref{ch:core-bsm}) fixes its position at index 3,
so it is included under a small set of aliases (\code{"s"}, \code{"sterile"},
\code{"nus"}, \code{"nu\_s"}) for callers that want to select or build it by
name instead of hand-indexing a length-4 tensor. \pyfunc{flavour\_index}
accepts these text aliases and validates out-of-range integer indices,
preventing silent indexing errors when building initial states;
\pyfunc{flavour\_state} returns the corresponding unit coherent-state
vector, sized to \code{n\_flavours} (3 by default, or 4 to select the
sterile flavour).
\end{implbox}

\chapter{\texttt{core.SM} --- Standard Model specialisation (3 flavours)}
\label{ch:core-sm}

\section{\texttt{pmns.py} --- the Standard Model mixing matrix}

\modulehead{core/SM/sm\_pmns.py}{\pyclass{PMNS\_SM}: the concrete 3-flavour
implementation of the abstract class \pyclass{core.common.pmns.PMNS}.}

\begin{theorybox}
\subsection*{The Standard Model product structure}
This is the only module that explicitly fixes a product structure for the
generic factors defined in \cref{ch:core-common}. The full 3-flavour PMNS
matrix is
\begin{equation}
  U = R_{23}(\theta_{23})\,\Delta(\delta)\,R_{13}(\theta_{13})\,\Delta(\delta)^\dagger\,R_{12}(\theta_{12}),
  \qquad \Delta(\delta) = \diag(1,1,e^{i\delta}),
\end{equation}
built from the generators $R_{12}$, $R_{13}$, $R_{23}$, $\Delta$ defined in
\cref{ch:core-common}. \pyfunc{PMNS\_SM.outer\_block} and
\pyfunc{PMNS\_SM.reduced} identify
\begin{equation}
  O \equiv R_{23}(\theta_{23})\,\Delta(\delta),
  \qquad
  U_{\rm red} \equiv R_{13}(\theta_{13})\,\Delta(\delta)^\dagger\,R_{12}(\theta_{12}),
  \qquad U = O\,U_{\rm red},
\end{equation}
so that this single factorisation is the concrete instance, for the SM,
of the abstract $U=O\,U_{\rm red}$ contract from \cref{ch:core-common}.

\subsection*{Why the outer block commutes with the matter Hamiltonian}
$O$ does not mix the electron-flavour index: $R_{23}$ acts only on the
$(\mu,\tau)$ block and $\Delta$ is diagonal, so $O$ has the block form
$O = \begin{psmallmatrix}1 & 0\\ 0 & O_{2\times2}\end{psmallmatrix}$ in the
flavour basis. Since the matter potential only has a non-zero $ee$ entry,
$H_{\rm mat}=\diag(V,0,0)$, it follows that
\begin{equation}
  O\,H_{\rm mat}\,O^\dagger = H_{\rm mat} \qquad \text{(they commute exactly).}
\end{equation}
This is precisely the property \cref{ch:core-common} assumes, without
proof, to derive the general factorisation \cref{eq:H-factorization}; here
it is verified explicitly for the concrete SM $O$.

\subsection*{Transforming the Hamiltonian to the reduced basis}
Given the flavour-basis Hamiltonian
$H_{\rm flavour}=H_{\rm kin}+H_{\rm mat}=U\diag(k_1,k_2,k_3)\,U^\dagger + H_{\rm mat}$
(\texttt{hamiltonian.py} section, \cref{ch:core-common}), substituting
$U=O\,U_{\rm red}$ and using $O\,H_{\rm mat}\,O^\dagger=H_{\rm mat}$ gives
\begin{equation}
  H_{\rm flavour} = O\underbrace{\Bigl[U_{\rm red}\diag(k)\,U_{\rm red}^\dagger + H_{\rm mat}\Bigr]}_{\displaystyle H_{\rm red}}O^\dagger = O\,H_{\rm red}\,O^\dagger,
\end{equation}
matching \cref{eq:H-factorization}: \pyfunc{PMNS\_SM.reduced\_basis}
applied to $H_{\rm flavour}$ (inherited from \pyclass{PMNS}, computing
$O^\dagger(\cdot)O$) yields exactly $H_{\rm red}$ above.

\subsection*{The CP phase drops out of the reduced kinetic term}
$R_{12}$ does not mix the $\tau$ index (third index) and $\Delta$ acts
non-trivially only on that same index, so the product
$R_{12}\diag(k)R_{12}^T$ has no non-zero entries connecting to index 3;
conjugation by $\Delta$ therefore acts trivially on it:
\begin{equation}
  \Delta^\dagger\bigl(R_{12}\diag(k)R_{12}^T\bigr)\Delta = R_{12}\diag(k)R_{12}^T,
\end{equation}
so that \textbf{the CP phase $\delta$ drops out completely from the reduced
kinetic term}, which reduces to the real matrix
$U_{\rm red}^{(0)} \equiv R_{13}R_{12}$:
\begin{equation}
  \keyresult{U_{\rm red}\diag(k)\,U_{\rm red}^\dagger = \bigl(R_{13}R_{12}\bigr)\,\diag(k)\,\bigl(R_{13}R_{12}\bigr)^T.}
\end{equation}
This is why \pyfunc{PMNS\_SM.reduced} returns $R_{13}R_{12}$ rather than
the formally-defined $R_{13}\Delta^\dagger R_{12}$: the two coincide once
contracted with the (real, diagonal) kinetic term, but $R_{13}R_{12}$ is
real, cheaper to evolve, and avoids carrying a phase that always cancels.

\subsection*{Recovering the flavour basis}
Once $H_{\rm red}$ (and its evolutor $S_{\rm red}$) has been obtained, the
flavour-basis result is recovered by dressing once with $O$
(\pyfunc{PMNS\_SM.flavour\_basis}, inherited from \pyclass{PMNS}, computing
$O(\cdot)O^\dagger$):
\begin{equation}
  H_{\rm flavour} = O\,H_{\rm red}\,O^\dagger,
  \qquad
  S_{\rm flavour}(x) = O\,S_{\rm red}(x)\,O^\dagger.
\end{equation}
\end{theorybox}

\begin{keyeqbox}
\textbf{Coherent propagation scheme.} In vacuum ($H_{\rm mat}=0$), a
coherent flavour state is most simply understood, and computed, entirely
in the mass basis:
\begin{enumerate}
  \item \textbf{Flavour $\to$ mass}: project onto mass eigenstates,
    $\psi_{\rm mass}(x_0) = U^\dagger\,\psi_{\rm flavour}(x_0)$;
  \item \textbf{Coherent propagation}: each mass eigenstate accumulates its
    own phase independently, since the mass basis diagonalises the kinetic
    term by construction,
    $\psi_{\rm mass}(x) = \diag\bigl(e^{-ik_1x},\dots,e^{-ik_3x}\bigr)\,\psi_{\rm mass}(x_0)$;
  \item \textbf{Mass $\to$ flavour}: project back,
    $\psi_{\rm flavour}(x) = U\,\psi_{\rm mass}(x)$.
\end{enumerate}
Composing the three steps gives the familiar vacuum evolutor
$S_{\rm vacuum}(x)=U\diag\bigl(e^{-ik_ix}\bigr)U^\dagger$
(\texttt{medium.vacuum.evolutor}). In matter, $H_{\rm mat}$ does not
commute with the full $U$, so this three-step picture does not apply
directly at the full-$U$ level; the reduction trick derived above is
exactly what restores an analogous, computationally cheap three-step
scheme -- flavour $\to$ reduced $\to$ propagate $\to$ flavour -- using $O$
and $U_{\rm red}$ in place of $U$.
\end{keyeqbox}

\begin{implbox}
\pyclass{PMNS\_SM(params: PMNSParams)} implements only the three abstract
methods that depend on the product structure: \pyfunc{pmns\_matrix} (the
full $3\times3$ $U$), \pyfunc{reduced} ($U_{\rm red}=R_{13}R_{12}$), and
\pyfunc{outer\_block} ($O=R_{23}\Delta$). Everything else (rotation
builders, \pyfunc{flavour\_basis}/\pyfunc{reduced\_basis}, antineutrino
selection via \pyfunc{select\_antinu}) is inherited unchanged from the base
class, without duplicating code. It is the default scenario for the whole
project: any function in \texttt{medium.*} that receives a generic
\pyclass{PMNS} object works unmodified with either \pyclass{PMNS\_SM} or
\pyclass{PMNS\_sterile} (\cref{ch:core-bsm}), thanks to the inheritance
design of \cref{ch:core-common}.
\end{implbox}

\begin{refbox}
Standard PMNS parametrisation: \cite{ParticleDataGroup}. Reference numerical
values (global best fits) used in the default presets:
\cite{Esteban2020NuFit,NuFit52}.
\end{refbox}

\section{\texttt{mass\_spectrum.py} --- Standard Model mass spectrum}

\modulehead{core/SM/sm\_mass\_spectrum.py}{\pyclass{MassSpectrum\_SM}, the
concrete 3-flavour implementation of \pyclass{MassSpectrum}
(\cref{ch:core-common}).}

\begin{theorybox}
The Standard Model scenario has no sterile state to accommodate, so
\pyfunc{difference\_vector()} simply returns
\pyfunc{difference\_vector\_base()} unchanged, sized to exactly 3
components:
\begin{equation}
  \keyresult{\bigl(\Delta m^2_1,\ \Delta m^2_2,\ \Delta m^2_3\bigr)
  =\begin{cases}
      \bigl(0,\ \Delta m^2_{21},\ \Delta m^2_{3\ell}\bigr) & \Delta m^2_{3\ell} > 0\ \text{(NO)},\\[2pt]
      \bigl(-\Delta m^2_{21},\ 0,\ \Delta m^2_{3\ell}\bigr) & \Delta m^2_{3\ell} < 0\ \text{(IO)}.
    \end{cases}}
\end{equation}
The ordering-dependent construction itself -- selecting between the two
branches above from the sign of $\Delta m^2_{3\ell}$ -- is entirely
inherited from \pyclass{MassSpectrum} (\cref{ch:core-common},
\texttt{mass\_spectrum.py} section): \pyclass{MassSpectrum\_SM} adds no
logic of its own beyond fixing the output size to 3, in contrast with
\pyclass{MassSpectrum\_BSM} (\cref{ch:core-bsm}), which appends a fourth,
sterile entry $\Delta m^2_{41}$ to the same base vector.
\end{theorybox}

\begin{implbox}
\pyclass{MassSpectrum\_SM(DeltamSq21, DeltamSq3l)}: adds no fields beyond
the base \pyclass{MassSpectrum}; \pyfunc{difference\_vector()} is a
single-line override, \code{return
self.difference\_vector\_base(context=context)}.
\end{implbox}

\chapter{\texttt{core.BSM} --- Beyond-the-Standard-Model extensions}
\label{ch:core-bsm}

This chapter covers the two supported BSM extensions: non-standard
interactions (NSI) in the 3-flavour sector, and a fourth light sterile
neutrino (the ``3+1'' scheme). Both reuse as much as possible of the
infrastructure from \cref{ch:core-common}: NSI only modifies $H_{\rm mat}$,
reusing whichever \pyclass{PMNS} object is otherwise in use, and the
sterile sector extends the factorisation $U=O\,U_{\rm red}$ to
$4\times4$. There is no separate \texttt{core.BSM} Hamiltonian-assembly
module: \pyfunc{hamiltonian\_reduced} and \pyfunc{hamiltonian\_flavour}
(\cref{ch:core-common}, \texttt{hamiltonian.py} section) already handle NSI,
the sterile extension, or both simultaneously.

\section{\texttt{bsm\_nsi.py} --- non-standard interactions}

\modulehead{core/BSM/bsm\_nsi.py}{parametric configuration of NSI couplings
and construction of the $\epsilon$ matrix.}

\begin{theorybox}
NSI are effective four-fermion operators that modify coherent forward
scattering beyond the standard MSW potential \parencite{Grossman1995}. The
matter Hamiltonian in the flavour basis generalises to
\begin{equation}
  \keyresult{H_{\rm mat}^{\rm NSI} = V\left(\begin{pmatrix}1&0&0\\0&0&0\\0&0&0\end{pmatrix} + \epsilon\right)},
  \qquad V = \pm\sqrt2\,G_F\,n_e\,L_{\rm scale},
\end{equation}
with $\epsilon$ a dimensionless $3\times3$ Hermitian matrix,
\begin{equation}
  \epsilon = \begin{pmatrix}
    \epsilon_{ee} & \epsilon_{e\mu} & \epsilon_{e\tau} \\
    \epsilon_{e\mu}^* & \epsilon_{\mu\mu} & \epsilon_{\mu\tau} \\
    \epsilon_{e\tau}^* & \epsilon_{\mu\tau}^* & \epsilon_{\tau\tau}
  \end{pmatrix}.
\end{equation}
Diagonal entries are real; off-diagonal entries are generally complex. The
Standard Model limit is $\epsilon=0$. Only the traceless part of $\epsilon$
matters for oscillations (an overall trace is an irrelevant global phase),
so fixing $\epsilon_{\mu\mu}=0$ is a valid convention for the diagonal
sector without loss of physical generality.

\textbf{No NSI-specific mixing matrix.} \pyclass{NSIConfig} touches only
$H_{\rm mat}$: it does not define, or require, any dedicated \pyclass{PMNS}
subclass. If the sterile extension is not also active, NSI propagation
reuses the plain 3-flavour \pyclass{PMNS\_SM} object from \cref{ch:core-sm}
completely unchanged -- the same $U$, $O$, $U_{\rm red}$ as the pure
Standard Model case. \pyclass{OscillationParameters.nsi} is simply an
optional, orthogonal attribute alongside \code{pmns}
(\cref{ch:core-common}, \texttt{oscillation.py} section): which mixing
scenario is active and whether NSI is active are two independent choices,
and every combination (SM+NSI, sterile alone, sterile+NSI) is valid.
\end{theorybox}

\begin{theorybox}
\subsection*{Transforming the matter Hamiltonian: Standard Model versus NSI}
In the plain Standard Model, the reduced-basis matter term needs no
rotation at all: $H_{\rm mat}^{\rm red}=H_{\rm mat}^{\rm flavour}=\diag(V,0,\dots,0)$,
because $O\,H_{\rm mat}\,O^\dagger=H_{\rm mat}$ exactly (\cref{ch:core-sm}) --
the matter term is identical in the flavour and reduced bases, so it is
simply added to the reduced kinetic term as-is.

When NSI are active this shortcut disappears: as shown below,
$O\,H_{\rm mat}^{\rm NSI}\,O^\dagger\neq H_{\rm mat}^{\rm NSI}$ in general.
\pyfunc{hamiltonian\_matter\_reduced} (\cref{ch:core-common}) therefore
takes the opposite route: it builds $H_{\rm mat}^{\rm NSI}$ directly in the
\emph{flavour} basis as above, and then rotates it explicitly \emph{down}
into the reduced basis,
\begin{equation}
  \keyresult{H_{\rm mat}^{\rm NSI,\,red} = O^\dagger\,H_{\rm mat}^{\rm NSI,\,flavour}\,O},
\end{equation}
the exact inverse of the flavour-basis dressing $O(\cdot)O^\dagger$
(\cref{ch:core-common}). Because the transformation is performed
explicitly here -- rather than assumed to be a no-op, as in the SM case
above -- the identity
$H_{\rm flavour}=\text{\pyfunc{pmns.flavour\_basis}}(H_{\rm red})$ holds
exactly for \emph{arbitrary} $\epsilon$, not only for the
diagonal-uniform Standard Model limit.
\end{theorybox}

\begin{notebox}
\textbf{Adjoint, not transpose.} $O$ is generally complex ($\Delta$ carries
$e^{i\delta}$), so $O^\dagger=(O^{*})^{T}$ and the plain transpose $O^{T}$
differ whenever $\delta\neq0$. The rotation above \textbf{must} use the
conjugate transpose (\code{O.conj().transpose(-2,-1)} in
\pyfunc{hamiltonian\_matter\_reduced}) to preserve Hermiticity: for
Hermitian $\epsilon$ (and hence Hermitian $H_{\rm mat}^{\rm NSI}$),
$O^\dagger\,H\,O$ is again Hermitian for any unitary $O$, whereas $O^{T}\,H\,O$
is not, in general, even Hermitian -- let alone the physically correct
transformed operator. This mirrors the identical requirement already
established for the kinetic term (\cref{ch:core-common},
\texttt{pmns.py}/\texttt{hamiltonian.py} sections): the Hermitian conjugate
is the only convention consistent with a unitary change of basis for an
operator; the plain transpose coincides with it only by accident, when the
matrix involved happens to be real.
\end{notebox}

\begin{notebox}
Since $H_{\rm mat}^{\rm NSI}$ is no longer proportional to $\diag(V,0,0)$,
the commutation property $O\,H_{\rm mat}\,O^\dagger=H_{\rm mat}$ assumed by
the general factorisation \cref{eq:H-factorization} \textbf{no longer holds
in general} once $\epsilon$ has entries outside the $ee$ block. The code
does not rely on it being true -- as shown above,
\pyfunc{hamiltonian\_matter\_reduced} builds $H_{\rm mat}^{\rm NSI}$ in the
flavour basis and rotates it explicitly, so the result is exact for any
$\epsilon$. What \emph{is} lost is the Standard Model \emph{simplification}
(\cref{ch:core-sm}) that $H_{\rm red}$ needs no reference to
$\theta_{23}/\delta$ and can be built with a real $U_{\rm red}$: for
generic NSI, $H_{\rm red}$ is genuinely complex and depends on every
mixing parameter, so the closed-form spectral formulas used elsewhere in
the project no longer apply, and the relevant diagonalisation (e.g.\ in
\texttt{medium.solar.probability}) is performed numerically instead, via
\code{torch.linalg.eigh}.
\end{notebox}

\begin{implbox}
\begin{itemize}
  \item \pyclass{NSIConfig} (frozen dataclass): stores the six independent
    parameters of $\epsilon$ (3 real diagonal + 3 complex off-diagonal) and
    exposes \pyfunc{epsilon\_tensor\_base()} (the raw complex $3\times3$
    tensor), \pyfunc{epsilon\_tensor(n\_flavours=...)} (embedded for any
    flavour count, zero-padded for the sterile rows/columns), and
    \pyfunc{is\_sm\_limit}.
  \item \pyfunc{NSIConfig.from\_preset(name)} builds named configurations
    from \code{tpeanuts.config.presets.NSI\_PRESETS}
    (\cref{sec:nsi-presets}), each with its physical justification
    documented.
  \item The $\epsilon$ tensor is read via
    \pyfunc{oscillation.nsi.epsilon\_tensor(...)}, passed through
    \pyfunc{pmns.select\_antinu} ($\epsilon\to\epsilon^{*}$ for
    antineutrinos, mirroring the $U\to U^{*}$ convention), and injected
    directly into \pyfunc{hamiltonian\_matter\_reduced}
    (\cref{ch:core-common}) -- the single Hamiltonian-assembly function
    handling this case; there is no separate \texttt{core.BSM} builder.
\end{itemize}
\end{implbox}

\begin{refbox}
Original NSI proposal: \cite{Grossman1995}. Model-independent bounds:
\cite{BiggioBlennowFM2009}. LMA-Dark degeneracy in the global fit:
\cite{Esteban2018LMADark}. IceCube DeepCore bounds: \cite{IceCube2022NSI}.
\end{refbox}

\section{NSI presets}
\label{sec:nsi-presets}

\modulehead{tpeanuts/config/presets.py}{\code{NSI\_PRESETS}: named,
literature-justified $\epsilon$ configurations, built via
\pyfunc{NSIConfig.from\_preset}.}

\begin{implbox}
Every preset below sets a handful of $\epsilon$ parameters to a
representative, literature-motivated value (all other components default
to zero) rather than to a unique global best fit: current data leaves large
NSI degeneracies, so no single preferred point exists. Their purpose is
threefold: (1) give the validation notebooks and tests physically motivated
$\epsilon$ values instead of arbitrary numbers, covering diagonal-only,
off-diagonal, single- and multi-parameter regimes; (2) reproduce, and stay
traceable to, specific published sensitivity/bound studies; and
(3) let users reach a standard benchmark scenario (e.g.\ the LMA-Dark
degeneracy) by name instead of re-deriving its $\epsilon$ values.

\begin{center}
\small
\begin{tabular}{@{}ll@{}}
\toprule
Preset & $\epsilon$ values and origin \\
\midrule
\parbox[t]{3.2cm}{\code{sm\_\allowbreak{}no\_\allowbreak{}nsi}} &
  \parbox[t]{9.4cm}{All $\epsilon=0$. Standard Model limit; equivalent to
  passing \code{nsi=None}.} \\[10pt]
\parbox[t]{3.2cm}{\code{nsi\_\allowbreak{}diagonal\_\allowbreak{}biggio2009}} &
  \parbox[t]{9.4cm}{$\epsilon_{ee}=0.30$, $\epsilon_{\tau\tau}=0.15$.
  Diagonal-only benchmark within the model-independent 90\% CL bounds of
  \textcite{BiggioBlennowFM2009} ($|\epsilon_{ee}|<0.50$,
  $|\epsilon_{\tau\tau}|<0.33$).} \\[18pt]
\parbox[t]{3.2cm}{\code{nsi\_\allowbreak{}lma\_\allowbreak{}dark\_\allowbreak{}esteban2018}} &
  \parbox[t]{9.4cm}{$\epsilon_{ee}=-2.0$. The LMA-Dark degenerate solar
  solution of \textcite{Esteban2018LMADark}: this large negative
  $\epsilon_{ee}$ flips the sign of the effective MSW potential, mimicking
  $\theta_{12}\to\pi/2-\theta_{12}\approx56.6^\circ$ (paired with the
  \code{"\_LMA\_DARK\_NUFIT52\_NO"} oscillation preset,
  \cref{ch:core-common}).} \\[24pt]
\parbox[t]{3.2cm}{\code{nsi\_\allowbreak{}offdiag\_\allowbreak{}icecube2022}} &
  \parbox[t]{9.4cm}{$\epsilon_{\mu\tau}=0.005$, $\epsilon_{\tau\tau}=0.015$.
  Sits near the 90\% CL boundary of the IceCube DeepCore atmospheric
  analysis \parencite{IceCube2022NSI}
  ($|\epsilon_{\mu\tau}|<0.0060$, $|\epsilon_{\tau\tau}-\epsilon_{\mu\mu}|<0.018$).} \\[18pt]
\parbox[t]{3.2cm}{\code{nsi\_\allowbreak{}dune\_\allowbreak{}etau}} &
  \parbox[t]{9.4cm}{$\epsilon_{e\tau}=0.15$ (real). Representative DUNE
  Near Detector sensitivity benchmark \parencite{DUNE2021BSM}: this
  flavour-changing coupling produces the strongest degeneracies with
  $\delta_{CP}$, the mass ordering, and $\theta_{23}$.} \\[18pt]
\parbox[t]{3.2cm}{\code{nsi\_\allowbreak{}hyperk\_\allowbreak{}etau}} &
  \parbox[t]{9.4cm}{$\epsilon_{ee}=0.20$, $\epsilon_{e\tau}=0.10$ (real).
  Combines a diagonal MSW correction with a flavour-changing coupling,
  representative of Hyper-Kamiokande sensitivity studies
  \parencite{HyperK2015Proposal}.} \\[18pt]
\parbox[t]{3.2cm}{\code{nsi\_\allowbreak{}globalfit\_\allowbreak{}esteban2018}} &
  \parbox[t]{9.4cm}{$\epsilon_{ee}=0.25$, $\epsilon_{e\tau}=0.12$,
  $\epsilon_{\mu\tau}=0.01$, $\epsilon_{\tau\tau}=0.08$. Multi-parameter
  benchmark within the allowed region of the global oscillation analysis of
  \textcite{Esteban2018LMADark}; not a unique best fit, since significant
  correlations remain.} \\[22pt]
\parbox[t]{3.2cm}{\code{nsi\_\allowbreak{}ee\_\allowbreak{}only}} &
  \parbox[t]{9.4cm}{$\epsilon_{ee}=0.30$. Simplest possible NSI scenario:
  only the effective MSW potential is modified, with no flavour-changing
  coupling \parencite{BiggioBlennowFM2009}.} \\[10pt]
\parbox[t]{3.2cm}{\code{nsi\_\allowbreak{}mutau\_\allowbreak{}only}} &
  \parbox[t]{9.4cm}{$\epsilon_{\mu\tau}=0.01$ (real). Atmospheric
  $\mu$--$\tau$ benchmark representative of IceCube/KM3NeT sensitivity to
  flavour-changing matter interactions \parencite{Coloma2017IceCubeNSI}.} \\[14pt]
\bottomrule
\end{tabular}
\end{center}
\end{implbox}

\begin{refbox}
\cite{BiggioBlennowFM2009,Esteban2018LMADark,IceCube2022NSI,DUNE2021BSM,HyperK2015Proposal,Coloma2017IceCubeNSI}.
\end{refbox}

\section{\texttt{bsm\_sterile.py} --- the 3+1 scheme}

\modulehead{core/BSM/bsm\_sterile.py}{\pyclass{PMNS\_sterile}: $4\times4$
mixing for one additional light sterile neutrino.}

\begin{theorybox}
A fourth mass eigenstate $\nu_4$ and a fourth sterile flavour eigenstate
$\nu_s$ are added; $\nu_s$ is a gauge singlet and therefore couples to
neither $W$ nor $Z$. The charged-current potential is consequently
unaffected by the extension, so, omitting the optional neutral-current term
(the default, \cref{ch:core-common}), the matter potential in the extended
basis $(e,\mu,\tau,s)$ is
\begin{equation}
  H_{\rm mat}=\diag(V_{CC},0,0,0),
\end{equation}
the same charged-current-only term as the 3-flavour case, simply padded
with a zero sterile entry. Since $\nu_s$ also does not receive the
neutral-current potential $V_{NC}$, optionally supplying a neutron density
activates the relative term derived in \cref{ch:core-common},
\begin{equation}
  H_{\rm mat}=\diag(V_{CC},0,0,-V_{NC}),
\end{equation}
which does affect sterile-active oscillation probabilities (see the
implementation box below for how this interacts with $O_4$). The $4\times4$
mixing matrix is parametrised by adding three active--sterile rotations,
$R_{14},R_{24},R_{34}$, to the standard structure, following the
parametrisation convention of \textcite{KoppMachadoMaltoniSchwetz2013}:
\begin{equation}
  U_4 = R_{23}\,\Delta_4\,R_{24}\,R_{34}\,R_{13}\,\Delta_4^\dagger\,R_{12}\,R_{14}.
\end{equation}
Identifying $O_4\equiv R_{23}\,\Delta_4\,R_{24}\,R_{34}$
(\pyfunc{PMNS\_sterile.outer\_block}) leaves a formal reduced factor
$U_{\rm red,4}\equiv R_{13}\,\Delta_4^\dagger\,R_{12}\,R_{14}$ satisfying
$U_4=O_4\,U_{\rm red,4}$ exactly, mirroring \cref{ch:core-sm}. As in the SM
case, $\Delta_4$ cancels from the reduced \emph{kinetic term} specifically:
\begin{equation}
  \keyresult{\Delta_4^\dagger\bigl(R_{12}R_{14}\diag(k)R_{14}^\dagger R_{12}^\dagger\bigr)\Delta_4
  = R_{12}R_{14}\diag(k)R_{14}^\dagger R_{12}^\dagger},
\end{equation}
valid because neither $R_{12}$ (the $e$--$\mu$ plane) nor $R_{14}$ (the
$e$--$s$ plane) mixes with the $\tau$ index, which is the only non-trivial
index of $\Delta_4$. This is why \pyfunc{PMNS\_sterile.reduced} returns the
simplified, $\Delta_4$-free
\begin{equation}
  U_{\rm red,4}^{\rm code}=R_{13}R_{12}R_{14}
\end{equation}
rather than the formal $U_{\rm red,4}$ above: the two agree once contracted
with the kinetic term (any diagonal unitary conjugating a diagonal matrix
leaves it invariant), but $U_{\rm red,4}^{\rm code}$ is cheaper. Note that,
consequently, $O_4\,U_{\rm red,4}^{\rm code}\neq U_4$ as bare matrices --
only $U_{\rm red,4}$ (with $\Delta_4^\dagger$ included) satisfies the exact
factorisation; \pyfunc{PMNS\_sterile.reduced} returns the code-optimised
form, valid specifically for building $H_{\rm kin}$, not a literal factor
of $U_4$ (see the limitations box below). The same commutation argument as
\cref{ch:core-sm} applies to $O_4$ for the plain charged-current matter
term: $R_{23},\Delta_4,R_{24},R_{34}$ all leave the $e$ index fixed, so
$O_4\,\diag(V_{CC},0,0,0)\,O_4^\dagger=\diag(V_{CC},0,0,0)$ exactly, and the
reduced Hamiltonian is assembled and dressed back to the flavour basis
exactly as in \cref{ch:core-common,ch:core-sm}, now for $n=4$. As
established in \cref{ch:core-common} (\texttt{hamiltonian.py} section),
this commutation is specific to the charged-current-only term: $O_4$ mixes
the $\{\mu,\tau,s\}$ block and therefore does \emph{not} leave the optional
neutral-current term $\diag(0,0,0,V_{NC})$ invariant, so the reduction-trick
shortcut used above breaks whenever that term is active (see the
implementation box below).

\textbf{SM limit.} When $\theta_{14}=\theta_{24}=\theta_{34}=0$, the
$4\times4$ matrices block-diagonalise exactly into the $3\times3$ SM matrix
\emph{embedded} in the larger space, recovering every result of
\cref{ch:core-common,ch:core-sm} without approximation.

\textbf{CP-phase counting.} For $N=4$ Dirac generations there are
$(N-1)(N-2)/2=3$ physical phases: $\delta_{13}$ (the standard phase, in the
$\Delta_4$ carried through $R_{13}$), $\delta_{14}$ ($e$--$s$ sector, in
$R_{14}$), and $\delta_{24}$ ($\mu$--$s$ sector, in $R_{24}$). $R_{34}$ is
always real: the fourth possible phase can be absorbed into a field
redefinition, so sterile presets do not include \code{delta34\_deg}.
\end{theorybox}

\begin{theorybox}
\subsection*{Transforming to the reduced basis: $H_{\rm kin}$ and $H_{\rm mat}$}
Exactly as in \cref{ch:core-sm}, the reduced-basis kinetic term uses the
code's $U_{\rm red,4}^{\rm code}=R_{13}R_{12}R_{14}$ together with its
\emph{Hermitian conjugate} (never the plain transpose, for the same reason
as \cref{ch:core-common,ch:core-sm}):
\begin{equation}
  \keyresult{H_{\rm kin,4}^{\rm red} = U_{\rm red,4}^{\rm code}\,\diag(k_1,k_2,k_3,k_4)\,\bigl(U_{\rm red,4}^{\rm code}\bigr)^\dagger}.
\end{equation}
Unlike the SM's $R_{13}R_{12}$, $U_{\rm red,4}^{\rm code}$ is genuinely
complex whenever the active-sterile CP phase $\delta_{14}\neq0$ (it carries
$R_{14}$, which -- unlike $R_{12}$/$R_{13}$ -- is built with a non-trivial
phase); the Hermitian-conjugate convention is therefore not merely a safe
default here, but the only correct choice.

\textbf{Without NSI and without the neutral-current term} (the default),
the matter term needs no basis change at all:
$O_4\,\diag(V_{CC},0,0,0)\,O_4^\dagger=\diag(V_{CC},0,0,0)$ exactly (shown
above), so
\begin{equation}
  H_{\rm mat,4}^{\rm red} = H_{\rm mat,4}^{\rm flavour} = \diag(V_{CC},0,0,0),
\end{equation}
and the reduced Hamiltonian is simply
$H_{\rm red,4}=H_{\rm kin,4}^{\rm red}+H_{\rm mat,4}^{\rm red}$
(\pyfunc{hamiltonian\_reduced}, \cref{ch:core-common}), dressed back to the
flavour basis with $O_4(\cdot)O_4^\dagger$ (\pyfunc{hamiltonian\_flavour}).
This case introduces no approximation beyond \cref{ch:core-sm}: the
reduced-basis simplification stays exact for any $\theta_{14},\theta_{24},\theta_{34}$.
Supplying a neutron density activates
$H_{\rm mat,4}^{\rm flavour}=\diag(V_{CC},0,0,-V_{NC})$
(\cref{ch:core-common}) instead: since this no longer commutes with $O_4$,
the shortcut above is unavailable and \pyfunc{hamiltonian\_matter\_reduced}
takes the general rotated path,
$H_{\rm mat,4}^{\rm red}=O_4^\dagger\,H_{\rm mat,4}^{\rm flavour}\,O_4$,
even when NSI is off.

\textbf{Simultaneous NSI} is fully supported by the same code path
(\cref{sec:nsi-presets}): \pyfunc{NSIConfig.epsilon\_tensor(n\_flavours=4)}
embeds the $3\times3$ $\epsilon$ block into the top-left corner of a
$4\times4$ zero matrix (the sterile row/column stay zero, since $\nu_s$ is
an $\mathrm{SU}(2)_L\times\mathrm{U}(1)_Y$ singlet and carries no NSI
coupling), and \pyfunc{hamiltonian\_matter\_reduced} rotates the resulting
$H_{\rm mat,4}^{\rm NSI,\,flavour}=V\bigl(\diag(1,0,0,0)+\epsilon_4\bigr)$
into the reduced basis with the \emph{same} $O_4^\dagger(\cdot)O_4$
transformation as the 3-flavour case:
\begin{equation}
  H_{\rm mat,4}^{\rm NSI,\,red} = O_4^\dagger\,H_{\rm mat,4}^{\rm NSI,\,flavour}\,O_4,
\end{equation}
again exact for arbitrary $\epsilon$, since it is computed explicitly
rather than assumed to commute with $O_4$ -- the same approximation
trade-off as the 3-flavour NSI case applies here too (previous sections):
the closed-form real-arithmetic shortcut is lost, and diagonalisation falls
back to \code{torch.linalg.eigh}. Nothing in
\pyfunc{hamiltonian\_reduced}/\pyfunc{hamiltonian\_flavour} branches
explicitly on ``sterile or NSI or both'': the flavour count $n\in\{3,4\}$
and the presence of \code{oscillation.nsi} are read directly off
\pyclass{OscillationParameters}, so all four combinations (SM, SM+NSI,
sterile, sterile+NSI) are handled by the exact same two functions.
\end{theorybox}

\begin{notebox}
\textbf{Limitation: $U=O\,U_{\rm red}$ holds exactly only for the
\emph{formal} $U_{\rm red}$, not for \pyfunc{reduced}.} Both here and in
\cref{ch:core-sm}, the abstract contract of \cref{ch:core-common} --
$U=O\,U_{\rm red}$ for the three methods \pyfunc{pmns\_matrix},
\pyfunc{outer\_block}, \pyfunc{reduced} -- is, strictly, a simplification
of what the code returns: \pyfunc{reduced} always drops a $\Delta$ (or
$\Delta_4$) factor that is provably irrelevant for building $H_{\rm kin}$
(any diagonal unitary conjugating $\diag(k)$ leaves it invariant) but whose
omission means \pyfunc{outer\_block}$(\,)\cdot$\pyfunc{reduced}$(\,)$ is
\emph{not} bit-for-bit equal to \pyfunc{pmns\_matrix}$(\,)$. This is
harmless for every use of \pyfunc{reduced} in this project (always as the
sandwiching matrix of a kinetic term), but it means \cref{ch:core-common}'s
statement of the contract should be read as ``$U$ and $O\,U_{\rm red}$
agree up to a diagonal phase that leaves $\diag(k)$ invariant'', not as a
literal matrix identity between the three named methods.
\end{notebox}

\begin{implbox}
\begin{itemize}
  \item \pyclass{PMNSSterileParams}: parameters for \emph{only} the sterile
    extension ($\theta_{14},\theta_{24},\theta_{34}$ and their phases); the
    SM sector lives in a separate \pyclass{PMNSParams} object passed
    alongside this one to the \pyclass{PMNS\_sterile} constructor.
    $\Delta m^2_{41}$ also does not live here: like
    $\Delta m^2_{21}/\Delta m^2_{3\ell}$, it is stored on
    \pyclass{OscillationParameters} because it never enters a mixing
    rotation.
  \item \pyclass{PMNS\_sterile}: subclass of \pyclass{core.common.pmns.PMNS}
    with \code{n\_flavours=4}; it inherits unchanged every piece that is
    generic in size (\pyfunc{R12}, \pyfunc{R13}, \pyfunc{R23},
    \pyfunc{Delta}, \pyfunc{flavour\_basis}, \pyfunc{reduced\_basis},
    \pyfunc{refresh}, \pyfunc{select\_antinu}, ...) and only adds the extra
    constructors \pyfunc{R14}, \pyfunc{R24}, \pyfunc{R34} and the concrete
    implementation of \pyfunc{pmns\_matrix}, \pyfunc{reduced}, and
    \pyfunc{outer\_block}.
  \item Direct consequence of this design: \textbf{every} piece of
    \texttt{medium.*} code that operates on a generic \pyclass{PMNS} object
    works without modification for 4 flavours, simply by passing a
    \pyclass{PMNS\_sterile} instance instead of \pyclass{PMNS\_SM}; the same
    holds for NSI (previous sections), which can be combined freely with
    the sterile extension since \pyfunc{NSIConfig.epsilon\_tensor} embeds
    the $3\times3$ $\epsilon$ block in the top-left corner of the
    $4\times4$ matter matrix, leaving the sterile row/column at zero.
  \item Named presets (e.g.\ \code{"sterile\_3p1\_bestfit\_giunti2017"}) are
    built via \pyfunc{PropagationConfig.oscillation\_parameters\_from\_preset}
    (\cref{ch:core-common}), which internally creates both parameter
    objects and sets $\Delta m^2_{41}$ on the resulting
    \pyclass{OscillationParameters}.
\end{itemize}
\end{implbox}

\begin{refbox}
3+1 mixing parametrisation convention: \cite{KoppMachadoMaltoniSchwetz2013}.
Global 3+1 fit: \cite{GiuntiMarronePalazzo2017}. IceCube atmospheric
$\nu_\mu$ disappearance search: \cite{IceCube2020SterileNumuDisappearance}.
\end{refbox}

\section{Sterile presets}
\label{sec:sterile-presets}

\modulehead{tpeanuts/config/presets.py}{\code{OSCILLATION\_PRESETS} entries
carrying \code{theta14\_deg}: named, literature-justified 3+1 parameter
sets, dispatched to \pyclass{PMNS\_sterile} automatically.}

\begin{implbox}
A preset is a \emph{sterile} preset purely by having a \code{theta14\_deg}
key; \pyfunc{PropagationConfig.oscillation\_parameters\_from\_preset}
(\cref{ch:core-common}) dispatches on that key alone to build a
\pyclass{PMNS\_sterile} instead of a \pyclass{PMNS\_SM}, with every other
field (the SM angles, $\Delta m^2_{21}/\Delta m^2_{3\ell}$) inherited from
the shared NuFIT 5.2 baseline (\cref{ch:core-common}, Default Configuration
box). As with the NSI presets (\cref{sec:nsi-presets}), these are named,
traceable benchmark points rather than a single global best fit -- current
3+1 data is in tension between different anomalies, so no universally
preferred point exists.

\begin{center}
\small
\begin{tabular}{@{}ll@{}}
\toprule
Preset & Parameters and origin \\
\midrule
\parbox[t]{3.2cm}{\code{sterile\_\allowbreak{}3p1\_\allowbreak{}null\_\allowbreak{}mixing}} &
  \parbox[t]{9.4cm}{$\theta_{14}=\theta_{24}=\theta_{34}=0$,
  $\Delta m^2_{41}=1.0\ \mathrm{eV^2}$. Null-hypothesis reference: with
  every sterile angle at zero, \pyclass{PMNS\_sterile} is numerically
  identical to the plain 3-flavour SM (the exact SM-limit embedding proved
  above), run through the 4-flavour machinery -- used to validate that
  embedding in \texttt{core/BSM/test/test3\_bsm\_sterile.py}.} \\[26pt]
\parbox[t]{3.2cm}{\code{sterile\_\allowbreak{}3p1\_\allowbreak{}bestfit\_\allowbreak{}giunti2017}} &
  \parbox[t]{9.4cm}{$\theta_{14}=8.5^\circ$, $\theta_{24}=7.5^\circ$,
  $\theta_{34}=0$, $\Delta m^2_{41}=1.7\ \mathrm{eV^2}$
  ($\sin^2 2\theta_{14}=0.085$, $\sin^2 2\theta_{24}=0.068$). Global 3+1
  fit combining the reactor and gallium anomalies with solar data
  \parencite{GiuntiMarronePalazzo2017}; sterile CP phases fixed to zero
  (no sensitivity in that data combination).} \\[30pt]
\parbox[t]{3.2cm}{\code{sterile\_\allowbreak{}3p1\_\allowbreak{}benchmark\_\allowbreak{}icecube}} &
  \parbox[t]{9.4cm}{$\theta_{14}=0$, $\theta_{24}\approx9.22^\circ$
  ($\sin^2 2\theta_{24}=0.10$), $\theta_{34}=0$,
  $\Delta m^2_{41}=0.3\ \mathrm{eV^2}$. IceCube atmospheric $\nu_\mu$
  disappearance benchmark \parencite{IceCube2020SterileNumuDisappearance};
  $\theta_{14},\theta_{34}$ set to zero since IceCube atmospheric neutrinos
  are primarily sensitive to $\theta_{24}$.} \\[26pt]
\parbox[t]{3.2cm}{\code{sterile\_\allowbreak{}3p1\_\allowbreak{}two\_\allowbreak{}flavour\_\allowbreak{}limit}} &
  \parbox[t]{9.4cm}{$\theta_{12}=\theta_{13}=\theta_{23}=0$,
  $\theta_{14}=15^\circ$, $\theta_{24}=12^\circ$, $\theta_{34}=0$,
  $\Delta m^2_{41}=1.0\ \mathrm{eV^2}$. A mathematical construction, not a
  data-driven benchmark: switching off every SM active angle makes the
  $\nu_e$ disappearance channel reduce exactly to the textbook 2-flavour
  formula $P_{ee}=1-\sin^2(2\theta_{14})\sin^2(1.267\,\Delta m^2_{41}L/E)$
  for any $\theta_{24}$, since $\theta_{24}$ lives only in $O_4$, which
  leaves the electron index untouched. Used to validate the 2-flavour
  approximation against the full 4-flavour evolutor in isolation, before
  quantifying the residual active-sector suppression present for realistic
  presets such as \code{sterile\_3p1\_bestfit\_giunti2017} (see
  \texttt{sterile2\_kinematics.ipynb} in
  \texttt{notebooks/validation/physics/BSM}).} \\[42pt]
\bottomrule
\end{tabular}
\end{center}
\end{implbox}

\begin{notebox}
\textbf{Approximation exposed by \code{sterile\_3p1\_two\_flavour\_limit}.}
The familiar 2-flavour disappearance formula is only \emph{exact} once the
SM active angles are switched off; for realistic presets (SM angles at
their measured values) the $\nu_\mu$ channel in particular picks up a small
residual correction, because $R_{14}$ and $R_{24}$ both act on the sterile
index and do not commute in general -- $|U_{\mu4}|^2=\sin^2\theta_{24}$
holds exactly only when $\theta_{23}=\theta_{14}=0$. This is a property of
the physics (the standard 2-flavour formula is itself an approximation once
more than one active-sterile angle is non-zero), not a code limitation; the
preset exists specifically to isolate and validate this effect
numerically.
\end{notebox}

\chapter{\texttt{core.numerical} --- general numerical propagation}
\label{ch:core-numerical}

This chapter covers the medium-independent numerical evolutor: exact
matrix-exponential integration of the propagation Hamiltonian over a
piecewise-constant-density discretisation of the path. It makes no
assumption on the size of matter effects and works unmodified for the
Standard Model, NSI, the 3+1 sterile extension, or any combination
(\cref{ch:core-common,ch:core-sm,ch:core-bsm}), since it consumes
\pyfunc{hamiltonian\_flavour} directly rather than any scenario-specific
formula. It is the general-purpose reference method against which the
perturbative approximation (\cref{ch:core-perturbative}) is validated.

\section{\texttt{geometry.py} --- trajectory discretisation}

\modulehead{core/numerical/geometry.py}{\pyclass{Trajectory}: the geometry
container passed from a medium profile to the numerical evolutor, and
\pyfunc{segment\_sample\_points}, the quadrature-rule helper that builds
it.}

\begin{theorybox}
Composing segment evolutors (\cref{ch:core-common}, \texttt{evolutor.py}
section) requires splitting a propagation path into $N$ segments delimited
by $N+1$ boundary points along some geometry coordinate $x$ (a chord
length, a radius, a zenith-angle parametrisation -- whatever the medium
finds natural). For each segment, the numerical core needs exactly two
further pieces of information, both computed once by the medium-specific
module that builds the \pyclass{Trajectory} and never revisited by
\texttt{core.numerical} itself:
\begin{itemize}
  \item the \textbf{dimensionless evolution length} of the segment,
    $\Delta x^{\rm evo}_j = \Delta x_j / L_{\rm scale}$, i.e.\ the physical
    segment length divided by the same \code{evolution\_scale\_m} used to
    non-dimensionalise the Hamiltonian (\cref{ch:core-common},
    \texttt{potential.py} section) -- this is what actually multiplies
    $H$ in the matrix exponential below, not $\Delta x_j$ itself;
  \item a \textbf{sample coordinate} $\bar x_j$ per segment, at which the
    medium's electron density (or any other property entering $H$) is
    evaluated once and treated as \emph{constant} over the whole segment
    -- the piecewise-constant-density approximation that makes the
    per-segment problem exactly solvable (next section).
\end{itemize}
\pyfunc{segment\_sample\_points} implements three quadrature rules for
$\bar x_j$ from the $N+1$ boundary coordinates $x_0,\dots,x_N$:
\begin{equation}
  \bar x_j =
  \begin{cases}
    \tfrac12\bigl(x_j+x_{j+1}\bigr) & \text{\code{"midpoint"}\ (default)}, \\
    x_j & \text{\code{"left"}}, \\
    x_{j+1} & \text{\code{"right"}}.
  \end{cases}
\end{equation}
\keyresult{The midpoint rule is second-order accurate} in the segment
width for a smoothly varying density (the segment's leading-order
mismatch between the true, continuously varying $n_e(x)$ and the
constant $n_e(\bar x_j)$ used on it cancels to first order when $\bar x_j$
is the segment centre, by the same Taylor-expansion argument as the
midpoint rule in quadrature); the endpoint rules (\code{"left"}/
\code{"right"}) are only first-order accurate, and exist mainly to match
legacy or directionally-biased sampling conventions.
\end{theorybox}

\begin{implbox}
\begin{itemize}
  \item \pyclass{Trajectory} (dataclass): \code{x} (shape $(N{+}1,)$, the
    geometry-coordinate boundaries), \code{dx\_evolution} (shape $(N,)$,
    the dimensionless evolution lengths $\Delta x^{\rm evo}_j$ above),
    \code{sample\_x} (shape $(N,)$, the $\bar x_j$ above), and \code{meta}
    (a plain \code{dict} of medium-specific bookkeeping -- layer indices,
    crossing flags, ... -- that \texttt{core.numerical} never interprets).
  \item \pyfunc{segment\_sample\_points(x, method)}: builds \code{sample\_x}
    from \code{x} via the three rules above; broadcasts over any leading
    batch dimensions, and returns an empty (zero-segment) result for a
    single-boundary-point input.
  \item Every medium-specific geometry module (e.g.\
    \texttt{medium.earth.geometry}, \texttt{medium.atmosphere.profile})
    builds its own \code{x} in whatever coordinate is natural for that
    medium, converts it to \code{dx\_evolution} using \emph{that medium's}
    evolution scale, and calls \pyfunc{segment\_sample\_points} to fill
    \code{sample\_x} -- \pyclass{Trajectory} is the single, medium-agnostic
    hand-off point to \texttt{core.numerical}.
\end{itemize}
\end{implbox}

\section{\texttt{evolutor.py} --- matrix-exponential composition}

\modulehead{core/numerical/evolutor.py}{\pyfunc{evolutor\_numerical\_segment}
and \pyfunc{evolutor\_numerical}: exact per-segment matrix exponentials,
composed into the total evolutor.}

\begin{theorybox}
\subsection*{Exact solution on a piecewise-constant-density segment}
The propagation Hamiltonian $H(x)$ varies continuously with $x$ through the
local electron density $n_e(x)$. Given the \pyclass{Trajectory} of the
previous section, \texttt{core.numerical} approximates $H(x)$ as
\textbf{piecewise constant}: on segment $j$ (between boundaries $x_j$ and
$x_{j+1}$), $H(x)\approx H_j$, where
\begin{equation}
  H_j \equiv \text{\pyfunc{hamiltonian\_flavour}}\bigl(\text{oscillation},\,E,\,n_e(\bar x_j)\bigr)
\end{equation}
is the full flavour-basis Hamiltonian (\cref{ch:core-common}) evaluated
once at the segment's sample point $\bar x_j$. \textbf{This is the
\emph{only} approximation in the method}: given a constant $H_j$, the
position-dependent Schr\"odinger equation
$i\,\mathrm{d}S/\mathrm{d}x=H_j\,S$ (\cref{ch:core-common},
\texttt{evolutor.py} section) integrates \emph{exactly} -- not
approximately -- over the segment, via the matrix exponential:
\begin{equation}
  \keyresult{S_j = \exp\bigl(-i\,H_j\,\Delta x^{\rm evo}_j\bigr)}.
\end{equation}
No eigendecomposition, diagonalisation, or perturbative expansion is
involved: \pyfunc{torch.linalg.matrix\_exp} evaluates this directly (via a
scaling-and-squaring / Pad\'e algorithm), batched over every segment and
every other broadcast dimension (energy, angle, ...) simultaneously, and
remains differentiable end-to-end. Because $H_j$ is built with
\pyfunc{hamiltonian\_flavour} -- the single, fully general
Hamiltonian-assembly function (\cref{ch:core-common,ch:core-sm,ch:core-bsm}) --
this formula is correct as written for the 3-flavour Standard Model, NSI,
the 3+1 sterile extension, or any combination, with $H_j$ simply changing
size ($3\times3$ or $4\times4$) and content accordingly; nothing in this
module branches on which scenario is active. In particular, an optional
neutron density $n_n(\bar x_j)$ sampled alongside $n_e(\bar x_j)$ is
forwarded straight through to \pyfunc{hamiltonian\_flavour} as
\code{n\_n\_mol\_cm3}, enabling the 3+1 sterile extension's neutral-current
term (\cref{ch:core-common}) with no change to this module's own logic:
omitting it (the default) reproduces the pre-existing CC-only behaviour
exactly.

\subsection*{Composing the total evolutor}
The $N$ per-segment evolutors are composed into the total operator via the
ordered product already established generically in
\cref{ch:core-common}:
\begin{equation}
  S(x_N,x_0) = S_{N-1}\,S_{N-2}\cdots S_1\,S_0,
\end{equation}
using \pyfunc{compose\_segment\_evolutors} (\cref{ch:core-common}) with its
binary-tree reduction for $O(\log N)$ parallel depth on GPU rather than
$O(N)$ sequential steps. Optionally, \pyfunc{evolutor\_numerical} can
instead return the \emph{history} of partial products
$S(x_j,x_0)$ for every $j=0,\dots,N$ (an accumulating loop rather than the
tree reduction, since every intermediate result is needed): useful for
path-resolved quantities, e.g.\ probabilities evaluated at every crossed
layer of a medium rather than only at the final boundary.
\end{theorybox}

\begin{implbox}
\begin{itemize}
  \item \pyfunc{evolutor\_numerical\_segment(oscillation, E\_MeV, n\_e\_mol\_cm3,
    dx\_evolution, *, n\_n\_mol\_cm3=None, ...)}: builds one $H_j$ per segment
    (the segment dimension is the final axis of
    \code{n\_e\_mol\_cm3}/\code{dx\_evolution}), broadcasting
    \code{oscillation.mass\_spectrum} and \code{antinu} to the segment shape
    via \code{dataclasses.replace} before calling \pyfunc{hamiltonian\_flavour}
    once for the whole segment batch, then exponentiates:
    \code{torch.linalg.matrix\_exp(-1j * H * dx[...,None,None])}.
  \item \pyfunc{evolutor\_numerical(oscillation, E\_MeV, n\_e\_mol\_cm3,
    dx\_evolution, *, n\_n\_mol\_cm3=None, return\_history=False, ...)}: calls
    the above, then either composes the segments into a single operator
    (\pyfunc{compose\_segment\_evolutors}) or accumulates and returns the
    full per-segment history.
  \item Both functions run under \code{@torch.no\_grad()}: the numerical
    evolutor is used for forward propagation, not for parameter inference
    through automatic differentiation.
\end{itemize}
\end{implbox}

\begin{notebox}
\textbf{Accuracy is entirely a discretisation choice, made outside this
module.} \texttt{core.numerical} performs no adaptive step-size control and
no local error estimate: the accuracy of $S(x_N,x_0)$ is fixed entirely by
how finely the \pyclass{Trajectory} passed in discretises the path (how
many segments $N$, and which quadrature rule fixed $\bar x_j$, previous
section) -- both decided by the calling medium-specific module before
\texttt{core.numerical} ever sees the problem. Refining the trajectory
(more, smaller segments) is the only way to improve accuracy; there is no
mechanism to detect that a given $N$ is insufficient. This is the
computational cost paid for generality: unlike the perturbative
approximation (\cref{ch:core-perturbative}), which is only valid in a
restricted regime but is $\mathcal{O}(1)$ in the number of density
evaluations, the numerical method is valid everywhere but its accuracy is
$\mathcal{O}(N)$ in the number of \pyfunc{hamiltonian\_flavour}
evaluations and matrix exponentials.
\end{notebox}

\chapter{\texttt{core.perturbative} --- perturbative approximation}
\label{ch:core-perturbative}

\begin{notebox}[title={This Chapter Describes an Approximation}]
Everything in this chapter is a \textbf{first-order perturbative
approximation} to the exact position-dependent Schr\"odinger equation
(\cref{ch:core-common}), offered as a cheaper alternative to the exact
numerical method of \cref{ch:core-numerical} \emph{within a restricted
validity range} -- it is not a second, independently correct propagation
method.

\textbf{The approximation.} Within one segment, the Hamiltonian is split
into a constant average and a position-dependent residual,
$H(x)=H_{\rm avg}+\delta H(x)$, where $H_{\rm avg}$ uses the
segment-\emph{average} density and $\delta H(x)$ captures the fluctuation
around it. The evolutor is expanded to \textbf{first order} in $\delta H$:
$S\approx U_0+U_1$, with $U_0=\exp(-iH_{\rm avg}L)$ exact for the constant
average (\cref{sec:perturbative-evolutor}), and $U_1$ the first Dyson-series
correction from $\delta H$.

\textbf{Range and conditions of validity.} The truncation is controlled by
the phase the \emph{residual} potential accumulates over the segment,
$|\delta V(x)|\cdot L$: the approximation is trustworthy while this stays
small compared to 1 throughout the segment (equivalently, while the
neglected second- and higher-order Dyson terms stay small relative to
$U_1$). In practice this means: (a) the profile model
(\cref{sec:perturbative-models}) must fit the true density well enough,
and/or segments must be short enough, that the density does not deviate
strongly from its local average within any one segment; (b) accuracy is
expected to degrade for segments spanning a large density contrast (e.g.\ a
sharp discontinuity treated as a single segment). \textbf{No automatic
error estimate or order check is built into this code path} -- unlike
\cref{ch:core-numerical}'s explicit accuracy note, there is no signal here
that a given (segment, profile-model) choice has left the valid regime; it
must be established externally, typically by cross-validating against the
exact numerical evolutor of \cref{ch:core-numerical} for representative
parameters before trusting a perturbative configuration in production.
\end{notebox}

This chapter covers \texttt{spectral.py} (the closed-form spectral
decomposition every perturbative segment evolutor is built from),
\texttt{evolutor.py} (the zeroth- and first-order evolution operators
themselves), and the family of interchangeable density-profile models that
supply the analytic inputs the first-order correction needs.

\section{\texttt{spectral.py} --- closed-form spectral decomposition}

\modulehead{core/perturbative/spectral.py}{Trace/traceless splitting,
characteristic-polynomial invariants, and closed-form spectral projectors
for a $3\times3$ or $4\times4$ Hermitian Hamiltonian.}

\begin{theorybox}
\subsection*{Trace and traceless parts}
Any Hermitian $H$ splits uniquely into a multiple of the identity and a
traceless remainder,
\begin{equation}
  \keyresult{T = H - \frac{\Tr(H)}{N}\,I_N}, \qquad N=H.\text{shape}[-1]\in\{3,4\}.
\end{equation}
Since $I_N$ commutes with everything, $H$ and $T$ share the same
eigenvectors, and $H$'s eigenvalues are simply $T$'s shifted uniformly by
$\Tr(H)/N$. Working with $T$ instead of $H$ removes the sub-leading term of
the characteristic polynomial (no $\lambda^{N-1}$ term, since $\Tr(T)=0$),
which is what makes the closed-form projector formulas below tractable.

\subsection*{Characteristic-polynomial invariants}
For $N=3$, $T$'s characteristic polynomial is $\lambda^3+c_1\lambda+c_0=0$
with
\begin{equation}
  c_1 = -\tfrac12\Tr(T^2), \qquad c_0 = -\tfrac13\Tr(T^3).
\end{equation}
The $N=4$ projector formula additionally needs
$e_3=\tfrac13\Tr(T^3)=-c_0$ (kept as a separate function in the code to
avoid a sign-confusion bug, since both invariants are used together for
$N=4$). The eigenvalues $\lambda_a$ themselves have no closed-form solver
in this module -- they are obtained numerically via
\pyfunc{torch.linalg.eigvalsh}; only the \emph{projectors} below are
closed-form, given the eigenvalues.

\subsection*{Closed-form spectral projectors}
The spectral theorem writes $T=\sum_a\lambda_a M_a$ with
$M_a=\prod_{b\neq a}\frac{T-\lambda_b I}{\lambda_a-\lambda_b}$ (the
Lagrange/Sylvester interpolation projector). Using Cayley--Hamilton and
$\Tr(T)=0$ to eliminate the "other" eigenvalues $\lambda_{b\neq a}$ from
this product, it reduces to an explicit polynomial in $T$ (and $T^2$, plus
$T^3$ for $N=4$) with coefficients depending only on $\lambda_a$ and the
invariants above:
For $N=3$:
\begin{equation}
  \keyresult{M_a = \frac{(\lambda_a^2+c_1)\,I + \lambda_a T + T^2}{3\lambda_a^2+c_1}}.
\end{equation}
For $N=4$:
\begin{equation}
  \keyresult{M_a = \frac{T^3+\lambda_a T^2+(c_1+\lambda_a^2)T+(\lambda_a^3+c_1\lambda_a-e_3)\,I}{4\lambda_a^3+2c_1\lambda_a-e_3}}.
\end{equation}
Both denominators equal $p'(\lambda_a)$, the derivative of the
characteristic polynomial at that root: they vanish exactly when
$\lambda_a$ is a repeated eigenvalue. Both formulas were verified
symbolically and numerically (against direct Lagrange interpolation and
against \pyfunc{torch.linalg.eigh} eigenvector outer products) to match to
floating-point precision.
\end{theorybox}

\begin{notebox}
\textbf{Degeneracy fallback.} Near a (near-)degenerate spectrum the
closed-form projector formulas lose precision well before their
denominator literally reaches zero, since $p'(\lambda_a)\to0$
continuously. \pyfunc{\_spectral\_degeneracy\_mask} flags batch entries
whose minimum pairwise eigenvalue gap is small \emph{relative} to the
spectrum scale (or whose scale itself is numerically zero, the $T\approx0$
case, where the relative test would be $0/0$); flagged entries have their
projectors replaced wholesale by $M_a=v_av_a^\dagger$ from a direct
\pyfunc{torch.linalg.eigh(T)} call -- exact and stable by construction,
independent of the closed-form formula's conditioning, computed only when
at least one batch entry actually needs it. This module runs entirely
under \code{@torch.no\_grad()}, so the well-known gradient instability of
degenerate eigenvectors is not a concern here (the perturbative path is
used for forward propagation, not for differentiating through the
spectrum).
\end{notebox}

\begin{implbox}
\begin{itemize}
  \item \pyfunc{hamiltonian\_traceless(H)}: returns $(T,\Tr H)$.
  \item \pyfunc{hamiltonian\_traceless\_c0/c1/e3(T)}: the invariants above.
  \item \pyfunc{hamiltonian\_traceless\_eigenvalues(T)}: $\lambda_a$ via
    \pyfunc{torch.linalg.eigvalsh}.
  \item \pyfunc{hamiltonian\_spectral\_projectors\_traceless(T, ...)}:
    builds $M_a$ (closed form, with the \pyfunc{eigh} fallback above);
    returns $(M,\lambda,c_1)$.
  \item \pyfunc{hamiltonian\_spectral\_data(H)}: the one-stop entry point
    used by \texttt{evolutor.py} -- computes $T$, $\Tr H$, $\lambda$, $c_1$,
    and $M$ in one call, reusing intermediate products ($T^2$, $T^3$)
    across every invariant and the projector formula to avoid redundant
    matrix multiplications.
  \item \pyfunc{spectral\_projector\_residuals(M, T, lam)}: diagnostic
    residual norms (completeness $\sum_aM_a=I$, idempotency
    $M_a^2=M_a$, orthogonality $M_aM_b=0$, unit trace, and
    reconstruction $\sum_a\lambda_aM_a=T$) -- used both by the test suite
    and by the degeneracy instrumentation itself.
\end{itemize}
\end{implbox}

\section{\texttt{evolutor.py} --- perturbative segment evolution}
\label{sec:perturbative-evolutor}

\modulehead{core/perturbative/evolutor.py}{\pyfunc{evolutor\_zero\_order}
(exact, for the segment-average Hamiltonian) and
\pyfunc{evolutor\_first\_order} (the first Dyson-series correction from the
density residual), combined into the segment evolutor below.}

\begin{theorybox}
\subsection*{Zeroth order: exact evolution under the average Hamiltonian}
Given the spectral decomposition $H_{\rm avg}=T+\Tr(H_{\rm avg})/N\cdot I$
from the previous section, the spectral theorem gives $\exp(-iH_{\rm avg}L)$
\emph{exactly} (not perturbatively) as
\begin{equation}
  \keyresult{U_0 = \sum_a \exp\!\Bigl[-i\Bigl(\lambda_a+\tfrac{\Tr H_{\rm avg}}{N}\Bigr)L\Bigr]\,M_a},
\end{equation}
since conjugating by the identity-shift only shifts every eigenvalue by the
same constant $\Tr(H_{\rm avg})/N$, leaving the projectors $M_a$
untouched. $U_0$ is the entire answer if the segment's density were exactly
constant at its average value; every approximation in this chapter lives
in the correction below.

\subsection*{First order: the Dyson-series correction from the residual}
Writing $H(x)=H_{\rm avg}+\delta H(x)$ with $\delta H(x)=\delta V(x)\,P$ (a
position-independent \emph{direction} matrix $P$ scaled by the
position-dependent residual potential $\delta V(x)=V(x)-V_{\rm avg}$; see
below for $P$), first-order time-dependent perturbation theory in the
interaction picture defined by $U_0$ gives
\begin{equation}
  U_1 = -i\int_{x_1}^{x_2} U_0(x_2,x')\,\delta H(x')\,U_0(x',x_1)\,\mathrm{d}x'.
\end{equation}
Expanding $U_0$ in the spectral basis, $U_0(x_2,x')=\sum_a
e^{-i\lambda_a(x_2-x')}M_a$, and inserting $\delta H(x')=\delta V(x')P$
turns the operator integral into a sum over eigenvalue \emph{pairs} of a
scalar oscillatory integral of the residual density weighted by the matrix
$M_aPM_b$:
\begin{equation}
  \keyresult{U_1 = -i\sum_{a,b} I_{ab}\;M_a\,P\,M_b},
  \qquad
  I_{ab} = \int_{x_1}^{x_2}\delta V(x')\,e^{-i\lambda_a(x_2-x')-i\lambda_b(x'-x_1)}\,\mathrm{d}x'.
\end{equation}
$I_{ab}$ (a ``spectral oscillatory integral'' of the residual profile) is
computed analytically by the active profile model
(\cref{sec:perturbative-models}, \pyfunc{residual\_integral}); it is the
only place the specific density parametrisation enters the perturbative
machinery -- \texttt{evolutor.py} itself never sees a profile's functional
form, only $I_{ab}(\lambda_a,\lambda_b)$.

\textbf{The direction matrix $P$.} For the plain Standard Model or 3+1
sterile case (no NSI), only the electron-flavour entry of $H_{\rm mat}$
varies with density, so $P=\diag(1,0,\dots,0)$ and the sandwich $M_aPM_b$
reduces to the rank-1 outer product of $M_a$'s first column with $M_b$'s
first row -- the closed-form shortcut used when the caller passes no
explicit $P$. When NSI is active, the coupling $\epsilon$ is
position-independent (only the electron density magnitude, not the NSI
coupling itself, varies along the segment), so the same first-order
structure applies with
\begin{equation}
  P = O^\dagger\bigl(\diag(1,0,\dots,0)+\epsilon\bigr)O
\end{equation}
in the reduced basis (\cref{ch:core-bsm}) -- fixed once per segment (built
at $V{=}1$) and reused for both $H_{\rm avg}$'s matter term and this
first-order sandwich, since it no longer reduces to the SM rank-1
shortcut once $\epsilon$ breaks the electron-only structure.

\textbf{The neutral-current direction $P_{NC}$.} When the optional 3+1
sterile neutral-current term is active (\cref{ch:core-common}), it
contributes a second, wholly independent first-order term, since the
electron- and neutron-density residuals fluctuate independently along the
segment:
\begin{equation}
  \keyresult{U_1 = -i\sum_{a,b} I_{ab}^{(e)}\;M_a\,P_{CC}\,M_b
             \;-i\sum_{a,b} I_{ab}^{(n)}\;M_a\,P_{NC}\,M_b},
\end{equation}
where $I_{ab}^{(n)}$ is the same spectral oscillatory integral as
$I_{ab}^{(e)}$ but of the \emph{neutron}-density residual
($\pyfunc{residual\_integral\_neutron}$, below), and
$P_{NC}=O^\dagger\diag(0,0,0,-1)\,O$ is the (always rotated) direction
matrix derived in \cref{ch:core-common}. Unlike $P_{CC}$, $P_{NC}$ never
reduces to a rank-1 shortcut regardless of whether NSI is active, since it
never commutes with $O$; the CC term above may still use its own rank-1
shortcut independently, since the two contributions are additive and
computed separately. Omitting the neutral-current term (the default)
recovers the CC(+NSI)-only expression exactly.

\subsection*{Combining orders, and the zero-length convention}
The segment evolutor is $U=U_0+U_1$
(\pyfunc{evolutor\_perturbative\_from\_H}), with two guards applied by the
code around this core formula: (1) if the profile model reports no
perturbation for a segment (a genuinely constant density, e.g.\
\pyfunc{has\_perturbation()} is \code{False}), $U_1$ is skipped entirely
--- both because it is identically zero and to avoid computing
$I_{ab}$ needlessly; (2) zero-length segments are forced to the identity
operator (\pyfunc{enforce\_identity\_for\_zero\_length}) regardless of the
formula above, since $L=0$ is a genuine edge case (e.g.\ a trajectory that
grazes but does not cross a layer) rather than a regime where the
perturbative expansion itself is being tested.
\end{theorybox}

\begin{notebox}
\textbf{Hermitian conjugate, again.} \pyfunc{evolutor\_perturbative\_segment}
builds $H_{\rm kin}$ via \pyfunc{hamiltonian\_kinetic\_reduced}
(\cref{ch:core-common}), which -- exactly as flagged there and in
\cref{ch:core-bsm} -- always uses $U_{\rm red}\diag(k)U_{\rm
red}^\dagger$, never the plain transpose. This matters even more visibly
here: for the 3+1 sterile case the code computes the segment's flavour-basis
trace directly as $\Tr(H_{\rm kin}+H_{\rm mat})$ rather than assuming
$\Tr(H_{\rm kin})=\sum_ik_i$ (only guaranteed for a \emph{real} $U_{\rm
red}$) or $\Tr(H_{\rm mat})=V$ (only true for the plain
$\diag(1,0,\dots,0)$ direction) -- both shortcuts used safely in the
pure-SM fast path, both invalid once the sterile CP phase $\delta_{14}$ or
a generic NSI $\epsilon$ is active.
\end{notebox}

\begin{implbox}
\begin{itemize}
  \item \pyfunc{evolutor\_zero\_order(H, L, ...)}: the exact $U_0$ formula
    above; can also return the spectral data for reuse by the first-order
    call, avoiding a second diagonalisation.
  \item \pyfunc{evolutor\_first\_order(M, lam, trace\_H, profile\_model, ..., P=None, P\_nc=None)}:
    the $U_1$ formula above, using the rank-1 shortcut for the CC term when
    \code{P} is \code{None}; when \code{P\_nc} is also supplied, adds the
    independent neutral-current sandwich term using
    \pyfunc{profile\_model.residual\_integral\_neutron} and
    \pyfunc{matter\_potential\_nc}.
  \item \pyfunc{evolutor\_perturbative\_from\_H(H, L, profile\_model, ..., P=None, P\_nc=None)}:
    combines both orders, with the no-perturbation and zero-length guards;
    the no-perturbation check also consults
    \pyfunc{profile\_model.has\_perturbation\_neutron()} when \code{P\_nc}
    is supplied.
  \item \pyfunc{evolutor\_perturbative\_segment(oscillation, E\_MeV,
    profile\_model, ..., include\_matter\_nc=False)}: the user-facing entry
    point -- builds the reduced Hamiltonian (SM, sterile, NSI, or any
    combination, mirroring \cref{ch:core-common,ch:core-bsm}) from the
    profile model's segment average, then evolves it perturbatively via the
    function above. The plain Standard Model case takes a dedicated fast
    path (bit-identical output, avoiding the general-$N$ machinery's
    overhead). \code{include\_matter\_nc=True} additionally requires
    \code{profile\_model.potential\_n} to be set (i.e.\ the profile was
    built with neutron-density coefficients, next section) and is silently
    ignored for a 3-flavour \code{oscillation.pmns}, mirroring
    \pyfunc{hamiltonian\_matter\_reduced}'s own convention
    (\cref{ch:core-common}).
\end{itemize}
\end{implbox}

\section{Comparing the perturbative density-profile models}
\label{sec:perturbative-models}

\modulehead{core/perturbative/models/*}{Three interchangeable profile
models, each supplying \code{average}, \code{length}, \code{zero\_mask},
\code{potential}, \pyfunc{residual\_integral}, and
\pyfunc{has\_perturbation} to the model-agnostic evolutor above.}

\begin{implbox}
None of the mathematics in \texttt{evolutor.py} or \texttt{spectral.py}
changes between models: they differ purely in \emph{how} the segment's
density profile $n_e(x)$ is parametrised, and hence in the closed-form
expression used for $I_{ab}$.

\begin{center}
\small
\begin{tabular}{@{}ll@{}}
\toprule
Model & Density parametrisation and usage \\
\midrule
\parbox[t]{3.2cm}{\code{even\_power}} &
  \parbox[t]{9.4cm}{$n_e(x)=a+bx^2+cx^4+\delta_1x^6+\dots$ (even powers only,
  arbitrary order) in the raw trajectory coordinate $x$. Appropriate
  whenever the segment is symmetric about the trajectory's point of
  closest approach to a spherically symmetric medium (e.g.\ an
  Earth-crossing chord through one shell), which forces $n_e(x)=n_e(-x)$
  and eliminates odd powers identically. The general-order model used for
  Earth propagation.} \\[30pt]
\parbox[t]{3.2cm}{\code{prem}} &
  \parbox[t]{9.4cm}{$n_e(x)=a+bx^2$: the 2-coefficient special case of
  \code{even\_power} truncated at the quadratic term, matching the PREM
  reference model's own native linear-in-$r^2$ parametrisation per shell
  \parencite{DziewonskiAnderson1981}. A lighter-weight specialisation, not
  an independent formulation -- it exists to avoid the general
  \code{even\_power} machinery's overhead when only two coefficients are
  ever needed.} \\[30pt]
\parbox[t]{3.2cm}{\code{atmosphere}} &
  \parbox[t]{9.4cm}{$n_e(q)=\sum_kc_kq^k$ (every power, even and odd) in a
  \emph{centred and rescaled} local coordinate $q=(x-\text{centre})/\text{half-width}$.
  Atmosphere segments are not symmetric about their own centre the way an
  Earth chord is, so odd terms are required; the local coordinate keeps the
  fitted coefficients well-conditioned regardless of the segment's absolute
  position or width.} \\[30pt]
\bottomrule
\end{tabular}
\end{center}

Every model computes its own \pyfunc{residual\_integral} in closed form for
ordinary (non-degenerate) eigenvalue gaps, switching to a Taylor expansion
in $\lambda_a-\lambda_b$ near degeneracy to avoid the numerical instability
of the closed form as the oscillation frequency approaches zero -- the same
degeneracy-awareness pattern as the spectral projectors themselves
(previous section), applied independently at the profile-model level.
\pyfunc{perturbative\_profile\_selection} (\code{core/perturbative/models/model\_selection.py})
maps a public model name (\code{"even\_power"}, \code{"prem"}, ...) to the
concrete layered-profile class, so medium code can request a model without
importing model-specific classes directly.

\textbf{Optional neutron-density coefficients.} All three segment classes
optionally accept a second, independent set of coefficients for the neutron
density $n_n$ -- same functional form and coordinate as the electron-density
fit above, fitted or supplied separately -- exposing \code{average\_n},
\code{potential\_n}, \pyfunc{residual\_integral\_neutron}, and
\pyfunc{has\_perturbation\_neutron} alongside the electron-only attributes,
required by \pyfunc{evolutor\_perturbative\_segment}'s
\code{include\_matter\_nc=True} (previous section). They default to
\code{None}, reproducing the pre-existing CC-only behaviour exactly; the
layered profile classes that build these segments
(\code{EvenPowerProfileLayered}, \code{PremTabulatedProfile}) load them from
a composition-derived neutron-density table when constructed with
\code{include\_neutron=True} (\cref{ch:earth}), and
\code{AtmospherePolynomialProfile} fits them from a separate neutron-density
sample at the same nodes (\cref{ch:atmosphere}).
\end{implbox}


% =============================================================================
\part{Physical Propagation Media --- \texttt{medium/}}
\chapter{\texttt{medium.vacuum} --- coherent propagation in vacuum}
\label{ch:vacuum}

\texttt{medium.vacuum} is the simplest propagation medium in the project:
$n_e\equiv0$ everywhere, so the matter term of the Hamiltonian
(\cref{ch:core-common}) never enters and every formula reduces to the
closed-form vacuum evolutor already derived, in general terms, in the
coherent-propagation scheme of \cref{ch:core-sm}. Its role is twofold:

\begin{itemize}
  \item \textbf{Base case and test bed.} With no density profile to model
    and a fully closed-form evolutor, this is the cheapest possible path
    through the entire propagation stack (\pyclass{PMNS} $\to$
    \pyfunc{kinetic\_eigenvalue\_vector} $\to$ evolutor $\to$ probability
    $\to$ flux, \cref{ch:core-common}), with an exact analytic answer to
    check every other piece against. It is, in practice, the first thing
    exercised when validating a new oscillation preset, a new BSM
    extension, or a refactor of the shared infrastructure: if
    \texttt{medium.vacuum} disagrees with the legacy \texttt{peanuts}
    reference (\cref{sec:vacuum-validation}), nothing built on top of it
    can be trusted either.
  \item \textbf{Physical use cases.} Vacuum propagation is the correct
    description whenever matter effects are genuinely negligible along the
    whole path: reactor and other short-baseline experiments, or any
    scenario where the accumulated matter phase $V\cdot L$ is small enough
    to ignore relative to the vacuum kinetic phases. It is also the
    $n_e\to0$ reference limit that the matter-propagation media of the
    following chapters must recover exactly in that limit.
\end{itemize}

\begin{notebox}
\textbf{The four-module pattern.} \texttt{medium.vacuum} fixes a module
layout that every subsequent medium in this part of the document repeats,
adapted to that medium's own needs: an \texttt{evolutor.py} building the
medium-specific $S$ from \pyclass{OscillationParameters} and geometry
inputs; a \texttt{probability.py} converting that evolutor into transition
and final-state probabilities via the shared
\texttt{core.common.probability} conventions (\cref{ch:core-common}); a
\texttt{flux.py} applying \texttt{core.common.flux} normalisations on top;
and, where a legacy NumPy reference implementation exists to check against,
a \texttt{validation.py} comparing the two numerically. Media with
additional physical structure add further modules on top of this
skeleton (\texttt{profile.py} for a density profile,
\texttt{geometry.py} for path geometry, \texttt{exposure\_*.py} for
detector exposure, ...) without changing this core four-piece contract.
\end{notebox}

\section{\texttt{evolutor.py} --- the vacuum evolutor}

\modulehead{medium/vacuum/evolutor.py}{\pyfunc{vacuum\_evolutor} and
\pyfunc{vacuum\_evolved\_state}: the closed-form coherent vacuum operator
and its action on a state.}

\begin{theorybox}
With $H_{\rm mat}=0$, the reduced and flavour Hamiltonians coincide with
the pure kinetic term, and the general coherent-propagation scheme of
\cref{ch:core-sm} is exact and complete on its own (no reduction trick is
needed, since there is no matter term to factor out):
\begin{equation}
  \keyresult{S(L,E) = U\,\diag\!\bigl(e^{-ik_1x},e^{-ik_2x},e^{-ik_3x}\bigr)\,U^\dagger},
  \qquad x = \frac{L\,[\mathrm{m}]}{L_{\rm scale}},
\end{equation}
where $U$ is the full PMNS matrix (\cref{ch:core-common}), $k_i$ the
dimensionless kinetic eigenvalues (\cref{ch:core-common},
\texttt{potential.py} section), and the baseline $L$ (given in km) is
converted to metres and non-dimensionalised by the same
\code{evolution\_scale\_m} used throughout the project. This is the same
formula the sterile extension recovers automatically for $n=4$
(\cref{ch:core-bsm}), simply by passing a $4\times4$ \pyclass{PMNS\_sterile}
object in place of \pyclass{PMNS\_SM}; nothing in this module is
3-flavour-specific.
\end{theorybox}

\begin{implbox}
\begin{itemize}
  \item \pyfunc{vacuum\_evolutor(oscillation, E\_MeV, L\_km, ...)}: builds
    $S$ exactly as above; broadcasts energy, baseline, and mass-splitting
    batch shapes together before the phase multiplication.
  \item \pyfunc{vacuum\_evolved\_state(nustate, oscillation, E\_MeV, L\_km, ...)}:
    applies \pyfunc{core.common.evolutor.apply\_evolutor\_to\_state}
    (\cref{ch:core-common}) to a coherent flavour-basis initial state.
  \item \code{legacy\_precision} is accepted for call-site symmetry with
    every matter-propagation evolutor in later chapters but is a genuine
    no-op here: there is no matter potential for the flag to affect.
\end{itemize}
\end{implbox}

\section{\texttt{probability.py} --- vacuum probabilities}

{\sloppy\modulehead{medium/vacuum/probability.py}{\pyfunc{vacuum\_\allowbreak{}probability\_\allowbreak{}transition}
(the full transition-probability matrix), \pyfunc{vacuum\_\allowbreak{}probability\_\allowbreak{}state}
(final-state probabilities for a given initial state), and
\pyfunc{vacuum\_\allowbreak{}probability\_\allowbreak{}integrated} (spectrum-weighted energy average).}}

\begin{implbox}
\begin{itemize}
  \item \pyfunc{vacuum\_probability\_transition(oscillation, E\_MeV, L\_km, ...)}:
    builds $S$ and returns $P_{\beta\alpha}=|S_{\beta\alpha}|^2$ via
    \pyfunc{core.common.probability.probability\_transition}
    (\cref{ch:core-common}) -- the full $3\times3$ (or $4\times4$) channel
    matrix, no initial state required.
  \item \pyfunc{vacuum\_probability\_state(nustate, oscillation, E\_MeV, L\_km, massbasis, ...)}:
    dispatches to the coherent or incoherent formula of
    \cref{ch:core-common} (\texttt{evolutor.py}/\texttt{probability.py}
    sections) via \pyfunc{probability\_state}, according to
    \code{massbasis}: \code{False} evolves \code{nustate} as coherent
    flavour amplitudes ($P_\beta=|(S\,\psi)_\beta|^2$); \code{True} (the
    default here) treats it as incoherent mass-basis weights
    ($P_\beta=\sum_i|(S\,U)_{\beta i}|^2w_i$). Either way the function
    returns a final flavour-probability vector, never a flux.
  \item \pyfunc{vacuum\_probability\_integrated(nustate, oscillation, E\_MeV, L\_km, spectrum, ...)}:
    builds \pyfunc{vacuum\_probability\_state} on an energy grid and averages
    it with \pyfunc{core.common.probability.probability\_integrated}
    (\cref{ch:core-common}), weighted by an explicit production
    \code{spectrum} -- there is no default weight, since a flat spectrum
    over an arbitrary grid is a modelling choice.
\end{itemize}
\end{implbox}

\section{\texttt{flux.py} --- vacuum flux}

\modulehead{medium/vacuum/flux.py}{\pyfunc{vacuum\_\allowbreak{}flux\_\allowbreak{}state} and
\pyfunc{vacuum\_\allowbreak{}flux\_\allowbreak{}integrated}: apply \texttt{core.common.flux}
normalisation and energy integration to
\pyfunc{vacuum\_\allowbreak{}probability\_\allowbreak{}state}'s output.}

\begin{implbox}
\pyfunc{vacuum\_flux\_state(nustate, oscillation, E\_MeV, L\_km, flux, spectrum, massbasis, ...)}
adds no new physics: it calls \pyfunc{vacuum\_probability\_state} above,
then \pyfunc{core.common.flux.flux\_state} (\cref{ch:core-common}) to apply
the flux normalisation and optional spectral weight. This is the minimal
instance of the flux pattern that every later medium's \texttt{flux.py}
extends. \pyfunc{vacuum\_flux\_integrated(...)} builds
\pyfunc{vacuum\_flux\_state} on an energy grid and integrates it with
\pyfunc{core.common.flux.flux\_integrated}, returning an unnormalised
physical rate rather than the normalised average of
\pyfunc{vacuum\_probability\_integrated}.
\end{implbox}

\section{\texttt{validation.py} --- cross-check against legacy \texttt{peanuts}}
\label{sec:vacuum-validation}

\modulehead{medium/vacuum/validation.py}{Numerical comparison against the
bundled legacy, NumPy-based \texttt{peanuts} package.}

\begin{implbox}
The legacy vacuum implementation is scalar (one energy, one baseline, one
initial state at a time); this module evaluates both implementations on
matching scalar inputs and compares the results directly, establishing the
pattern reused by \texttt{medium.solar.validation} and
\texttt{medium.earth.validation} (\cref{ch:solar,ch:earth}).
\begin{itemize}
  \item \pyfunc{ensure\_legacy\_importable()} / \pyfunc{legacy\_modules()}:
    locate and import the bundled \code{peanuts} package (kept in the
    repository specifically as this validation reference, not as a runtime
    dependency of \texttt{tpeanuts} itself).
  \item \pyfunc{legacy\_pmns\_from\_torch(...)}: rebuilds a legacy
    \code{peanuts.pmns.PMNS} object from a torch \pyclass{PMNS} instance's
    parameters, so both implementations start from identical inputs.
  \item \pyfunc{compare\_vacuum\_probability\_state\_with\_legacy(...)} and
    \pyfunc{compare\_vacuum\_evolved\_\allowbreak{}state\_with\_\allowbreak{}legacy(...)}:
    run both implementations and return their outputs (and/or a residual)
    for the test suite to assert against.
\end{itemize}
\end{implbox}

\chapter{\texttt{medium.solar} --- adiabatic MSW propagation in the Sun}
\label{ch:solar}

Solar neutrinos are produced deep in the high-density solar core and exit
through a smoothly, monotonically decreasing electron density to vacuum.
This chapter's physics is dominated by one fact exploited throughout: the
density gradient is slow enough, on the scale of the relevant MSW
resonances, that the \textbf{adiabatic approximation} applies almost
everywhere in the standard oscillation parameter space -- a neutrino
produced in a given local matter eigenstate stays in it as the density
drops, so propagation reduces to an incoherent mixture of eigenstates
rather than a coherent evolutor (see the note in
\cref{sec:solar-no-evolutor}).

\section{\texttt{profile.py} --- the tabulated solar model}
\label{sec:solar-profile}

\modulehead{medium/solar/profile.py}{\pyclass{SolarProfile}: tabulated
solar radius, electron density, per-source production fractions, and
per-source total fluxes.}

\begin{implbox}
\pyclass{SolarProfile} stores four tabulated quantities on a common radius
grid $r/R_\odot\in[0,1]$: the electron density $n_e(r)$ (mol/cm$^3$), a
production-fraction profile $f_{\rm source}(r)$ per neutrino-producing
reaction (pp, $^8$B, $^7$Be, ...), and a total integrated flux
$\Phi_{\rm source}$ per source. \pyfunc{SolarProfile.default} loads these
from the package's bundled CSV tables (\code{tpeanuts/data/solar}), whose
paths are fixed by \code{tpeanuts.config.default.solar\_model\_filename}
and \code{solar\_fluxes\_filename} -- currently the \textbf{zenodo
SF3-AGSS09} tables, a standard-solar-model dataset combining helioseismic
constraints with the AGSS09 heavy-element abundances
\parencite{Vinyoles2017B16}, tabulated over the full solar radius (2001
rows). \pyclass{SolarParameters} allows overriding both file paths
explicitly, but the defaults are the project's validated reference
composition and should not be changed casually.
\begin{itemize}
  \item \pyfunc{electron\_density(r\_query)} / \pyfunc{production\_fraction(source, r\_query)}:
    linear interpolation on the tabulated grid (\pyfunc{util.math.interp1d\_linear}),
    clamped to the boundary values outside $[0,1]$.
  \item \pyfunc{source\_fractions(sources)}: stacks one or several tabulated
    (not interpolated) production-fraction profiles.
  \item \pyfunc{normalized\_fraction(source)}: normalises a production
    profile to unit trapezoidal integral over $r$
    (\pyfunc{util.math.normalize\_trapz}).
  \item \pyfunc{flux(source)}: the total (radius-integrated) flux
    normalisation for one source, as tabulated.
  \item \code{use\_LZ}: boolean flag (default \code{False}) selecting
    whether \texttt{probability.py} (\cref{sec:solar-probability}) applies
    the Landau-Zener correction of \cref{sec:solar-lz}.
\end{itemize}
\end{implbox}

\begin{notebox}
\textbf{Validity of the tabulated model.} The bundled solar model is a
\emph{standard} solar model: a hydrostatic, helioseismically calibrated
stellar-structure calculation, not a direct measurement. Its accuracy is
therefore tied to the input microphysics (opacities, nuclear reaction
rates, heavy-element abundances) and is typically quoted at the few-percent
level for the core density and composition profiles that set $n_e(r)$ and
the production-fraction shapes; solar-neutrino and helioseismic data
provide independent, largely consistent cross-checks of this precision.
Because $n_e(r)$ only ever enters through the dimensionless ratio $V_k$
(\cref{sec:solar-matter-mixing}), a uniform rescaling error in the model's
overall density normalisation would bias the resonance radius, not the
qualitative adiabatic/non-adiabatic character of the propagation.
\end{notebox}

\section{\texttt{matter\_mixing.py} --- the adiabatic matter-mixing angles}
\label{sec:solar-matter-mixing}

\modulehead{medium/solar/matter\_mixing.py}{Closed-form matter-modified
mixing angles $\theta_{12}^M(r)$, $\theta_{13}^M(r)$ used by the adiabatic
solar-probability approximation.}

\begin{theorybox}
Because $\Delta m^2_{3\ell}\gg\Delta m^2_{21}$ and $\theta_{13}$ is small,
the solar MSW problem separates hierarchically into two sequential
effective 2-flavour resonances (the $1$--$3$ sector, then the $1$--$2$
sector), each with the standard closed-form matter-mixing-angle solution
\parencite{GiuntiKim2007} -- no numerical diagonalisation of the full
$3\times3$ Hamiltonian is needed in the Standard Model case. Define the
dimensionless matter-potential ratio and the effective atmospheric
splitting
\begin{align}
  V_k(\Delta m^2,E,n_e) &= \frac{V(n_e)}{k(\Delta m^2,E)}
  = \frac{2E\hbar c\,V_{CC}}{\Delta m^2}, \\
  \Delta m^2_{ee} &= \cos^2\theta_{12}\,\Delta m^2_{31} + \sin^2\theta_{12}\,\Delta m^2_{32},
\end{align}
using $V$ and $k$ exactly as defined in \cref{ch:core-common}
(\texttt{potential.py} section) -- $V_k$ is simply their ratio, evaluated
here for a given splitting rather than embedded in the Hamiltonian. The
matter-modified reactor angle is
\begin{equation}
  \keyresult{\theta_{13}^M = \tfrac12\arccos\!\left(
    \frac{\cos2\theta_{13}-V_k(\Delta m^2_{ee},E,n_e)}
         {\sqrt{\bigl(\cos2\theta_{13}-V_k(\Delta m^2_{ee},E,n_e)\bigr)^2+\sin^22\theta_{13}}}
  \right)},
\end{equation}
and, using $\theta_{13}^M$ as an input, the matter-modified solar angle is
\begin{equation}
  \keyresult{\theta_{12}^M = \tfrac12\arccos\!\left(
    \frac{\cos2\theta_{12}-V_k'}
         {\sqrt{(\cos2\theta_{12}-V_k')^2+\sin^22\theta_{12}\cos^2(\theta_{13}^M-\theta_{13})}}
  \right)},
\end{equation}
\begin{equation}
  V_k' = V_k(\Delta m^2_{21},E,n_e)\cos^2\theta_{13}^M
       + \frac{\Delta m^2_{ee}}{\Delta m^2_{21}}\sin^2(\theta_{13}^M-\theta_{13}).
\end{equation}
Each resonance occurs where the corresponding $V_k=\cos2\theta$, i.e.\
where the vacuum mixing angle is driven to $\pi/4$ by the matter
potential; away from resonance $\theta^M\to\theta$ (low density) or
$\theta^M\to\pi/2$ (high density), recovering the vacuum and
fully-suppressed limits respectively.
\end{theorybox}

\begin{notebox}
\code{legacy\_precision=True} reproduces the legacy \texttt{peanuts}
combined prefactor for $V_k$ ($3.868\times10^{-7}/2.533$, both factors
independently rounded to 4 significant digits in the original
implementation) rather than deriving the matter and kinetic prefactors
separately from full-precision constants (\cref{ch:core-common}); the two
choices differ by a relative $\sim3.6\times10^{-4}$ -- an order of
magnitude larger than the matter-potential-only rounding error of
\cref{ch:core-common}, since both legacy prefactors compound. Intended
purely for bit-comparable validation against legacy \texttt{peanuts}.
\end{notebox}

\begin{implbox}
\pyfunc{Vk(Deltam2, E, ne, antinu, legacy\_precision)},
\pyfunc{DeltamSqee(oscillation)}, \pyfunc{th13\_M(oscillation, E, ne, ...)},
\pyfunc{th12\_M(oscillation, E, ne, ..., th13m=None)} implement the four
formulas above directly; \pyfunc{th12\_M} accepts a precomputed
\code{th13m} to avoid recomputing it when both angles are needed together
(as \pyfunc{Tei} below does).
\end{implbox}

\section{\texttt{landau\_zener.py} --- non-adiabatic corrections}
\label{sec:solar-lz}

\modulehead{medium/solar/landau\_zener.py}{Landau-Zener transition
probability $P_{LZ}(E)$ quantifying the breakdown of the adiabatic
approximation near the MSW resonance.}

\begin{theorybox}
At \emph{finite} adiabaticity, a neutrino crossing the MSW resonance can
jump between the two local matter eigenstates instead of adiabatically
following one of them. The exponential Landau-Zener approximation
\parencite{Parke1986} gives the jump probability as
\begin{equation}
  \keyresult{P_{LZ} = \exp\!\left(-\frac{\pi}{2}\gamma_{\rm res}\right)},
\end{equation}
with the adiabaticity parameter evaluated at the resonance
\parencite{GiuntiKim2007}
\begin{equation}
  \gamma_{\rm res} = \frac{\Delta m^2_{21}\sin^22\theta_{12}}{2E\hbar c\cos2\theta_{12}}\,L_n(r_{\rm res}),
  \qquad
  L_n = \left|\frac{n_e}{\mathrm{d}n_e/\mathrm{d}l}\right|_{\rm res},
\end{equation}
where $L_n$ is the electron-density scale height (in metres) at the
resonance radius $r_{\rm res}(E)$, defined by $V_k(\Delta m^2_{21},E,n_e(r_{\rm
res}))=\cos2\theta_{12}$ (i.e.\ where $\theta_{12}^M=\pi/4$). Large
$\gamma_{\rm res}$ means a slowly varying density on the scale of the
oscillation length at resonance, hence a vanishing jump probability
(adiabatic); small $\gamma_{\rm res}$ means a comparatively abrupt density
change, hence a sizeable jump probability.
\end{theorybox}

\begin{notebox}
\textbf{Scope and limitations, as documented in the module itself.} Only
the $\theta_{12}$-sector resonance is tracked ($\theta_{13}$'s resonance
sits deep in the solar core at very high density and is negligible for
every standard solar source). $P_{LZ}=0$ is returned wherever no resonance
exists in the solar volume: below-threshold energies, antineutrinos (no MSW
resonance for standard parameters), and LMA-Dark parameters
($\theta_{12}>\pi/4$, resonance at unphysical negative density). For
standard LMA parameters and pp/$^8$B solar neutrinos,
$\gamma_{\rm res}\sim10^3$--$10^4$, so $P_{LZ}\sim e^{-1500}\approx0$ and
the pure adiabatic approximation is excellent; this module becomes
relevant mainly for non-standard parameter explorations or precision
studies in the 1--3 MeV range. Solar NSI is \textbf{not} included here: the
resonance condition uses the standard $V_k$, and \pyfunc{Tei}
(\cref{sec:solar-probability}) silently ignores \code{p\_lz} once
\code{oscillation.nsi} is set, since both the resonance condition and the
adiabaticity parameter would need to be rederived for off-diagonal
$\epsilon$.
\end{notebox}

\begin{implbox}
\pyfunc{density\_gradient(solar\_profile)}: $\mathrm{d}n_e/\mathrm{d}\hat r$
on the tabulated grid via central differences (one-sided at the
boundaries). \pyfunc{resonance\_radius(oscillation, E, solar\_profile)}:
locates $r_{\rm res}(E)$ by linear interpolation at the first
sign-change of $V_k-\cos2\theta_{12}$ along the (monotonically decreasing)
density profile; returns \code{NaN} where no crossing exists.
\pyfunc{plz(oscillation, E, solar\_profile)}: the $P_{LZ}(E)$ formula
above, masked to exactly zero wherever \pyfunc{resonance\_radius} returns
\code{NaN}.
\end{implbox}

\begin{notebox}[title={No Coherent Medium Evolutor}]
\label{sec:solar-no-evolutor}
Unlike every other medium in this document, \texttt{medium.solar} has
\textbf{no \texttt{evolutor.py}}. This is a direct consequence of the
adiabatic approximation, not an omission: once a neutrino is assigned to a
local matter eigenstate at its production radius (\pyfunc{Tei}, next
section), the adiabatic theorem guarantees it stays in the
\emph{corresponding} eigenstate all the way to the solar surface, where the
density has dropped to (effectively) zero and the local matter eigenstates
have smoothly become the \emph{vacuum} mass eigenstates. No further
coherent phase evolution needs computing: the production-point projection
already \emph{is} the final incoherent mass-basis composition, and only the
final flavour projection (\pyfunc{solar\_probability\_state}) remains.
Consequently \texttt{medium.solar} exposes no
\pyfunc{solar\_probability\_transition}: unlike the other media
(\cref{ch:vacuum,ch:earth,ch:atmosphere}), there is no coherent evolutor $S$
to build a transition matrix from, so \pyfunc{solar\_probability\_state} is
the lowest-tier public probability function for this medium.
\end{notebox}

\section{\texttt{probability.py} --- production weights and flavour probabilities}
\label{sec:solar-probability}

\modulehead{medium/solar/probability.py}{\pyfunc{Tei} (production-point
matter-eigenstate weights), \pyfunc{solar\_\allowbreak{}probability\_\allowbreak{}mass} (integrated
over the production profile), \pyfunc{solar\_\allowbreak{}probability\_\allowbreak{}state} (final
flavour probabilities), and \pyfunc{solar\_\allowbreak{}probability\_\allowbreak{}integrated}
(spectrum-weighted energy average).}

\begin{theorybox}
\subsection*{\pyfunc{Tei}: flavour-to-matter-basis production weights}
$T_{ei}$ answers: given an electron neutrino produced at electron density
$n_e$ and energy $E$, what is the probability it is projected onto local
matter eigenstate $i$? In the Standard Model (adiabatic, no LZ) case this
is read directly off the matter-mixing angles of
\cref{sec:solar-matter-mixing}:
\begin{equation}
  \keyresult{
  \bigl(w_1,w_2,w_3\bigr) =
  \Bigl(\cos^2\theta_{13}^M\cos^2\theta_{12}^M,\ \
        \cos^2\theta_{13}^M\sin^2\theta_{12}^M,\ \
        \sin^2\theta_{13}^M\Bigr)
  },
\end{equation}
i.e.\ $|\langle\nu_e|\nu_i^M\rangle|^2$ for the standard PMNS-like
parametrisation applied to the matter-modified angles. When a
Landau-Zener probability $P_{LZ}$ is supplied (\cref{sec:solar-lz}), the
$1$--$2$ sector weights mix according to the LZ jump probability,
\begin{equation}
  w_1 = (1-P_{LZ})\cos^2\theta_{12}^M + P_{LZ}\sin^2\theta_{12}^M,
  \qquad
  w_2 = (1-P_{LZ})\sin^2\theta_{12}^M + P_{LZ}\cos^2\theta_{12}^M,
\end{equation}
while $w_3=\sin^2\theta_{13}^M$ is untouched (the $\theta_{13}$ resonance
is negligible, \cref{sec:solar-lz}). \textbf{When NSI is active}
(\code{oscillation.nsi} set), this entire analytic construction is
replaced by direct numerical diagonalisation: the full flavour-basis
Hamiltonian $H=H_{\rm kin}+H_{\rm mat}^{\rm NSI}$
(\pyfunc{hamiltonian\_flavour}, \cref{ch:core-common,ch:core-bsm}) is built
at each $(E,n_e)$ point and diagonalised with \pyfunc{torch.linalg.eigh};
the production weights are then simply the squared electron-flavour
components of each matter eigenvector,
\begin{equation}
  w_i(E,n_e) = \bigl|\langle\nu_e|\nu_i^M(E,n_e)\rangle\bigr|^2 = |{\rm eigvec}[\dots,0,i]|^2,
\end{equation}
which supports arbitrary off-diagonal $\epsilon$ (including LMA-Dark-type
degenerate scenarios) at the cost of losing the closed form; \code{p\_lz}
is silently ignored in this path.

\subsection*{From production weights to a final flavour probability}
\pyfunc{solar\_probability\_mass} integrates $T_{ei}(E,n_e(r))$ over the
production-radius distribution of one or more sources, weighted by each
source's normalised production-fraction profile $f_{\rm source}(r)$
(\cref{sec:solar-profile}):
\begin{equation}
  \keyresult{P_i(E) = \frac{\int_0^1 T_{ei}\bigl(E,n_e(r)\bigr)\,f_{\rm source}(r)\,\mathrm{d}r}
                          {\int_0^1 f_{\rm source}(r)\,\mathrm{d}r}},
\end{equation}
evaluated by the trapezoidal rule, giving an energy-dependent (not
radius-resolved) incoherent \emph{vacuum} mass-basis probability vector
$P_i(E)$ for that source -- by the argument of the note above, this
average over production radius already \emph{is} the mass-basis
composition reaching the solar surface, with no further coherent step
required. \pyfunc{solar\_probability\_state} performs the last, purely
kinematic conversion to flavour, using the ordinary incoherent-mixture
formula of
\cref{ch:core-common} (\texttt{probability.py} section) with the
\emph{identity} evolutor (there is nothing left to evolve coherently):
\begin{equation}
  \keyresult{P_\alpha(E) = \sum_i |U_{\alpha i}|^2\,P_i(E)},
\end{equation}
where $U$ is the ordinary \emph{vacuum} PMNS matrix
(\cref{ch:core-common,ch:core-sm}) -- the matter-modified angles only ever
appear inside $T_{ei}$, at the production point; every subsequent step in
this chapter works with the vacuum mixing matrix.
\end{theorybox}

\begin{implbox}
\begin{itemize}
  \item \pyfunc{Tei(oscillation, E, ne, p\_lz=None, legacy\_precision=False)}:
    the weight formulas above.
  \item \pyfunc{solar\_probability\_mass(oscillation, E, profile, sources, ...)}:
    builds the spatially resolved LZ mask (production radii above the
    resonance density only, \cref{sec:solar-lz}) when
    \code{profile.use\_LZ} and no NSI is active, then integrates $T_{ei}$
    against the source's production-fraction profile as above.
  \item \pyfunc{solar\_probability\_state(oscillation, E, profile, sources, ...)}:
    the final flavour-probability formula above, via
    \pyfunc{core.common.probability.probability\_incoherent}
    (\cref{ch:core-common}) with the identity operator standing in for
    ``no further coherent evolution''.
  \item \pyfunc{solar\_probability\_integrated(oscillation, E, profile, sources, spectrum, ...)}:
    builds \pyfunc{solar\_probability\_state} on an energy grid and
    averages it with
    \pyfunc{core.common.probability.probability\_integrated}
    (\cref{ch:core-common}), weighted by an explicit production
    \code{spectrum}.
\end{itemize}
\end{implbox}

\section{\texttt{flux.py} --- solar flux}

\modulehead{medium/solar/flux.py}{\pyfunc{solar\_\allowbreak{}flux\_\allowbreak{}state} and
\pyfunc{solar\_\allowbreak{}flux\_\allowbreak{}integrated}: apply per-source flux normalisations,
optional spectra, and energy integration to
\pyfunc{solar\_\allowbreak{}probability\_\allowbreak{}state}'s output.}

\begin{implbox}
\pyfunc{solar\_flux\_state(sources, profile, oscillation, E\_MeV, source\_spectrum=None, ...)}
adds no new physics: it calls \pyfunc{solar\_probability\_state} above for
the requested source(s), then \pyfunc{core.common.flux.flux\_state}
(\cref{ch:core-common}) to apply each source's total flux normalisation
(\pyfunc{SolarProfile.flux}) and optional energy spectrum -- the same
minimal pattern established in \cref{ch:vacuum}.
\pyfunc{solar\_flux\_integrated(...)} builds \pyfunc{solar\_flux\_state} on
an energy grid and integrates it with
\pyfunc{core.common.flux.flux\_integrated}, returning an unnormalised
physical rate.
\end{implbox}

\chapter{\texttt{medium.earth} --- Earth-crossing matter propagation}
\label{ch:earth}

Earth-crossing neutrinos (atmospheric, astrophysical, or terrestrial-source)
see a matter density that varies by four orders of magnitude between the
crust and the inner core, along a chord whose length and depth of
penetration depend entirely on the detector's nadir angle. This chapter
combines the perturbative machinery of \cref{ch:core-perturbative} (dense,
shell-crossing trajectories) with the exact numerical evolutor of
\cref{ch:core-numerical} (used as the reference and as an alternative
pipeline), plus a dedicated geometric layer converting detector angles into
trajectory coordinates, and a time-exposure layer averaging over a year of
Earth rotation and orbital motion.

\section{\texttt{profile.py} --- the Earth density model}
\label{sec:earth-profile}

\modulehead{medium/earth/profile.py}{\pyclass{EarthProfile}: shell radii and
a delegated perturbative density model, transformed into trajectory
coordinates on demand.}

\begin{implbox}
\textbf{Sources for the density profile.} Every option ultimately traces
back to the Preliminary Reference Earth Model, PREM
\parencite{DziewonskiAnderson1981} -- a seismologically constrained,
spherically symmetric radial density model, tabulated in $\sim$500 layers
from the inner core to the surface. \texttt{tpeanuts} offers three ways to
consume it, selected by \code{EarthParameters.profile\_perturbative\_name}
(\cref{ch:core-perturbative}, \texttt{model\_selection.py}):
\begin{itemize}
  \item \code{"even\_power"} (\textbf{the default}): the legacy
    \texttt{peanuts} even-power polynomial fit per shell,
    $n_e(r)=\alpha+\beta r^2+\gamma r^4+\dots$, read from
    \code{data/density/earth\_density.csv}. This is what every perturbative
    Earth evolutor formula in this chapter uses unless overridden.
  \item \code{"prem"}: the PREM500 table's own native piecewise-linear-in-$r^2$
    parametrisation per shell, consumed directly (converting tabulated mass
    density to electron density via layer-appropriate $\langle Z/A\rangle$
    ratios) rather than through a polynomial fit
    (\cref{ch:core-perturbative}).
  \item \textbf{Tabulated} (\code{earth\_tabulated\_density=True}): bypasses
    the perturbative polynomial representation entirely and interpolates the
    raw tabulated density directly -- at the cost of losing the closed-form
    residual integral that the perturbative pipeline needs, so this mode is
    only usable together with the numerical evolutor
    (\cref{ch:core-numerical}), not the analytical one below.
\end{itemize}
\end{implbox}

\begin{implbox}
\begin{itemize}
  \item \pyclass{EarthParameters}: bundles the profile-model selection
    above, \code{profile\_scale\_m}/\code{evolution\_scale\_m} (normally both
    $R_\oplus$, kept independent for testing), detector \code{depth\_m}, the
    \code{method} choice (\code{"analytical"}/\code{"numerical"},
    \cref{sec:earth-probability}), and the numerical-mode step count
    \code{nsteps} and sampling rule \code{ode\_method}
    (\cref{ch:core-numerical}).
  \item \pyfunc{EarthProfile.shells\_x(eta)}: for a surface-equivalent
    nadir angle, returns each shell's chord-coordinate crossing point via
    $x_j=\sqrt{r_j^2-\sin^2\eta}$ (clamped at 0), and which shells are
    physically crossed ($r_j>\sin\eta$).
  \item \pyfunc{EarthProfile.trajectory\_profile(eta)}: re-expresses the
    selected radial model in the trajectory coordinate $x$ via the chord
    identity $r^2=x^2+\sin^2\eta$ (\cref{sec:earth-geometry}).
  \item \pyfunc{EarthProfile.density\_x\_eta(x, eta)}: pointwise electron
    density at a trajectory coordinate, used directly by the numerical
    pipeline (\cref{sec:earth-probability}).
\end{itemize}
\end{implbox}

\begin{implbox}
\textbf{Optional neutron density, for the 3+1 sterile neutral-current
term.} Both the \code{"even\_power"} and \code{"prem"} layered profile
models optionally expose a neutron-density companion table alongside their
default electron-density one, built with \code{include\_neutron=True}:
\begin{itemize}
  \item \textbf{\code{"prem"}}: the PREM500 loader already computes a
    layer-appropriate $\langle Z/A\rangle$ to convert tabulated mass density
    to $n_e$ (implbox above); the complementary neutron fraction
    $1-\langle Z/A\rangle$ gives $n_n$ from the very same tabulated mass
    density, at no extra physical assumption --
    \pyfunc{load\_prem500\_neutron\_profile} fits it with the same
    piecewise-linear-in-$r^2$ shells (identical \code{rj} boundaries) as the
    electron-density loader.
  \item \textbf{\code{"even\_power"}}: has no composition information of its
    own (it is a bare polynomial fit to tabulated $n_e(r)$, next section),
    so its neutron companion is instead read from a separate, independently
    fitted table, \code{data/density/earth\_density\_nn.csv}, matching the
    default electron table's shell boundaries exactly.
\end{itemize}
\pyfunc{EarthProfile.density\_n\_x\_eta(x, eta)} / \pyfunc{call\_neutron(r,
eta)} evaluate this companion table pointwise, mirroring
\pyfunc{density\_x\_eta}/\pyfunc{call} above; they raise a clear
\code{ValueError} if the selected profile was not built with
\code{include\_neutron=True} (e.g.\ the default \code{"even\_power"}
construction). See \cref{ch:core-common} (\texttt{potential.py} section) for
the physics $n_n$ enables, and the \texttt{evolutor.py}/\texttt{probability.py}
sections below for how it is threaded through the analytical and numerical
Earth pipelines via \code{include\_matter\_nc}.
\end{implbox}

\section{\texttt{geometry.py} --- nadir angle and trajectory coordinates}
\label{sec:earth-geometry}

\modulehead{medium/earth/geometry.py}{Pure geometric conversions: detector
depth and nadir/zenith angle to dimensionless trajectory coordinates.
Computes no Hamiltonians, evolutors, or probabilities.}

\begin{notebox}[title={Nadir Angle and Its Relation to the Zenith Angle}]
Two angle conventions coexist in this chapter, both measured at the
detector. The \textbf{zenith angle} $\theta$ is measured from the local
upward vertical: $\theta=0$ is a neutrino arriving from directly overhead
(shortest possible path, no Earth crossing), $\theta=180^\circ$ is a
neutrino arriving from directly below (having crossed the full Earth
diameter). The \textbf{nadir angle} $\eta$ used internally throughout this
chapter is measured from the opposite direction (straight down, toward the
Earth's centre), so the two are exact complements:
\begin{equation}
  \keyresult{\eta = \pi - \theta}.
\end{equation}
$\eta=0$ is therefore the longest possible path (straight through the
centre) and $\eta=\pi/2$ is the horizon (grazing incidence, zero Earth
material for a surface detector). \pyfunc{earth\_evolutor\_from\_zenith}
(\cref{sec:earth-evolutor}) applies exactly this conversion for callers
who prefer to work with $\theta$ in degrees.
\end{notebox}

\begin{theorybox}
For a detector displaced below the surface by depth $d$, define the
dimensionless \textbf{detector radius}
\begin{equation}
  r_d = 1 - \frac{d}{R_\oplus}\ \ (\le 1),
\end{equation}
and the \textbf{surface-equivalent nadir angle} $\eta'$ -- the nadir angle
at which a ray from the Earth's centre through the detector would exit the
(undisplaced) surface, obtained from the sine rule applied to the chord
geometry:
\begin{equation}
  \keyresult{\eta' = \arcsin\bigl(r_d\sin\eta\bigr)}.
\end{equation}
$\eta'$ is what \pyfunc{EarthProfile.shells\_x}/\pyfunc{trajectory\_profile}
(\cref{sec:earth-profile}) actually consume, so that a buried detector sees
exactly the shells a surface detector at the corresponding $\eta'$ would.
The detector's own position along the trajectory chord is
\begin{equation}
  x_d = r_d\cos\eta,
\end{equation}
and, for shallow trajectories that never reach a tabulated shell boundary
(Case B, \cref{sec:earth-evolutor}), the total matter path length from the
surface to the detector is
\begin{equation}
  \keyresult{\Delta x_B = r_d\cos\eta + \sqrt{1-r_d^2\sin^2\eta}}.
\end{equation}
\pyfunc{classify\_eta\_regions} partitions any $\eta\in[0,\pi]$ into three
disjoint outcomes: \textbf{above horizon} ($\eta\in[\pi/2,\pi]$ for a
surface detector only -- no Earth material at all, only relevant when
$d=0$), \textbf{Case A} ($\eta\in[0,\pi/2)$, shell-crossing), and
\textbf{Case B} ($\eta\in[\pi/2,\pi]$ for $d>0$, shallow detector-near
paths) -- see \cref{sec:earth-evolutor}.
\end{theorybox}

\section{\texttt{evolutor.py} --- Case A and Case B evolutors}
\label{sec:earth-evolutor}

\modulehead{medium/earth/evolutor.py}{\pyfunc{earth\_evolutor}: the full
flavour-basis Earth evolution operator, dispatching every trajectory to one
of two geometrically distinct constructions.}

\begin{notebox}[title={Case A versus Case B}]
\begin{center}
\small
\begin{tabular}{@{}lll@{}}
\toprule
 & \textbf{Case A} ($0\le\eta<\pi/2$) & \textbf{Case B} ($\eta\ge\pi/2$, $d>0$) \\
\midrule
Path & \parbox[t]{4.6cm}{Deep, crosses one or more tabulated Earth shells} & \parbox[t]{4.6cm}{Shallow, surface to a buried detector, no shell crossed} \\[16pt]
Density & \parbox[t]{4.6cm}{Piecewise, one perturbative segment per crossed shell} & \parbox[t]{4.6cm}{Single constant value, sampled at the mid-depth radius} \\[16pt]
Segments & \parbox[t]{4.6cm}{Many (near half + mirrored far half)} & \parbox[t]{4.6cm}{One} \\[10pt]
Evolutor & \parbox[t]{4.6cm}{Composed product of many perturbative segment evolutors (\cref{ch:core-perturbative})} & \parbox[t]{4.6cm}{A single perturbative segment evolutor} \\[16pt]
\bottomrule
\end{tabular}
\end{center}
\end{notebox}

\begin{theorybox}
\subsection*{Case A: shell-crossing trajectories}
The chord identity $r^2=x^2+\sin^2\eta$ makes the physical density profile
\emph{even} in $x$: $n_e(x)=n_e(-x)$, since both $\pm x$ map to the same
$r$. This lets the chord be built from two mirror-symmetric halves without
recomputing the shell geometry twice:
\begin{enumerate}
  \item \textbf{Near half} (point of closest approach, $x=0$, to the
    detector, $x=x_d$): the crossed shells between $0$ and $x_d$ are
    ordered outermost-first; each is turned into a perturbative segment
    evolutor (\cref{ch:core-perturbative}) via the selected density model
    (\cref{sec:earth-profile}), and non-crossed shells are masked to the
    identity. These compose (\cref{ch:core-common}) into the detector-side
    half-evolutor $U_{\rm half,det}$.
  \item \textbf{Far half} (mirrored, $x\in[-x_d^{\max},0]$): recomputed
    \emph{directly} from the mirrored segment endpoints
    ($x_1,x_2)\to(-x_2,-x_1)$, \textbf{not} derived from the near half by
    transpose or reordering. Although the density profile is even, the
    first-order perturbative correction's oscillatory residual integral
    (\cref{ch:core-perturbative}) depends on the segment endpoints through a
    complex phase that is \emph{not} itself even, so the mirrored segment's
    evolutor genuinely differs and must be built through the same
    segment-evolutor pipeline independently. A transpose-based shortcut
    would only be valid when the reduced Hamiltonian is real (i.e.\
    $\delta_{14}=0$ for the 3+1 sterile extension, \cref{ch:core-bsm}); this
    direct recomputation is correct unconditionally, at the cost of one
    extra pass.
\end{enumerate}
The two halves combine into the reduced-basis Earth evolutor,
\begin{equation}
  \keyresult{U_{\rm red}^{\rm Earth} = U_{\rm half,det}\;U_{\rm half,far}},
\end{equation}
which is then dressed into the flavour basis via
\pyfunc{pmns.flavour\_basis} (\cref{ch:core-common,ch:core-sm,ch:core-bsm}),
correct for the Standard Model, NSI, the 3+1 sterile extension, or any
combination.

\subsection*{Case B: shallow detector-near trajectories}
For $\eta\ge\pi/2$ with a buried detector, the path from the surface to the
detector never reaches a tabulated shell boundary. The density is
approximated by its value at the mid-depth radius $r_{\rm mid}=1-d/(2R_\oplus)$,
treated as a single constant-density perturbative segment of length
$\Delta x_B$ (\cref{sec:earth-geometry}):
\begin{equation}
  \keyresult{U_{\rm red}^{\rm Earth} = \text{\pyfunc{evolutor\_perturbative\_segment}}\bigl(n_e(r_{\rm mid}),\ \Delta x_B\bigr)},
\end{equation}
again dressed to the flavour basis exactly as in Case A. Since the profile
is exactly constant on this single segment, the perturbative first-order
correction (\cref{ch:core-perturbative}) is identically zero and $U_0$
alone is exact for this segment.

Entries with $\eta\in[\pi/2,\pi]$ \emph{and} $d=0$ (above the horizon for a
surface detector) never enter either construction: they are left as the
identity operator directly by \pyfunc{earth\_evolutor}, since no Earth
material is crossed at all.
\end{theorybox}

\begin{implbox}
\pyfunc{earth\_evolutor(profile\_earth, oscillation, E, eta, depth\_m, ...)}:
classifies every $(E,\eta)$ entry (\pyfunc{classify\_eta\_regions}), routes
Case A entries through \pyfunc{\_earth\_evolutor\_case\_a\_batched} and
Case B entries through \pyfunc{\_earth\_evolutor\_case\_b\_batched}, and
optionally re-projects each result onto the nearest unitary matrix
(\code{reunitarize}, via \pyfunc{util.math.project\_to\_unitary}) to absorb
the small numerical drift the perturbative expansion can introduce.
\pyfunc{earth\_evolutor\_from\_zenith} is the zenith-angle convenience
wrapper (\cref{sec:earth-geometry}). Both support the 3-flavour Standard
Model and the 3+1 sterile extension transparently, sizing the identity
operator and every intermediate tensor by \code{oscillation.pmns.n\_flavours}.
Both also accept \code{include\_matter\_nc=False}, forwarded down to every
\pyfunc{evolutor\_perturbative\_segment} call in both Case A and Case B
(\cref{ch:core-perturbative}): when \code{True}, the Case B constant segment
additionally samples \pyfunc{EarthProfile.call\_neutron} at the mid-depth
radius (mirroring the electron-density sample above it), enabling the 3+1
sterile neutral-current term end-to-end. \code{False} (the default)
reproduces the pre-existing CC-only behaviour exactly, and requires no
neutron-density support from the selected profile model.
\end{implbox}

\section{\texttt{probability.py} --- analytical versus numerical}
\label{sec:earth-probability}

{\sloppy\modulehead{medium/earth/probability.py}{\pyfunc{earth\_\allowbreak{}probability\_\allowbreak{}state}:
dispatches between the perturbative-analytical and the exact-numerical
Earth probability pipeline. \pyfunc{earth\_\allowbreak{}probability\_\allowbreak{}transition} and
\pyfunc{earth\_\allowbreak{}probability\_\allowbreak{}integrated} build on it (full transition matrix
and spectrum-weighted energy average, respectively).}}

\begin{notebox}[title={Analytical versus Numerical}]
\begin{center}
\small
\begin{tabular}{@{}lll@{}}
\toprule
 & \textbf{\code{method="analytical"}} & \textbf{\code{method="numerical"}} \\
\midrule
Evolutor &
  \parbox[t]{4.4cm}{\pyfunc{earth\_evolutor} (\cref{sec:earth-evolutor}): a handful of perturbative segments per shell, closed-form residual integrals} &
  \parbox[t]{4.4cm}{\pyfunc{evolutor\_numerical} (\cref{ch:core-numerical}): \code{nsteps} matrix exponentials sampled along a fine \pyclass{Trajectory}} \\[36pt]
Accuracy driver &
  \parbox[t]{4.4cm}{First-order truncation per segment (\cref{ch:core-perturbative}); exact only if the density is well fit by the chosen model} &
  \parbox[t]{4.4cm}{Exact per segment; accuracy set purely by \code{nsteps} (\cref{ch:core-numerical})} \\[36pt]
Density model &
  \parbox[t]{4.4cm}{Perturbative even-power/PREM shell coefficients (\cref{sec:earth-profile})} &
  \parbox[t]{4.4cm}{Any \pyfunc{EarthProfile.density\_x\_eta}, including the tabulated (interpolated) option} \\[26pt]
Batching &
  \parbox[t]{4.4cm}{Fully vectorised over $(E,\eta)$} &
  \parbox[t]{4.4cm}{Currently scalar $\eta$ per call} \\[10pt]
Extras &
  \parbox[t]{4.4cm}{\code{reunitarize} drift correction} &
  \parbox[t]{4.4cm}{\code{full\_oscillation} returns the whole path, not just the endpoint} \\[20pt]
\bottomrule
\end{tabular}
\end{center}
The numerical pipeline is the reference: it makes no perturbative
assumption, so it is what the analytical pipeline is validated against
(matching the general relationship between \cref{ch:core-numerical,ch:core-perturbative}).
\end{notebox}

\begin{implbox}
\begin{itemize}
  \item \pyfunc{earth\_probability\_transition(profile\_earth, oscillation, E\_MeV, eta, depth\_m, ...)}:
    builds $U_{\rm Earth}$ via \pyfunc{earth\_evolutor} and returns the full
    $|U_{\rm Earth}[\beta,\alpha]|^2$ matrix via
    \pyfunc{core.common.probability.probability\_transition}
    (\cref{ch:core-common}); uses the analytical evolutor only, since the
    numerical segment pipeline below only ever evolves an explicit initial
    state and has no matrix-level counterpart.
  \item \pyfunc{earth\_probability\_state\_analytical(nustate, profile\_earth, oscillation, E\_MeV, eta, depth\_m, massbasis, ...)}:
    builds $U_{\rm Earth}$ via \pyfunc{earth\_evolutor} and applies the
    coherent or incoherent probability formula of \cref{ch:core-common}
    depending on \code{massbasis}.
  \item \pyfunc{earth\_probability\_state\_numerical(nustate, profile\_earth, oscillation, E\_MeV, eta, depth\_m, ...)}:
    builds a scalar \pyclass{Trajectory} (\pyfunc{build\_earth\_trajectory},
    \cref{ch:core-numerical}) -- Case A samples
    \pyfunc{EarthProfile.density\_x\_eta} along the chord, Case B repeats
    the constant mid-depth density -- then calls \pyfunc{evolutor\_numerical}
    and the same probability dispatch; \code{full\_oscillation=True}
    additionally returns the probability at every sampled point, not only
    the final one.
  \item \pyfunc{earth\_probability\_state(...)}: the public dispatcher
    selecting one of the two pipelines by the string \code{method}.
  \item \pyfunc{earth\_probability\_integrated(nustate, profile\_earth, oscillation, E\_MeV, eta, depth\_m, spectrum, ...)}:
    builds \pyfunc{earth\_probability\_state} on an energy grid at a fixed
    nadir angle and averages it with
    \pyfunc{core.common.probability.probability\_integrated}
    (\cref{ch:core-common}), weighted by an explicit production
    \code{spectrum}. This time-averages over energy, not over nadir angle,
    and is therefore distinct from \pyfunc{earth\_probability\_exposure}
    below (\cref{sec:earth-exposure}).
\end{itemize}
\end{implbox}

\begin{implbox}
Every function above (and \pyfunc{earth\_probability\_transition}) accepts
\code{include\_matter\_nc=False}, forwarded to \pyfunc{earth\_evolutor} for
the analytical pipeline or sampled directly from
\pyfunc{profile\_earth.call\_neutron} alongside the electron-density sample
for the numerical one (both Case A and the Case B constant-density
fallback). Requesting it with a profile model that lacks neutron-density
coefficients (\cref{sec:earth-profile}) raises a clear \code{ValueError}
rather than silently falling back to CC-only; requesting it for a
3-flavour \code{oscillation.pmns} is silently ignored (\cref{ch:core-common}).
\end{implbox}

\section{\texttt{flux.py} --- Earth flux}

{\sloppy\modulehead{medium/earth/flux.py}{\pyfunc{earth\_\allowbreak{}flux\_\allowbreak{}state},
\pyfunc{earth\_\allowbreak{}flux\_\allowbreak{}integrated}, and \pyfunc{earth\_\allowbreak{}flux\_\allowbreak{}exposure}: apply
\texttt{core.common.flux} normalisations, energy integration, and
exposure integration (respectively) on top of
\pyfunc{earth\_\allowbreak{}probability\_\allowbreak{}state} and \pyfunc{earth\_\allowbreak{}probability\_\allowbreak{}exposure}
(\cref{sec:earth-exposure}).}}

\begin{implbox}
\pyfunc{earth\_flux\_state(nustate, profile\_earth, oscillation, E\_MeV, eta, depth\_m, flux, spectrum, method, ...)}
calls \pyfunc{earth\_probability\_state} then
\pyfunc{core.common.flux.flux\_state} (\cref{ch:core-common}), propagating
the \code{(probabilities\_along\_path, x)} pair through unchanged when
\code{method="numerical"} and \code{full\_oscillation=True}.
\pyfunc{earth\_flux\_integrated(nustate, profile\_earth, oscillation, E\_MeV,
eta, depth\_m, flux, spectrum, ...)} builds \pyfunc{earth\_flux\_state} on an
energy grid at a fixed nadir angle and integrates it with
\pyfunc{core.common.flux.flux\_integrated}, returning an unnormalised
physical rate. \pyfunc{earth\_flux\_exposure(nustate, profile\_earth,
oscillation, E\_MeV, depth\_m, flux, spectrum, exposure, ...)} instead starts
from \pyfunc{earth\_probability\_exposure} (\cref{sec:earth-exposure}): the
angular exposure average happens at the probability level, before the flux
normalisation is applied. All three accept \code{include\_matter\_nc=False},
forwarded unchanged to the underlying probability function.
\end{implbox}

\section[Nadir-angle exposure]{Nadir-angle exposure: \texttt{exposure\_math.py}, \texttt{exposure\_table.py}, \texttt{exposure\_integration.py}, \texttt{exposure\_io.py}}
\label{sec:earth-exposure}

\modulehead{medium/earth/exposure\_*.py}{Together, these four modules build
a time-averaged nadir-angle exposure weight $W(\eta)$ for a detector at
fixed geographic latitude, and integrate Earth probabilities against it.}

\begin{theorybox}
\subsection*{The physical exposure problem}
A detector at fixed geographic latitude $\lambda$ observes a given nadir
angle $\eta$ only during specific windows of each day, because the visible
nadir angle towards a fixed astrophysical/atmospheric source direction
sweeps as the Earth rotates and as the Sun's apparent declination
$\delta_\odot(t)$ (approximated as varying sinusoidally with the orbital
inclination $\iota\approx23.44^\circ$ over the year) changes with the
season. The exposure weight
\begin{equation}
  \keyresult{W(\eta) = \int_{d_1}^{d_2}\!\!\bigl[\text{fraction of day } d \text{ spent at nadir angle }\eta\bigr]\,\mathrm{d}d}
\end{equation}
integrates that daily observability window over a requested day-of-year
range $[d_1,d_2]\subseteq[0,365]$ (the full year by default), giving a
weight function used to time-average Earth-crossing probabilities:
\begin{equation}
  \keyresult{P_{\rm exp}(E) = \int_0^\pi W(\eta)\,P_{\rm Earth}(E,\eta)\,\mathrm{d}\eta}.
\end{equation}

\subsection*{From an elementary rate to $W(\eta)$: the elliptic-integral machinery}
The instantaneous rate at which nadir angle $\eta$ is observed has no
elementary closed-form antiderivative in the day-phase variable; the
project evaluates it via \textbf{incomplete elliptic integrals of the
first kind}, $F(\varphi,m)$ (\pyfunc{util.math.ellipf\_incomplete}), through
three nested closed-form functions, entirely in complex arithmetic to
remain well-defined across the geometry's branch cuts (the physical result
is always the real part of a difference):
\begin{itemize}
  \item \pyfunc{IndefiniteIntegralDay(T, eta, lam)}: the antiderivative,
    in the day-phase variable $T\in[-1,1]$, of the instantaneous
    nadir-observation rate, expressed through $F(\varphi,m)$ with an
    elliptic amplitude $\varphi$ and parameter $m$ built from
    $\sin\iota$, $\sin(\eta\mp\lambda)$, and $T$ (the exact closed form is
    a several-term rational combination of square roots of these
    quantities, verified numerically against direct integration in the
    test suite rather than reproduced symbol-by-symbol here).
  \item \pyfunc{IntegralAngle(eta, lam, a1, a2)}: evaluates
    \pyfunc{IndefiniteIntegralDay} at the two endpoints of the
    sub-interval of $T$ consistent with \emph{both} the day/night geometry
    at nadir angle $\eta$ \emph{and} a single year-angle window
    $[a_1,a_2]$ (found via \pyfunc{util.math.intersection} of the relevant
    validity intervals), returning zero when no valid sub-interval exists.
  \item \pyfunc{IntegralDay(eta, lam, d1, d2)}: converts the requested
    day-of-year window to year-angle radians ($a=2\pi d/365$), splits it at
    the half-year boundaries (the day/night geometry is evaluated
    separately on each half via a sign reflection), and sums the
    \pyfunc{IntegralAngle} contributions -- this \emph{is} $W(\eta)$ for one
    nadir angle.
\end{itemize}
\pyfunc{make\_eta\_grid(ns, daynight=...)} builds the uniform grid on which
$W(\eta)$ is tabulated, optionally restricted to the day half (larger
$\eta$, closer to $\pi$) or night half (smaller $\eta$, closer to $0$,
where core-crossing trajectories concentrate) of the full $[0,\pi]$ range.
\end{theorybox}

\begin{implbox}
\begin{itemize}
  \item \pyclass{ExposureParameters}: selects the exposure backend
    (\code{exposure\_source}: \code{"math"} direct integration,
    \code{"cache"} a precomputed torch file, \code{"csv"} a tabulated file
    in any of the Nadir/Zenith/CosZenith conventions, or \code{"legacy"}
    delegating to legacy \texttt{peanuts} for validation), the detector
    latitude, the day-of-year window, the day/night filter, and grid size.
  \item \pyfunc{build\_nadir\_exposure(exposure, context, normalized)}
    (\texttt{exposure\_table.py}): dispatches to the selected backend and
    returns a \pyclass{NadirExposureTable} (an $(\eta, W)$ pair),
    optionally normalised to unit trapezoidal integral so it reads as a
    probability density over nadir angle.
  \item \pyfunc{integrate\_exposure(probabilities\_eta, eta, exposure)}:
    the $P_{\rm exp}(E)$ formula above, evaluated by the trapezoidal rule
    (\code{torch.trapezoid}).
  \item \pyfunc{earth\_probability\_exposure(...)} (\texttt{exposure\_integration.py}):
    the end-to-end entry point -- builds the exposure table, evaluates
    \pyfunc{earth\_probability\_state} (\cref{sec:earth-probability}) over
    the full $\eta$ grid (optionally chunked via \code{chunk\_eta} to bound
    memory), and integrates against $W(\eta)$. Accepts
    \code{include\_matter\_nc=False}, forwarded to every
    \pyfunc{earth\_probability\_state} call.
  \item \texttt{exposure\_io.py}: CSV and torch-cache read/write helpers
    for the \code{"csv"}/\code{"cache"} backends above -- pure I/O, no
    additional physics.
\end{itemize}
\end{implbox}

\chapter{\texttt{medium.atmosphere} --- atmospheric neutrino propagation}
\label{ch:atmosphere}

Atmospheric neutrinos are produced by cosmic-ray air showers at a range of
altitudes above the Earth's surface and propagate down through the
atmosphere -- a comparatively low-density medium, but one whose geometry
(production height, zenith angle, detector depth) is considerably richer
than the single-scale exponential often assumed. This chapter follows the
same profile/geometry/evolutor/probability/flux structure established in
\cref{ch:vacuum}, adapted to that geometry, plus a
production-altitude density layer (\texttt{density.py}) and an atmospheric
slant-depth layer (\texttt{depth.py}) with no direct analogue in the
previous media chapters.

\section{\texttt{profile.py} and \texttt{density.py} --- the atmosphere density model}
\label{sec:atmosphere-profile}

\modulehead{medium/atmosphere/density.py, profile.py}{Pluggable
production-to-surface electron-density backends, and
\pyclass{AtmosphereProfile}, the trajectory container consumed by the
numerical evolutor (\cref{ch:core-numerical}).}

\begin{implbox}
Unlike the Earth chapter's single PREM-derived source
(\cref{sec:earth-profile}), \pyfunc{atmosphere\_density} offers a genuine
choice of density backends, selected by \code{atmosphere\_density\_source}:
\begin{itemize}
  \item \code{"exponential"} (\textbf{the default}): the barometric formula
    $\rho(h)=\rho_0 e^{-h/H}$, with $\rho_0=1.225\times10^{-3}\ \mathrm{g/cm^3}$
    (U.S. Standard Atmosphere sea-level density) and scale height
    $H=8.4\ \mathrm{km}$ -- a single-parameter approximation, adequate when
    only the coarse exponential fall-off matters.
  \item \code{"nusquids"}/\code{"earthatm"}: the external nuSQuIDS
    \code{EarthAtm} density formula (\texttt{external.nusquids}), for
    bit-comparable cross-checks against that reference implementation.
  \item \code{"file"}: linear interpolation of a user-supplied two-column
    altitude/density table.
  \item \code{"mceq"}: densities consistent with an initialised MCEq
    atmosphere object \parencite{Fedynitch2015MCEq}, for internal
    consistency with an MCEq-generated production flux.
  \item \code{"msis"}/\code{"pymsis"}: the NRLMSIS 2.0 empirical
    atmosphere model \parencite{Emmert2021NRLMSIS2}, the most physically
    detailed option (season-, location-, and space-weather-dependent), via
    the optional \texttt{pymsis} backend.
\end{itemize}
Every backend returns mass density $\rho(h)$ in g/cm$^3$; electron density
is obtained via $n_e=Y_e\,\rho\times\text{(mol/g conversion)}$, with the
electron fraction $Y_e$ read from the backend's own configuration (PyMSIS)
or from the shared default $Y_e=0.5$ (air, elsewhere) -- exactly the same
mass-to-electron-density conversion pattern as
\cref{ch:core-common}/\cref{sec:earth-profile}. \pyfunc{atmosphere\_density}
also accepts \code{density\_type="neutron\_density"}, returning
$n_n=(1-Y_e)\,\rho\times\text{(mol/g conversion)}$ from the very same
$Y_e$/$\rho$ already computed for $n_e$ -- no new backend logic, since air's
composition is fully captured by the single scalar $Y_e$ already in use.
This enables the 3+1 sterile neutral-current term (\cref{ch:core-common})
for every density backend uniformly.
\end{implbox}

\begin{implbox}
\pyclass{AtmosphereParameters} bundles the density-backend selection above
with the numerical trajectory settings (\code{nsteps}, sampling
\code{method}) and the analytical-mode fit settings
(\code{perturbative\_segments}, \code{perturbative\_degree}, next section);
\code{matter=False} forces zero density everywhere (a pure-vacuum check).
\pyclass{AtmosphereProfile} builds the production-to-surface
\pyclass{Trajectory} (\cref{ch:core-numerical}) from
(\code{h\_km},\code{theta\_deg},\code{depth\_km}) via the geometry of the
next section, and samples the selected density backend once per segment.
\code{AtmosphereParameters.include\_matter\_nc=False} additionally makes
\pyclass{AtmosphereProfile} sample \code{density\_type="neutron\_density"}
alongside $n_e$, exposed as \code{n\_n\_molcm3} (\code{None} when not
requested); since \pyfunc{atmosphere\_density} supports every backend
uniformly (implbox above), unlike \cref{ch:earth} there is no profile model
for which this can fail to be available.
\end{implbox}

\section{\texttt{geometry.py} and \texttt{depth.py} --- angles, altitude, and atmospheric depth}
\label{sec:atmosphere-geometry}

\modulehead{medium/atmosphere/geometry.py, depth.py}{Zenith/nadir angle
conversions, surface/detector angle mapping, path-length geometry, and the
altitude $\leftrightarrow$ atmospheric-depth relation.}

\begin{notebox}[title={Two Angles and Two ``Depths''}]
This chapter uses \textbf{three} angle labels and \textbf{two} unrelated
quantities both called ``depth'' -- keeping them apart is essential:
\begin{itemize}
  \item $\theta$ (\textbf{detector zenith angle}, the MCEq convention):
    measured at the detector from the local upward vertical;
    $\theta=0^\circ$ is straight down (shortest path), $\theta=90^\circ$ is
    horizontal.
  \item $\eta$ (\textbf{nadir angle}, \cref{ch:earth}'s convention):
    $\eta=\pi-\theta$, identical relation and identical meaning to the one
    used throughout the Earth chapter.
  \item $\alpha$ (\textbf{surface zenith angle}): the zenith angle the same
    ray would have if measured by an observer \emph{at the Earth's
    surface} directly above/along the trajectory, rather than at the
    (possibly buried) detector -- related to $\theta$ by impact-parameter
    conservation (below), not simple equality.
  \item \code{depth\_km} or \code{depth\_m} (\textbf{detector depth}): how far
    below the Earth's surface the detector sits -- a \emph{geometric}
    length, identical in meaning to Earth's $d$ (\cref{ch:earth}).
  \item $X$ (\textbf{atmospheric depth}, \texttt{depth.py}): a
    \emph{column density} in g/cm$^2$ (mass per unit area integrated along
    the path), \textbf{not} a length -- the physical quantity most
    directly tied to how many interaction lengths of atmosphere a
    cosmic-ray shower (and hence its neutrino production point) has
    traversed. It is unrelated to detector depth beyond sharing the English
    word.
\end{itemize}
\end{notebox}

\begin{theorybox}
\subsection*{Angle conversions}
The nadir/zenith relation is exactly the one used throughout
\cref{ch:earth}:
\begin{equation}
  \keyresult{\eta = \pi - \theta}.
\end{equation}
The detector zenith angle $\theta$ and the surface zenith angle $\alpha$ of
the \emph{same} straight-line ray are related by conservation of the
impact parameter (the perpendicular distance from Earth's centre to the
ray is the same whether measured from the detector's radius $r_d=R_\oplus-d$
or the surface radius $R_\oplus$):
\begin{equation}
  \keyresult{r_d\sin\theta = R_\oplus\sin\alpha}
  \qquad\Longleftrightarrow\qquad
  \theta=\arcsin\!\Bigl(\frac{R_\oplus}{r_d}\sin\alpha\Bigr),\
  \alpha=\arcsin\!\Bigl(\frac{r_d}{R_\oplus}\sin\theta\Bigr).
\end{equation}
Because $R_\oplus/r_d>1$ for a buried detector, not every surface angle
maps to a valid detector ray: $\alpha$ must stay below
$\alpha_{\max}=\arcsin(r_d/R_\oplus)$ for the first arcsine to remain
defined (\pyfunc{alpha\_max\_for\_detector\_depth}). For a surface detector
($d=0$), $r_d=R_\oplus$ and $\alpha\equiv\theta$ trivially.

\subsection*{Path lengths}
For a neutrino produced at altitude $h$ and detected at zenith angle
$\theta$ with detector depth $d$ (so $r_d=R_\oplus-d$), the straight-line
law of cosines applied to the triangle (Earth centre, detector, production
point) gives the total detector-to-production distance,
\begin{equation}
  \keyresult{L_{\rm total} = -r_d\cos\theta+\sqrt{(R_\oplus+h)^2-r_d^2\sin^2\theta}},
\end{equation}
and, setting $h=0$, the purely \textbf{underground} part of the same ray
(detector to surface),
\begin{equation}
  L_{\rm underground} = -r_d\cos\theta+\sqrt{R_\oplus^2-r_d^2\sin^2\theta}.
\end{equation}
The \textbf{atmospheric} path length actually carrying matter effects is
the difference, $L_{\rm atm}=L_{\rm total}-L_{\rm underground}$
(\pyfunc{atmosphere\_path\_length}); \pyfunc{atmosphere\_path\_grid} samples
altitude along it via the same law-of-cosines relation evaluated at
distance $s$ from the detector,
\begin{equation}
  \keyresult{r(s)=\sqrt{r_d^2+s^2+2r_ds\cos\theta}},\qquad h(s)=r(s)-R_\oplus.
\end{equation}

\subsection*{Vertical and slant atmospheric depth}
Independently of the path-length geometry above, the atmospheric depth $X$
(the column density integral, used e.g.\ to tabulate cosmic-ray shower
development or MCEq/Honda fluxes rather than oscillation propagation
directly) is
\begin{equation}
  \keyresult{X_{\rm vertical}(h) = \int_h^\infty \rho(h')\,\mathrm{d}h'},
  \qquad
  X_{\rm slant}(h,\alpha) = \frac{X_{\rm vertical}(h)}{\cos\alpha},
\end{equation}
the second equation being the plane-parallel approximation (valid for
$\alpha$ not too close to the horizon, where the Earth's curvature would
need to be accounted for); $\alpha$ here is the \emph{surface} zenith
angle defined above, since the plane-parallel picture is anchored to a
flat patch of atmosphere.
\end{theorybox}

\begin{implbox}
\begin{itemize}
  \item \pyfunc{theta\_to\_eta} / \pyfunc{eta\_to\_theta}: $\eta=\pi-\theta$
    and its inverse.
  \item \pyfunc{alpha\_surface\_to\_theta\_detector} /
    \pyfunc{theta\_detector\_to\_alpha\_surface} /
    \pyfunc{alpha\_max\_for\_detector\_depth}: the impact-parameter mapping
    above and its domain limit.
  \item \pyfunc{total\_path\_length} / \pyfunc{underground\_path\_length} /
    \pyfunc{atmosphere\_path\_length}: the three path-length formulas
    above.
  \item \pyfunc{altitude\_along\_detector\_path} /
    \pyfunc{atmosphere\_path\_grid}: the $r(s)$ sampling formula above,
    and the resulting distance/altitude grid used to build an
    \pyclass{AtmosphereProfile} trajectory.
  \item \pyfunc{atmosphere\_vertical\_depth(h\_km, rho\_gcm3)} /
    \pyfunc{atmosphere\_slant\_depth(X\_vertical, alpha\_deg)}
    (\texttt{depth.py}): the $X_{\rm vertical}$/$X_{\rm slant}$ formulas
    above, the first via trapezoidal integration from the top of the
    supplied grid downward.
  \item \pyfunc{compute\_dXdh} / \pyfunc{interpolate\_flux\_at\_Xobs}
    (\texttt{depth.py}): the depth-profile derivative and a
    depth-tabulated-flux interpolator, both backend-neutral utilities for
    working with $\Phi(E,X,\alpha)$ tables from any external source.
\end{itemize}
\end{implbox}

\section{\texttt{evolutor.py} --- atmosphere evolution operators}
\label{sec:atmosphere-evolutor}

\modulehead{medium/atmosphere/evolutor.py}{\pyfunc{atmosphere\_evolutor}:
dispatches between a numerical (default) and an analytical perturbative
atmosphere evolutor.}

\begin{implbox}
The same numerical/analytical distinction of \cref{ch:earth}
applies here, with the \textbf{default reversed} (\code{method="numerical"}
by default, versus Earth's \code{"analytical"} default): the atmosphere's
low density and short path make the numerical matrix-exponential pipeline
(\cref{ch:core-numerical}) cheap enough to use by default.
\begin{itemize}
  \item \pyfunc{atmosphere\_evolutor\_numerical(...)}: builds an
    \pyclass{AtmosphereProfile} trajectory (\cref{sec:atmosphere-profile})
    and calls \pyfunc{evolutor\_numerical} (\cref{ch:core-numerical})
    directly on its sampled density; entries with $L_{\rm atm}\le0$
    (grazing/no-path geometries) are masked to the identity.
  \item \pyfunc{atmosphere\_evolutor\_analytical(...)}: \emph{fits} a local
    polynomial density model on the fly -- splits the atmospheric path
    into \code{perturbative\_segments} equal sub-intervals, samples the
    density backend at \code{perturbative\_degree}$+1$ Chebyshev-like nodes
    per sub-interval (\code{q}$\in[-1,1]$, centred/rescaled), and builds an
    \pyclass{AtmospherePolynomialProfile}
    (\cref{ch:core-perturbative}, \texttt{atmosphere} model) from those
    samples, then evolves it with
    \pyfunc{evolutor\_perturbative\_segment} (\cref{ch:core-perturbative})
    exactly as the Earth analytical pipeline does. Unlike Earth's
    tabulated shell coefficients, the polynomial here is fitted
    \emph{per call}, from whichever density backend is selected
    (\cref{sec:atmosphere-profile}) -- the price of supporting arbitrary
    density sources with the same closed-form perturbative machinery.
  \item \pyfunc{atmosphere\_evolutor(...)}: the public dispatcher, mirroring
    \pyfunc{earth\_probability\_state} (\cref{sec:earth-probability}) in
    spirit.
\end{itemize}
\end{implbox}

\begin{implbox}
Both evolutors accept \code{atmosphere.include\_matter\_nc=False}
(\cref{sec:atmosphere-profile}), enabling the 3+1 sterile neutral-current
term (\cref{ch:core-common}). The numerical evolutor forwards
\pyclass{AtmosphereProfile}'s sampled \code{n\_n\_molcm3} straight through
to \pyfunc{evolutor\_numerical} as \code{n\_n\_mol\_cm3}. The analytical
evolutor fits a \emph{second}, independent
\pyclass{AtmospherePolynomialProfile} to a neutron-density sample at the
same nodes as the electron-density fit, and passes its coefficients as
\code{coefficients\_n} into the same segment model
(\cref{ch:core-perturbative}), then calls
\pyfunc{evolutor\_perturbative\_segment} with
\code{include\_matter\_nc=True}. Both default to \code{False}, reproducing
the pre-existing CC-only behaviour exactly.
\end{implbox}

\section{\texttt{probability.py} --- atmosphere-surface probabilities}
\label{sec:atmosphere-probability}

{\sloppy\modulehead{medium/atmosphere/probability.py}{\pyfunc{atmosphere\_\allowbreak{}probability\_\allowbreak{}transition}
and \pyfunc{atmosphere\_\allowbreak{}probability\_\allowbreak{}state}: final flavour probabilities at
the Earth surface from an atmosphere evolution operator.
\pyfunc{atmosphere\_\allowbreak{}probability\_\allowbreak{}integrated},
\pyfunc{atmosphere\_\allowbreak{}probability\_\allowbreak{}integrated\_\allowbreak{}angular}, and
\pyfunc{atmosphere\_\allowbreak{}probability\_\allowbreak{}integrated\_\allowbreak{}height} build on
\pyfunc{atmosphere\_\allowbreak{}probability\_\allowbreak{}state} to average over energy, the full
zenith range, and production height, respectively.}}

\begin{implbox}
Structurally identical to \cref{ch:vacuum,ch:earth}:
\pyfunc{atmosphere\_\allowbreak{}probability\_\allowbreak{}transition(oscillation, E\_MeV, h\_km, theta\_deg, depth\_km, method, ...)}
returns the full transition-probability matrix $P=|S_{\rm atm}|^2$
(\cref{ch:core-common}); \pyfunc{atmosphere\_\allowbreak{}probability\_\allowbreak{}state(nustate, ..., massbasis, ...)}
dispatches between the coherent formula
($\psi_{\rm surface}=S_{\rm atm}\psi_{\rm initial}$,
$P_\alpha=|\psi_{\rm surface,\alpha}|^2$) and the incoherent formula
($P_\alpha=\sum_i|(S_{\rm atm}U_{\rm PMNS})_{\alpha i}|^2w_i$) of
\cref{ch:core-common}, exactly as \pyfunc{vacuum\_probability\_state} and
\pyfunc{earth\_probability\_state} do,
using whichever evolutor \pyfunc{atmosphere\_evolutor}
(\cref{sec:atmosphere-evolutor}) returns. Because every function in this
section takes the whole \pyclass{AtmosphereParameters} bundle as its
\code{atmosphere} argument rather than exploding it into individual
keywords, \code{include\_matter\_nc} (previous section) is already
available at every level here with no further plumbing.
\end{implbox}

\begin{notebox}
\textbf{A real, previously uncaught 3-flavour-only bug.}
\pyfunc{atmosphere\_probability\_state} used to crash outright for a
4-flavour (3+1 sterile) initial state, even though
\pyfunc{atmosphere\_evolutor} underneath it already built the correct
$4\times4$ operator: a shared broadcasting helper
(\code{util.type.broadcast\_last3}, used identically in
\cref{ch:vacuum}) hardcoded its input's final dimension to exactly 3. The
fix generalises the helper to accept 3 \emph{or} 4 (renamed
\code{broadcast\_flavour\_vector} to match, since it is no longer
3-specific) and is shared by \texttt{medium.vacuum.probability},
\texttt{medium.vacuum.evolutor}, and this module. No test exercised a
sterile initial state through this function before the fix.
\end{notebox}

\begin{implbox}
Atmospheric probabilities depend on three axes at once ($E$, production
height $h$, zenith angle $\theta$), so this medium is the only one with
three dedicated \code{*\_integrated} reductions, each delegating to the
corresponding \cref{ch:core-common} primitive:
\begin{itemize}
  \item \pyfunc{atmosphere\_probability\_integrated(nustate, ..., spectrum, integrate\_angular, integrate\_height, ...)}:
    \textbf{always} averages over energy (spectrum-weighted, via
    \pyfunc{core.common.probability.probability\_integrated}), then
    optionally chains a height average and/or an angular average by hand,
    each a further, independent reduction rather than one fused operation.
  \item \pyfunc{atmosphere\_probability\_integrated\_angular(nustate, ...)}:
    averages over the full zenith range via
    \pyfunc{core.common.probability.probability\_integrated\_angular}
    (geometric $\sin\theta$ weight).
  \item \pyfunc{atmosphere\_probability\_integrated\_height(nustate, ..., production\_flux, ...)}:
    averages over production height, weighted by the height-resolved
    production flux, again via
    \pyfunc{core.common.probability.probability\_integrated} -- height is
    not a universal \cref{ch:core-common} concept, so this wrapper lives
    here, even though it reuses that module's generic weighted-average
    primitive.
\end{itemize}
Every reduction's \code{*\_dim} parameter defaults to $-2$ (``the axis
immediately before flavour''), matching \cref{ch:core-common}'s own
convention; for the composite function's fixed energy-then-height-then-angle
chaining order to line up with that default at every step, the input grid
should be laid out as $(\ldots,\theta,h,E,\text{flavour})$.
\end{implbox}

\section{\texttt{flux.py} --- atmosphere flux}
\label{sec:atmosphere-flux}

{\sloppy\modulehead{medium/atmosphere/flux.py}{\pyfunc{atmosphere\_\allowbreak{}flux\_\allowbreak{}state}
(surface flux from \pyfunc{atmosphere\_\allowbreak{}probability\_\allowbreak{}state}) and three
\code{*\_integrated} variants mirroring the probability-layer reductions
above.}}

\begin{implbox}
Like \texttt{probability.py} above, every function here takes the whole
\code{atmosphere} (\pyclass{AtmosphereParameters}) bundle, so
\code{include\_matter\_nc} is available with no further plumbing.
\begin{itemize}
  \item \pyfunc{atmosphere\_flux\_state(nustate, oscillation, E\_MeV, h\_km, theta\_deg, flux, spectrum, depth\_km, ...)}:
    the standard pattern -- \pyfunc{atmosphere\_probability\_state} then
    \pyfunc{core.common.flux.flux\_state} (\cref{ch:core-common}).
  \item \pyfunc{atmosphere\_flux\_integrated(..., integrate\_angular, integrate\_height, ...)}:
    always integrates over energy via
    \pyfunc{core.common.flux.flux\_integrated} (an unnormalised rate,
    unlike the probability-layer average), then optionally chains a plain
    trapezoidal height integral and/or an angular integral
    (\pyfunc{core.common.flux.flux\_integrated\_angular}) -- no separate
    height weight is needed here, since \code{flux} already carries
    whatever height dependence it has.
  \item \pyfunc{atmosphere\_flux\_integrated\_angular(...)}: integrates over
    the full zenith range via
    \pyfunc{core.common.flux.flux\_integrated\_angular}.
  \item \pyfunc{atmosphere\_flux\_integrated\_height(..., production\_flux, ...)}:
    integrates over production height using the tabulated production-flux
    data. Atmosphere-specific and, unlike the three functions above, calls
    no \texttt{core.common} function directly: it combines
    \pyfunc{atmosphere\_probability\_integrated\_height} (the
    production-flux-weighted average probability over height, $\langle
    P\rangle_h$) with the plain trapezoidal integral of
    \code{production\_flux} itself over height,
    $\Phi_h=\int\text{production\_flux}(h)\,\mathrm{d}h$; the product
    $\langle P\rangle_h\,\Phi_h$ equals $\int P(h)\,\text{production\_flux}(h)\,\mathrm{d}h$
    exactly, by the definition of the weighted average.
\end{itemize}
\end{implbox}

\begin{notebox}[title={Detector Composition Lives in the Pipeline Layer}]
Unlike \cref{ch:earth}, \texttt{medium.atmosphere} only ever propagates
production $\to$ surface: it has no notion of a detector, and no
\texttt{validation.py}. Weighting a surface or detector probability by
production-flux data, integrating a detector flux over production height,
and summing the incoherent contributions from several produced flavours are
all pipeline-level concerns, composed in
\texttt{pipeline.atmosphere\_earth} (\pyfunc{detector\_flux\_from\_production},
\pyfunc{integrate\_detector\_flux\_over\_height},
\pyfunc{sum\_detected\_flavours}) on top of
\pyfunc{pipeline.atmosphere\_earth.propagate\_surface\_to\_detector}, which
applies the Earth evolutor of \cref{ch:earth} to this medium's coherent
surface state via ordinary matrix-vector composition, not the matrix-matrix
composition used internally by \cref{sec:earth-evolutor}.
\end{notebox}


% =============================================================================
% \part{External Engines and Data --- \texttt{external/}}
% \input{chapters/10_external}            % TODO

% =============================================================================
% \part{Infrastructure, Orchestration, and Legacy Code}
% \input{chapters/11_peanuts_legacy}      % TODO
% \input{chapters/12_pipeline}            % TODO
% \input{chapters/13_util}                % TODO

% =============================================================================
% \appendix
% \input{chapters/99_apendices}           % TODO

\clearpage
\printbibliography[title={References}]

\end{document}
